\documentclass[11pt]{article}
\usepackage[a4paper,margin=27mm]{geometry}
\usepackage{amsmath,amssymb,amsthm,mathtools}
\usepackage{mathrsfs}
\usepackage[T1]{fontenc}
\usepackage{lmodern}
\usepackage{microtype}
\usepackage[hidelinks]{hyperref}
\urlstyle{same}
\hypersetup{pdfauthor={Tseng},
  pdftitle={Finite-time breakdown of Euler flows from smooth, compactly supported initial data in every fixed dimension d >= 3}}
\newcommand{\R}{\mathbb R}
\newcommand{\N}{\mathbb N}
\newcommand{\esssup}{\operatorname*{ess\,sup}}
\newcommand{\dd}{\,\mathrm d}
\newtheorem{theorem}{Theorem}[section]
\newtheorem{lemma}[theorem]{Lemma}
\newtheorem{proposition}[theorem]{Proposition}
\theoremstyle{definition}
\newtheorem{definition}[theorem]{Definition}
\title{Finite-time breakdown of Euler flows\\from smooth, compactly supported initial data\\in every fixed dimension $d\geq3$}
\author{Tseng}
\date{8 October 2026}
\setlength{\emergencystretch}{2em}
\begin{document}
\maketitle
\begin{abstract}
For every fixed integer $d\geq3$, we construct smooth, compactly supported,
divergence-free initial data for the incompressible Euler equation on
$\mathbb R^d$ whose maximal smooth Sobolev solution has finite lifetime.
The Frobenius norm of its velocity gradient has infinite endpoint limsup,
and the time integral of the $L^\infty$ norm of
$(\Omega_{ij})_{i<j}$ diverges. The pressure has the standard Newton normalization.
The proof extends the localized oscillatory construction and scalar
amplification mechanism of OpenAI. The higher-dimensional argument retains the full transverse
space and the full pressure coupling, replaces curl by antisymmetric
potentials, uses reflection only for the central endpoint selection, and
proves a quantitative ambient relative estimate. A variational history
solver, polynomial bounds for the lifted particle flow, and a
fixed physical mean cutoff yield a single frequency sequence for which
the initial increments are summable at every Sobolev order. Constants
may depend on the fixed dimension.
\end{abstract}

\section{Statement and conventions}
For each fixed integer $d\geq3$, let
$C^\infty_{c,\sigma}(\R^d)$ denote smooth compactly supported vector fields
with zero divergence. All dimension-dependent constants, derivative orders,
starting indices and scale thresholds may depend on this fixed $d$.
\begin{theorem}[$H_d$]\label{target:Hd}
There exists $u_0\in C^\infty_{c,\sigma}(\R^d)$ whose unique maximal smooth
Sobolev solution of
\[
 u_t+(u\cdot\nabla)u+\nabla P=0,\qquad \operatorname{div}u=0,\qquad u(0)=u_0
\]
has $0<T^*(u_0)<\infty$. Here the pressure is exactly
\[
 P=(-\Delta)^{-1}\sum_{i,j=1}^d\partial_i\partial_j(u_i u_j),
\]
and, on every $[0,S]$ with $S<T^*$, the solution lies in
$C([0,S];H^m)\cap C^1([0,S];H^{m-1})$ for every integer $m\geq1$.
With $\Omega_{ij}=\partial_i u_j-\partial_j u_i$, it satisfies
\[
 \limsup_{t\uparrow T^*}\esssup_x
   \left(\sum_{i,j}|\partial_i u_j(t,x)|^2\right)^{1/2}=\infty,
 \qquad
 \int_0^{T^*}\esssup_x
   \left(\sum_{i<j}|\Omega_{ij}(t,x)|^2\right)^{1/2}\,dt=\infty.
\]
\end{theorem}
The three-dimensional conclusion is the theorem of OpenAI~\cite{OAI},
with the Newton pressure and smooth Sobolev class specified above.
The present argument applies in that dimension with a transverse space
of dimension two. For $d>3$, every velocity and every localization is
defined on $\R^d$. In particular, the construction retains compact
support and finite energy in all spatial variables.

The scalar amplification equation, the variational history method, and
the finite oscillatory expansion follow the structure of~\cite{OAI}.
Their scalar and Hilbert-space estimates are reproduced below with their
hypotheses. The dimension-dependent steps are the antisymmetric-potential
localization, the full $(d-1)$-dimensional transverse evolution, its
coupling to the pressure, and the Sobolev orders in the nonlinear
correction and continuation argument. These steps are proved in the
ambient dimension. The parameter estimates then provide one frequency
sequence for which all initial Sobolev increments are summable.

Throughout, $C_d$ denotes a positive constant depending on the fixed
dimension and on fixed auxiliary cutoffs. Constants may change between
estimates. Dependence on a derivative order is indicated explicitly.
The notation $A\lesssim_d B$ means $A\leq C_d B$, and
$A\asymp_d B$ means both inequalities. The principal stage parameters are
listed below; each is defined precisely at its first use.
\begin{center}
\renewcommand{\arraystretch}{1.15}
\begin{tabular}{@{}ll@{}}
$t_0,\ \tau$ & physical activation time and scaled forward time\\
$\rho,\ k,\ \kappa=k^{-1}$ & physical localization length and frequency parameters\\
$\delta,\ f_\delta$ & angular width and periodic primary profile\\
$P$ & fixed angular period\\
$h_*,\ h_{\rm child}$ & inherited shear size and prescribed new shear size\\
$\Theta,\ e$ & scaled time bound and geometric comparison error\\
$\alpha,\ A_{\rm tar}$ & primary amplitude and unnormalized target size\\
$\mathcal P$ & bound for the fixed source coefficients and inverse estimates\\
$\ell_*,\ r_*$ & fixed physical mean-cutoff parameters
\end{tabular}
\end{center}

\section{The ambient equation and the endpoint criteria}
Fix an integer $d\geq3$. Constants may depend on $d$.
We use the Euclidean norm on vectors, the operator norm on linear maps when
estimating compositions, and the Frobenius norm in the statement of the
singularity conclusions. Put $s=d+2$. In particular $s>d/2+1$, so
$H^s(\R^d)$ embeds continuously into the space of fields with bounded
continuous velocity gradients. We use the
ordered-word Sobolev energy
\[
 E_m(u)=\|u\|_{H^m}^2
 =\sum_{r=0}^m\ \sum_{i_1,\ldots,i_r=1}^d
   \|\partial_{i_1}\cdots\partial_{i_r}u\|_2^2.
\]
It is equivalent, with constants depending only on $d,m$, to the usual
multi-index norm. The logarithmic estimate below uses the higher fixed
order $s_{\rm log}=d+5$. Both orders depend on $d$.

\begin{lemma}[Local theory and persistence]\label{lem:local}
Every solenoidal datum in all integer Sobolev spaces has a unique smooth
Sobolev Euler solution on an interval of length bounded below in terms of
its $H^s$ norm. The solution belongs to
$C_tH^m\cap C^1_tH^{m-1}$ on every shorter closed interval, for every
$m\geq1$. For $m\geq s$ it satisfies
\begin{equation}\label{eq:tame-energy}
 \frac{d}{dt}\|u\|_{H^m}\leq C_{d,m}\|\nabla u\|_\infty\|u\|_{H^m}.
\end{equation}
The same solution is obtained at every regularity level.
\end{lemma}
\begin{proof}
Let $J_N$ be convolution with an even, nonnegative, smooth compactly
supported kernel of integral one and support radius tending to zero.
Set $S_N=J_N\Pi$, where $\Pi$ is the orthogonal solenoidal projection,
and solve on $L^2(\R^d)$
\[
 \partial_tu_N=-S_N\bigl((S_Nu_N\cdot\nabla)S_Nu_N\bigr),
 \qquad u_N(0)=u_0.
\]
The smoothed quadratic right-hand side is locally Lipschitz on $L^2$.
The operators $S_N$ are self-adjoint contractions, commute with every
coordinate derivative, and have solenoidal range. Hence the solution
remains solenoidal. The integral equation and smoothing give all spatial
Sobolev derivatives for this fixed-scale solution. If $v_N=S_Nu_N$,
then for every derivative word $D^w$ the energy pairing is
\[
 \langle D^wu_N,D^w\partial_tu_N\rangle
 =-\langle D^wv_N,D^w((v_N\cdot\nabla)v_N)\rangle.
\]
This identity follows from self-adjointness of the mollifier.
The top differentiated transport term cancels by $\operatorname{div}v_N=0$.
For a smooth solution of the limiting Euler equation the same cancellation
holds, and the pressure term pairs to zero against each solenoidal
velocity derivative.
At order one the remaining integral is bounded directly by
$C_d\|\nabla u\|_\infty\|\nabla u\|_2^2$. For $m\geq2$ the
remaining terms are products with derivative orders $j$ and $m-j+1$. The Gagliardo--Nirenberg inequalities with endpoints
$\nabla u\in L^\infty$ and $D^m u\in L^2$ give, for $1\leq j\leq m$,
\[
 \|D^j u\|_{L^{2(m-1)/(j-1)}}
 \leq C_{d,m}\|\nabla u\|_\infty^{1-(j-1)/(m-1)}
                 \|D^m u\|_2^{(j-1)/(m-1)}.
\]
The exponent is interpreted as infinity at $j=1$. The two reciprocal
Lebesgue exponents sum to $1/2$, and the two interpolation parameters
sum to one. This proves the homogeneous top-order estimate. Summing the
lower orders and retaining conservation of $L^2$ energy proves
\eqref{eq:tame-energy}. The argument uses coordinate differentiation and integration by parts
in the ambient dimension.

Combining the preceding pairings, Sobolev embedding, and contraction of
$S_N$ in every derivative word gives, for $m\geq s$,
\[
 E_m(u_N)'\leq C_{d,m}\sqrt{E_s(u_N)}E_m(u_N),\qquad
 E_s(u_N)'\leq C_d(1+E_s(u_N))^2.
\]
An explicit common lifespan is
\[
 T_0=\frac{1}{2(1+C_d)(1+E_s(u_0))}.
\]
On this closed interval $E_s(u_N)\leq2E_s(u_0)+1$, uniformly in $N$.
Gronwall then bounds every higher $E_m(u_N)$ on the same interval.
These bounds include the $L^2$ norm required for Banach-space
continuation of the regularized quadratic equation.

To pass to the limit, translation followed by the fundamental theorem
of calculus gives
$\|J_Nf-f\|_2\leq\epsilon_N\|\nabla f\|_2$, with
$\epsilon_N\to0$. Apply this estimate to both inner smoothed factors
and the outer smoothing of the quadratic term. The uniform bounds
through order $d+4$ imply that
\[
 r_N:=\partial_tu_N+\Pi((u_N\cdot\nabla)u_N)
 \quad\hbox{satisfies}\quad
 \sup_{t\leq T_0}\|r_N(t)\|_2\longrightarrow0.
\]
For $u_N-u_M$, the solenoidal transport cancellation leaves an $L^2$
difference inequality with uniformly bounded gradient coefficient and
forcing $r_N-r_M$. All initial values equal $u_0$, so Gronwall's inequality implies that
$u_N$ is Cauchy in $C_tL^2$. Interpolation with the uniform higher bounds
gives convergence in $C_tH^m$ for every fixed $m$. The limits of the
derivative words are the weak derivatives of the same velocity field;
Sobolev embedding at arbitrarily high orders supplies its smooth
representative. Continuity of the projected quadratic term then gives
the limiting integral equation and its time derivative in every
$H^{m-1}$. It also supplies the one-sided derivatives at the endpoints.
The same difference estimate proves uniqueness. Consequently
the solutions produced using different cutoff orders or Sobolev indices
agree on their common intervals. This proves persistence and defines a
single maximal smooth Sobolev solution.
\end{proof}

\begin{lemma}[Newton pressure with fixed normalization]\label{lem:newton}
For a smooth solenoidal $u$ in all Sobolev spaces, set
$F=\sum_{i,j}\partial_i\partial_j(u_i u_j)$ and
$N_d(x)=|x|^{2-d}/((d-2)|\mathbb S^{d-1}|)$ for $x\ne0$.
Since $|\mathbb S^{d-1}|=d\,|B_1|$, this coefficient equals
$[d(d-2)|B_1|]^{-1}$. The integral
$P=N_d*F$ is absolutely convergent at every point and smooth.
Its gradient is the $L^2$ gradient part of $-(u\cdot\nabla)u$.
Thus the projected equation gives the Euler equation with the pressure
normalization specified by this convolution.
\end{lemma}
\begin{proof}
Every derivative of $F$ is both bounded and integrable: expand by the
product rule and use the $L^2$ bounds for the two factors of each product.
The kernel is locally integrable and bounded outside the unit ball, so
the Newton integral is absolutely convergent. Equivalently, with
$G_t(x)=(4\pi t)^{-d/2}e^{-|x|^2/(4t)}$,
\[
 N_d*F=\int_0^\infty G_t*F\,dt.
\]
Indeed $\int_0^\infty G_t(x)\,dt=N_d(x)$ for $x\ne0$, by the substitution
$r=|x|^2/(4t)$ and the Gamma integral. Absolute Fubini follows from
$\|G_t*|D^\gamma F|\|_\infty\leq
\min(\|D^\gamma F\|_\infty,C_dt^{-d/2}\|D^\gamma F\|_1)$.
This envelope is time integrable because $d>2$. It also permits
differentiation under the integral to every spatial order.

For each coordinate derivative, the small-time $L^2$ bound is
$\|\partial_i F\|_2$, while the large-time bound is
$\|\partial_iG_t\|_2\|F\|_1\leq C_dt^{-d/4-1/2}\|F\|_1$.
The latter is integrable at infinity for $d>2$. Hence $\nabla P\in L^2$,
and the same argument applies to all its derivatives. The heat equation
and its limits at zero and infinity give $-\Delta P=F$.
Writing $B=(u\cdot\nabla)u$, incompressibility gives
$\operatorname{div}B=F$, so $\nabla P+B$ is solenoidal.
Moreover $\nabla P$ lies in the closed $L^2$ gradient space. Indeed its
distributional curl vanishes. After Fourier transformation, its $L^2$
transform is parallel to $\xi$ for almost every $\xi\ne0$; this is
exactly the range of the orthogonal gradient multiplier
$\xi\otimes\xi/|\xi|^2$. The zero frequency is a set of measure zero.
Cutting off in frequency away from zero and then approximating the
potentials by compact smooth functions identifies this range with the
closed $L^2$ gradient space. Orthogonal
Helmholtz projection therefore yields $\nabla P=\Pi B-B$.

This normalization is also preserved under the limiting procedure in
Lemma~\ref{lem:local}. More generally, suppose
$u_n\to u$ in $C([0,S];H^m)$ for every finite $m$, and define $F_n,P_n$
by the same formulas. For each multi-index $\gamma$, the product rule,
Cauchy--Schwarz, and Sobolev embedding give
\[
 \sup_{t\le S}\bigl(\|D^\gamma(F_n-F)\|_1+
                    \|D^\gamma(F_n-F)\|_\infty+
                    \|\nabla D^\gamma(F_n-F)\|_2\bigr)\longrightarrow0.
\]
Splitting the heat integral at time one in the two estimates above yields
\begin{align*}
 \|D^\gamma(P_n-P)\|_\infty
 &\le C_d\bigl(\|D^\gamma(F_n-F)\|_\infty+
                         \|D^\gamma(F_n-F)\|_1\bigr),\\
 \|\nabla D^\gamma(P_n-P)\|_2
 &\le C_d\bigl(\|\nabla D^\gamma(F_n-F)\|_2+
                         \|D^\gamma(F_n-F)\|_1\bigr).
\end{align*}
These estimates are uniform on $[0,S]$. Thus the Newton scalars converge
with every spatial derivative uniformly, and their gradients converge
in every Sobolev space. The same estimates applied to two nearby times
give their time continuity. In particular the normalization is specified
by the convolution before and after passage to the limit.
\end{proof}

\begin{lemma}[Stability on varying intervals]\label{lem:stability}
Let $u$ be a smooth Euler solution on $[0,S]$ and let $v_n$ be smooth
solutions on $[0,S_n]$, where $0<S_n\leq S$. If
$v_n(0)\to u(0)$ in $H^s$, then
\[
 \sup_{0\leq t\leq S_n}\|v_n(t)-u(t)\|_{H^s}\longrightarrow0.
\]
No uniform higher norm of $v_n$ is assumed.
\end{lemma}
\begin{proof}
Put $w=v_n-u$. Its equation is
$w_t+(u+w)\cdot\nabla w+w\cdot\nabla u+\nabla q=0$.
Differentiation through order $s$, the transport cancellation, and the
Sobolev product estimate give
\[
 y'\leq C_d\|u\|_{H^{s+1}}y+C_d y^2,
 \qquad y=\|w\|_{H^s}.
\]
All coefficients involving $u$ are bounded on $[0,S]$ independently of
$n$. While $y\leq1$, Gronwall gives $y(t)\leq y(0)e^{CS}$.
For large $n$ this is less than $1/2$; continuity excludes a first time
when $y=1$. The estimate therefore holds on the full interval
$[0,S_n]$. Sobolev embedding also gives uniform convergence of the
velocity gradients there.
\end{proof}

\begin{lemma}[Continuation and the exact vorticity norm]\label{lem:continuation}
Let $u$ be the maximal solution from Lemma~\ref{lem:local}, and suppose
its maximal lifetime $T^*$ is finite. Write
$\Omega_{ij}=\partial_i u_j-\partial_j u_i$ and
$|\Omega|=(\sum_{i<j}|\Omega_{ij}|^2)^{1/2}$. Then
\[
 \limsup_{t\uparrow T^*}\|\nabla u(t)\|_{L^\infty,\mathrm F}=\infty,
 \qquad \int_0^{T^*}\|\Omega(t)\|_\infty\,dt=\infty.
\]
\end{lemma}
\begin{proof}
If $\|\nabla u\|_\infty$ were bounded, \eqref{eq:tame-energy} would
bound the $H^s$ norm up to $T^*$. Restarting the local solution at times
approaching $T^*$ gives a common positive extension length.
The same energy estimate propagates every higher norm on a smaller
common interval, so uniqueness gives an extension in the full smooth
Sobolev class, contrary to maximality. A finite endpoint limsup would
imply such a bound, since the solution is smooth on shorter intervals.
The inequalities $\|A\|_{\rm op}\leq\|A\|_{\rm F}\leq
\sqrt d\,\|A\|_{\rm op}$ identify the required Frobenius conclusion.

For the second assertion, incompressibility gives
$\sum_i\partial_i\Omega_{ij}=\Delta u_j$. Use the Gaussian averages
\[
 Q_tf(x)=(\pi t)^{-d/2}\int_{\R^d}e^{-|y|^2/t}f(x+y)\,dy,
 \qquad \partial_tQ_tf=\tfrac14 Q_t\Delta f.
\]
For $f=u_j$, $g_i=\Omega_{ij}$, and $0<\epsilon<1$, write
\[
 \partial_k f=Q_1\partial_k f+
       (Q_\epsilon-Q_1)\partial_k f+
       (I-Q_\epsilon)\partial_k f.
\]
The first term is at most $C_d\|f\|_2$, by moving its derivative to
the Gaussian kernel and using its $L^2$ norm. In the middle term use
$\Delta f=\sum_i\partial_i g_i$ and move two spatial derivatives to
the Gaussian kernel. Its second derivative has $L^1$ norm at most
$C_d/t$, so time integration from $\epsilon$ to one bounds this term by
$C_d W(-\log\epsilon)$ whenever $|g_i(x)|\leq W$.
The final term is bounded by
$C_d\epsilon\|D^2\partial_k f\|_\infty$: integrate the displayed heat
identity from zero to $\epsilon$, using contraction of Gaussian averaging
on bounded functions. All integrations by parts are justified by the
Gaussian decay and the bounded smooth derivatives of the fields.

Let
$\mathcal H(u)=\sum_{r=0}^{d+5}\|D^ru\|_2$, with the operator tensor
norm in each $L^2$ term. Sobolev embedding applied to each third derivative
gives $\|D^2\partial_k u_j\|_\infty\leq C_d\mathcal H(u)$.
Moreover $\mathcal H(u)\leq C_d\sqrt{E_{s_{\rm log}}(u)}$.
Choose $\epsilon=(e+C_d\mathcal H(u))^{-4}$ in the three bounds.
The high remainder is bounded by a dimensional constant, while the middle
term is logarithmic. Each entry of $\Omega$, including a lower-triangular
entry by antisymmetry, is bounded by the once-counted norm
$|\Omega|=(\sum_{i<j}|\Omega_{ij}|^2)^{1/2}$. Summing the finite
coordinate estimates therefore proves
\begin{equation}\label{eq:log-vorticity}
 \|\nabla u\|_{\infty,\mathrm{op}}\leq
 C_d\bigl(1+\|u(0)\|_2+
       \|\Omega\|_\infty\log(e+\mathcal H(u))\bigr).
\end{equation}
Here $L^2$ energy conservation has replaced $\|u(t)\|_2$ by
$\|u(0)\|_2$.

Put $K(t)=\|\nabla u(t)\|_{\infty,\mathrm{op}}$,
$W(t)=\|\Omega(t)\|_\infty$, and $G(t)=\int_0^tK(r)\,dr$.
The gradient energy inequality at order $s_{\rm log}$ implies
\[
 E_{s_{\rm log}}(u(t))\leq E_{s_{\rm log}}(u(0))e^{C_dG(t)},
 \qquad \log(e+\mathcal H(u(t)))\leq C_0+C_dG(t),
\]
where $C_0$ depends only on the initial high norm and $d$.
Consequently \eqref{eq:log-vorticity} gives
$G'\leq C(1+W)(1+G)$, and hence
\[
 1+G(t)\leq\exp\!\left(C\left(t+\int_0^tW(r)\,dr\right)\right).
\]
A finite vorticity integral on $[0,T^*)$ would bound $G$ uniformly.
The energy inequality at the lower order $s=d+2$ would then bound
$E_s$ uniformly, permitting the restart already proved and contradicting
maximality. This argument differentiates the integral $G$, and does not
require differentiability of a supremum norm or of $\mathcal H$.

On every shorter closed time interval, the bounded gradient and vorticity
tensors depend continuously on time in the uniform norm: their differences
are controlled by the $H^{d+2}$ difference of the velocities. Their norm
paths are therefore continuous and nonnegative. The ordinary integrals
on these shorter intervals equal the corresponding nonnegative Lebesgue
integrals; monotonicity makes a finite integral on $[0,T^*)$ a common
bound for all of them. Thus the contradiction proves divergence in the
extended nonnegative sense used in the statement.
For smooth fields the pointwise and essential suprema agree: a strict
pointwise lower bound persists on an open ball of positive measure.
Thus these are the essential-supremum norms in Theorem~\ref{target:Hd}.
\end{proof}

\section{Full-dimensional solenoidal geometry}
Let $E=\operatorname{span}(e_1,e_2,e_3)$ and let
$\mathcal R=I_E\oplus(-I_{E^\perp})$. The fields used below have odd
parity and satisfy $u(\mathcal R x)=\mathcal R u(x)$.
The fields are constructed directly on $\R^d$.

\begin{lemma}[Antisymmetric potentials]\label{lem:potential}
If $A_{ij}=-A_{ji}$ is a smooth compactly supported tensor and
$v_j=\sum_i\partial_iA_{ij}$, then $\operatorname{div}v=0$.
For any trace-free matrix $L$,
\[
 A_{ij}(x)=\frac{x_i(Lx)_j-x_j(Lx)_i}{d}
 \quad\Longrightarrow\quad \sum_i\partial_i A_{ij}=(Lx)_j.
\]
Multiplying this tensor by an even compact Gevrey cutoff equal to one
near zero produces an odd compact solenoidal velocity equal to $Lx$
near zero. If $L\mathcal R=\mathcal R L$ and the cutoff is radial, the
velocity is $\mathcal R$-equivariant.
\end{lemma}
\begin{proof}
Commutation of mixed derivatives and exchange of $i,j$ make
$\sum_{i,j}\partial_j\partial_iA_{ij}$ equal to its negative.
For the displayed tensor, differentiation of the first numerator term
gives $(d+1)(Lx)_j$, and differentiation of the second gives
$(Lx)_j+x_j\operatorname{tr}L$. Their difference is $d(Lx)_j$.
The cutoff preserves antisymmetry and compact support. The tensor is
even, so its divergence is odd. Under $\mathcal R$ the tensor transforms
as a two-tensor; applying divergence gives the stated equivariance.
Gevrey bounds follow from the finite product rule with a polynomial.
\end{proof}

Antisymmetric potentials also provide oscillatory solenoidal
corrections. Put $D_i=\kappa\partial_{y_i}+m_i\partial_\theta$.
These constant-coefficient operators commute. Thus
$\sum_iD_i\sum_jD_j A_{ji}=0$ for every antisymmetric $A$.
When $m\cdot v=0$, $|m|=1$ and $F'=f$,
\[
 A_{ij}=(m_i v_j-m_j v_i)F(\theta)
 \quad\Longrightarrow\quad
 \sum_iD_i A_{ij}=v_jf(\theta)+\kappa\sum_i\partial_{y_i}A_{ij}.
\]
The last term is the spatial corrector and is retained at the
terminal truncation order. This identity treats every ambient component.

\begin{lemma}[Symmetry of the ambient pressure and central blocks]\label{lem:reflection}
If $u$ is $\mathcal R$-equivariant, its Newton pressure obeys
$P(\mathcal R x)=P(x)$. At $x\in E$, both $\nabla u(x)$ and
$\nabla^2P(x)$ preserve $E$ and $E^\perp$.
\end{lemma}
\begin{proof}
The chain rule and orthogonality show that
$\sum_{i,j}\partial_i\partial_j(u_i u_j)$ is invariant under
$\mathcal R$. The Newton kernel is radial and Lebesgue measure is
orthogonally invariant, so a change of variables in its absolutely
convergent convolution proves the pressure identity. Differentiate the
velocity identity once and the pressure identity twice at a fixed point.
The resulting maps commute with $\mathcal R$, and therefore preserve its
two eigenspaces. The trace of the full velocity gradient vanishes; its restriction to
$E$ need not be trace free.
\end{proof}

\begin{lemma}[The endpoint operator preserves the active plane]\label{lem:endpoint-symmetry}
Along the central trajectory of the equivariant preceding solution,
suppose the normal
$m(t)$ belongs to $E$. By Lemma~\ref{lem:reflection}, its gradient $M(t)$
and pressure Hessian $H(t)$ commute with $\mathcal R$. For the
stationary displacement problem on $[0,t_0]$ with
$m(t)\cdot\eta(t)=0$, $\eta(0)=0$, prescribed $\eta(t_0)=Y$, and
coercive action
$\mathcal A(\eta)=\int_0^{t_0}(|\eta'|^2-\langle H\eta,\eta\rangle)dt$,
the endpoint map $\Lambda Y=\operatorname{proj}_{m(t_0)^\perp}\eta'(t_0)$
commutes with $\mathcal R$. In particular a prescribed endpoint in
$E\cap m(t_0)^\perp$ produces a displacement and velocity in $E$.
\end{lemma}
\begin{proof}
Reflection preserves the constraint, the two endpoints and the action.
Uniqueness for the coercive variational problem therefore gives
$\eta_{\mathcal RY}=\mathcal R\eta_Y$. Taking the terminal derivative
and its orthogonal projection proves the claim about $\Lambda$.
If $Y=\mathcal RY$, uniqueness also gives
$\eta_Y=\mathcal R\eta_Y$ at every time. The identity remains true for
the physical velocity $\eta'-M\eta$ because $M$ commutes with the same
reflection. Hence its $E^\perp$ components vanish, including after endpoint
normalization.
\end{proof}
\section{Polynomial bounds from the lifted particle flow}
The estimates in this section control the particle flow through the
small lifted transport field. The resulting bounds depend polynomially
on the physical frequency.

Write $\mathcal L=\R^d\times\R$, with its product norm, and put
$J_kx=(x,k\langle m_0,x\rangle)$, $|m_0|=1$. For a lifted packet solution
$z(t,y,\theta)$ the lifted transport field is
\[
 b(t,y,\theta)=\bigl(\kappa z(t,y,\theta),\,
                         \langle m_0,z(t,y,\theta)\rangle\bigr),
 \qquad k\kappa=1.
\]
Throughout this section, $z$ is the finite packet together with the
correction that eliminates its residual. The normal component in
the second entry retains the cancellation that gives the small drift estimate.

\begin{lemma}[Invariance of the physical graph]\label{lem:graph-invariant}
Let $b$ be a smooth bounded Lipschitz field on $\mathcal L$, continuous
in time, of the displayed form. Its global flow $\Psi_{s,t}$ satisfies
\[
 \Psi_{s,t}(J_kx)=J_k\psi_{s,t}(x),\qquad
 \psi_{s,t}(x)=\pi\Psi_{s,t}(J_kx),
\]
where $\pi(y,\theta)=y$. The map $\psi$ is the flow of
$v(t,x)=\pi b(t,J_kx)$ and $\psi_{t,s}$ is its inverse.
For any $\ell>0$, the flow of
$v_\ell(t,x)=\ell v(t,x/\ell)$ is
$\ell\psi_{s,t}(x/\ell)$.
\end{lemma}
\begin{proof}
The linear functional $q(y,\theta)=\theta-k\langle m_0,y\rangle$
satisfies $q(b)=0$. Its derivative along every trajectory is zero.
Therefore the level set $q=0$, which is precisely the image of $J_k$,
is invariant. Projection of the lifted ODE gives the ODE for
$v$; uniqueness identifies its flow and inverse. Substitution and the chain rule prove the scaling identity in every
ambient dimension.
\end{proof}

\begin{lemma}[Small lifted flow bounds]\label{lem:small-lift}
Suppose, for $0\leq t\leq T$ and every $n\geq0$,
\[
 \|D^n b(t)\|_\infty\leq B R^n(n!)^2,
 \quad B\geq0,\ R>0,\quad BRT\leq\tfrac18.
\]
Then the displacement $f_t=\Psi_{0,t}-\mathrm{id}$ obeys
\begin{equation}\label{eq:lift-disp}
 \|D^n f_t\|_\infty\leq Bt(4R)^n(n!)^2 \quad(n\geq0).
\end{equation}
Put $\mathscr R(S)=(4R+1)((1+BT)S+2)$.
If also $\|D^n b_t\|_\infty\leq B_1R_1^n(n!)^2$ with
$B_1\ge0$ and $R_1\ge0$, then the
material velocity and acceleration satisfy
\begin{align}
 \|D^n(b(t)\circ\Psi_{0,t})\|_\infty
   &\leq B\mathscr R(R)^n(n!)^2,\label{eq:lift-velocity}\\
 \|D^n((b_t+Db\,b)(t)\circ\Psi_{0,t})\|_\infty
   &\leq(B_1+3B^2R)\mathscr R(4R+R_1)^n(n!)^2.
   \label{eq:lift-acceleration}
\end{align}
These estimates hold in every finite-dimensional normed space.
\end{lemma}
\begin{proof}
Apply the finite generating-sum argument of~\cite[\S3.8]{OAI} to
the integral flow equation. For a finite derivative cutoff $N$, let
\[
 F_N(t)=\sum_{n=1}^N\frac{r^n}{(n!)^2}
                 \|D^n f_t\|_\infty,\qquad r=(4R)^{-1}.
\]
The differentiated equation $f_t=\int_0^t b(s)\circ
(\mathrm{id}+f_s)\,ds$ and the finite Fa\`a di Bruno convolution give
\[
 F_N(t)\leq\int_0^t
  \frac{BR(r+F_N(s))}{1-R(r+F_N(s))}\,ds
\]
as long as the denominator is positive. The algebra behind this
inequality is the scalar factorial majorant for a composition: after
division by $(n!)^2$, each partition coefficient is bounded by the
corresponding geometric-series coefficient. It applies to operator
norms of multilinear derivatives, independently of the ambient
dimension. Under the bootstrap assumption $F_N\leq Bt$, one has
$R(r+F_N)\leq1/4+BRT\leq3/8$, so the integrand is at most $3B/5$.
Continuity and a first-contact argument give $F_N(t)\leq Bt$ for all
$t\leq T$, uniformly in $N$. The $n$th summand gives
\eqref{eq:lift-disp} for $n>0$; order zero follows by directly integrating
$\|b\|_\infty\leq B$. The finite generating sums may be estimated pointwise before taking
the supremum, so differentiation of the supremum is unnecessary.

Adding the identity gives, for $n>0$,
$\|D^n\Psi_{0,t}\|_\infty\leq
(1+BT)(4R+1)^n(n!)^2$. The same finite composition formula implies
that composing a field bounded by $C S^n(n!)^2$ with this map is bounded
by $C\mathscr R(S)^n(n!)^2$. This gives
\eqref{eq:lift-velocity}. Differentiating a factorial bound once and
using $(n+1)^2\leq4^n$ bounds $Db$ by $BR(4R)^n(n!)^2$.
The factorial Leibniz convolution is at most three times the product
of the two amplitudes. At the common radius $4R+R_1$ this bounds
$b_t+Db\,b$ by $(B_1+3B^2R)(4R+R_1)^n(n!)^2$.
The composition estimate proves \eqref{eq:lift-acceleration}.
\end{proof}

\begin{lemma}[Physical labels and their inverse]\label{lem:physical-labels}
The graph and physical scaling in Lemma~\ref{lem:graph-invariant}
convert a bound $\|D^n f\|_\infty\leq C S^n(n!)^2$ into
\[
 \left\|D^n\bigl[\ell\pi f(J_k(\cdot/\ell))\bigr]\right\|_\infty
 \leq C\bigl(\ell^{-1}S(1+|k|)\bigr)^n(n!)^2
 \quad(0<\ell\leq1).
\]
Consequently the physical displacement, material velocity and material
acceleration have factorial bounds with amplitudes and radii polynomial
in $1+B+R+B_1+R_1+T+\ell^{-1}+|k|$. If the physical field is
solenoidal, its inverse deformation, evaluated in particle labels, has
factorial bounds polynomial in
the same parameters, with exponents allowed to depend on $d$.
\end{lemma}
\begin{proof}
The chain rule for a linear inner map gives the factor
$\|J_k/\ell\|^n$, and $\|J_k\|\leq1+|k|$. The outer map has norm at
most $\ell\leq1$. This proves the first assertion, including order
zero. Apply it to \eqref{eq:lift-disp}--\eqref{eq:lift-acceleration}.
The first spatial derivative of a flow is the identity plus the
derivative of its displacement, so it has the same type of polynomial
factorial bound after enlarging the parameters.

Solenoidality gives $\det D\psi=1$ by the variational equation and
Jacobi's formula. Thus $(D\psi)^{-1}=\operatorname{adj}(D\psi)$.
Every entry of the adjugate is a sum of $(d-1)!$ products of $d-1$
entries. For a product of $r$ factorially bounded fields, the
multinomial Leibniz formula and
$\prod_{i=1}^r n_i!\leq n!$ show that its $n$th derivative is bounded
by $C_d C^r(rR)^n(n!)^2$. Here $r=d-1$ is fixed.
Hence the inverse deformation has a polynomial factorial bound, with
no exponential in the physical frequency. For composition with the preceding flow, the finite composition and
product formulas combine its factorial bounds with the graph-flow bounds
above. Two time derivatives are controlled by the material velocity and
acceleration estimates.
\end{proof}

The trace identity needed in this construction is also
dimension independent. If a lifted field is tangent to the graph,
$b=J_k\pi b$. Differentiating and cyclically commuting the trace gives
\[
 \operatorname{tr}_{\R^d}(\pi Db\,J_k)
   =\operatorname{tr}_{\mathcal L}(J_k\pi Db)
   =\operatorname{tr}_{\mathcal L}(Db).
\]
Thus lifted divergence zero implies ordinary divergence zero for the
graph velocity. Scaling preserves it. The graph flow and the preceding particle map both preserve
$d$-dimensional volume, as does their composition.

\subsection{Cylinder and graph bounds in \texorpdfstring{$L^2$}{L2}}
To derive global $L^2$ estimates, let $b$ be periodic in its last
variable with period $P$ and assume its cylinder divergence is zero. Suppose the input bounds
of Lemma~\ref{lem:small-lift} hold both in supremum norm and in
$L^2(\R^d\times\mathbb T_P)$, with common amplitudes $B,B_1$ and
radii $R,R_1$. The weighted Sobolev estimates for the packet and its time derivative
provide these input $L^2$ bounds.

The flow satisfies
$\Psi_t(y,\theta+P)=\Psi_t(y,\theta)+(0,P)$ by uniqueness.
It therefore descends to the cylinder. Its variational Jacobian obeys
\[
 \partial_t\det D\Psi_t
       =\bigl(\operatorname{div}b\bigr)(t,\Psi_t)\det D\Psi_t=0.
\]
The inverse flow exists by the same bounded Lipschitz ODE, so change of
variables gives
\begin{equation}\label{eq:lift-volume-l2}
 \|a\circ\Psi_t\|_{L^2(\R^d\times\mathbb T_P)}=\|a\|_2.
\end{equation}
For the lifted velocity of the corrected packet,
$\operatorname{div}b=\kappa\operatorname{div}_y z+
\partial_\theta(m_0\cdot z)=0$ is the closed solenoidal constraint, which holds pointwise by
smoothness.

Set $S_v=\mathscr R(R)$ and $S_a=\mathscr R(4R+R_1)$, with
$\mathscr R$ as in Lemma~\ref{lem:small-lift}. Then the fields
satisfy
\begin{align}
 \|D^n(b(t)\circ\Psi_t)\|_2
    &\le B S_v^n(n!)^2,\label{eq:lift-velocity-l2}\\
 \|D^n((b_t+Db\,b)(t)\circ\Psi_t)\|_2
    &\le(B_1+3B^2R)S_a^n(n!)^2,\label{eq:lift-acceleration-l2}\\
 \|D^n(\Psi_t-\mathrm{id})\|_2
    &\le Bt S_v^n(n!)^2.\label{eq:lift-displacement-l2}
\end{align}
To prove these assertions, expand the derivative of a composition into
its finite Fa\`a di Bruno sum. In each term place the single factor
$D^r a\circ\Psi_t$ in $L^2$ using \eqref{eq:lift-volume-l2}; place
all derivatives of the inner flow in $L^\infty$ using
Lemma~\ref{lem:small-lift}. The numerical composition convolution is
the same as for the supremum bound, so its radius is
$\mathscr R(S)$ for an outer radius $S$. For $Db\,b$, put one factor
in $L^2$ and the other in $L^\infty$ in each Leibniz term.
The factorial convolution is bounded by three, giving amplitude
$3B^2R$ at radius $4R$. This proves the first two estimates.
Finally integrate the identity
$\Psi_t-\mathrm{id}=\int_0^t b(s)\circ\Psi_s\,ds$ and apply
Minkowski's inequality to each derivative. This proves the last one,
including order zero, with radius $S_v$.

The graph trace estimate includes the physical scaling. If a smooth
periodic field $a$ has
$\|D^n a\|_2\le C S^n(n!)^2$, then
\begin{equation}\label{eq:graph-l2-jet}
 \left\|D^n\{\ell\pi a(J_k(\cdot/\ell))\}\right\|_{L^2(\R^d)}
 \le C_P\ell^{1+d/2}C(1+S)
       \left(\frac{4S(1+|k|)}\ell\right)^n(n!)^2.
\end{equation}
Indeed the one-dimensional periodic estimate
\[
 \sup_\theta|v(\theta)|^2
   \le\frac2P\int_{\mathbb T_P}|v|^2
                   +2P\int_{\mathbb T_P}|v_\theta|^2
\]
follows from the fundamental theorem of calculus, the triangle
inequality and Cauchy--Schwarz. The same constants apply to any
finite-dimensional normed target, in particular to the operator norm
of $D^n a$, independently of $n$. Applied at each spatial point it
bounds evaluation on any phase graph by
$C_P(\|a\|_2+\|\partial_\theta a\|_2)$; the phase slope does not
enter this trace estimate. Differentiation of the affine graph introduces the factor
$\|J_k\|^n\le(1+|k|)^n$ before applying that trace. The additional
angular derivative is bounded using
$(n+1)^2\le4^n$. Finally $x=\ell y$ gives the exact factor
$\ell^{1+d/2-n}$. For $0<\ell\le1$ the prefactor
$\ell^{1+d/2}$ is bounded above by one and may be omitted. This proves
\eqref{eq:graph-l2-jet}. The argument applies to the
displacement, material velocity and material acceleration in
\eqref{eq:lift-velocity-l2}--\eqref{eq:lift-displacement-l2}.

For the parameter induction, use the source polynomial
$C_{\rm rel}(W)\geq1$ obtained by inserting the packet bounds into
the preceding finite composition formulas. Here $W\geq1$ bounds the
source coefficients, the common packet radius, and its profile amplitude.
The correction that eliminates the residual has amplitude
$\delta_{\rm corr}(X)$, where $X$ is its expansion parameter.
The lifted drift bound is a source polynomial times
$k^{-1}+\delta_{\rm corr}(X)$. Thus, with $T\leq1$ and the physical
scale $\ell=\rho\leq1$, the same explicit polynomial can be enlarged to
give
\[
 BRT\leq C_{\rm rel}(W)(k^{-1}+\delta_{\rm corr}(X)),\qquad
 M_{\rm rel}\leq C_{\rm rel}(W),\qquad
 S_{\rm rel}\leq C_{\rm rel}(W)\rho^{-1}k.
\]
Here $M_{\rm rel},S_{\rm rel}\geq1$ are common factorial amplitudes and
radii for the three relative fields in $L^2$, and for the
displacement and material velocity also in supremum norm. The latter
supremum bounds come directly from the small lifted flow estimates;
applying a physical Sobolev embedding here would introduce unnecessary
powers of $\rho^{-1}k$ into the amplitude.

The prescribed scales satisfy the smallness conditions
\[
 16C_{\rm rel}(W)\leq k^{1/2},\qquad
 \rho^{-1}\leq k^{1/2},\qquad
 \delta_{\rm corr}(X)\leq k^{-1},\qquad k\geq1.
\]
Consequently $BRT\leq1/8$, $M_{\rm rel}\leq k$, and
$S_{\rm rel}\leq k^2$. These bounds will be used in the parameter induction.
All three fields are restrictions of the same lifted flow. In original
physical labels their relative map is
$\psi_t(x)=\rho\,\psi_t^{\rm norm}(x/\rho)$, and the updated
particle map is $X_{\rm parent}(t,\psi_t(x))$. The composition uses the label bounds of the preceding map in its
original physical labels.

\subsection{Fixed Sobolev shifts and label renewal}
For the evaluation step fix $q=d+2$, and suppose a smooth Euclidean
vector-valued field satisfies
$\|D^na\|_2\leq C R^n(n!)^2$ for every $n\geq0$, with $C,R\geq0$.
Define
\[
 A_d(C,R)=c_d C\sum_{j=0}^q(4R)^j(j!)^2,
 \qquad R_d(R)=4\max(1,d)R,
\]
where $c_d$ is the fixed Sobolev evaluation constant. Then
\begin{equation}\label{eq:fixed-sobolev-jet-shift}
 \|D^na\|_\infty\leq A_d(C,R)R_d(R)^n(n!)^2.
\end{equation}
Indeed, apply the fixed $H^q$ evaluation inequality to each coordinate
word of length $n$. Its derivative of order $j\leq q$ has $L^2$ norm
bounded by $\|D^{n+j}a\|_2$. The factorial inequality
\[
 ((n+j)!)^2\leq4^{n+j}(n!)^2(j!)^2
\]
gives the amplitude $A_d$ and coordinate-word radius $4R$.
The operator norm of an $n$-linear derivative is at most the sum of its
$d^n$ coordinate evaluations, which gives exactly $R_d(R)$.
Thus the derivative loss is fixed by $d$, independently of packet grade
or truncation.

The composition with the preceding flow has the following
chain-rule expressions. Write $X_{\rm parent}(t,x)=x+A(t,x)$, with material
velocity $U=\partial_tA$ and acceleration $W_0=\partial_tU$. Write the
relative fields as $a,b,c$ with $\partial_ta=b$, $\partial_tb=c$, and
$\psi=x+a$. The updated displacement, velocity, and acceleration are
\begin{align*}
 A_{\rm new}&=a+A\circ\psi,\\
 U_{\rm new}&=U\circ\psi+(I+DA\circ\psi)b,\\
 W_{\rm new}&=W_0\circ\psi+2(DU\circ\psi)b
       +(D^2A\circ\psi)[b,b]+(I+DA\circ\psi)c.
\end{align*}
Suppose the three preceding fields have both $L^2$ and supremum factorial
bounds with amplitude and radius $K\leq k$, $K\geq1$. Let
$M=M_{\rm rel}\leq k$ and $S=S_{\rm rel}\leq k^2$.
Since $\psi$ preserves volume, each factor composed with the preceding flow retains its
$L^2$ norm before applying the derivative composition formula. The
relative displacement and velocity provide the needed supremum factors;
the relative acceleration is placed in $L^2$. Differentiating a
factorial bound changes amplitude $C$ and radius $R$ to $CR$ and $4R$.
The first and second derivative amplitudes of the preceding fields are
consequently
$K^2$ and $4K^3$. The finite product and composition estimates bound all
three updated fields in $L^2$ by the common amplitude and radius
\[
 C_{\rm new}=K+M+9K^2M+36K^3M^2,\qquad
 R_{\rm new}=(1+S)\bigl((1+M)16K+2\bigr)+S.
\]
For $k\geq69$ these satisfy
$C_{\rm new}\leq47k^5\leq k^6$ and
$R_{\rm new}\leq69k^4\leq k^5$.

The evaluation step \eqref{eq:fixed-sobolev-jet-shift} now applies to
these composed fields. Enlarge the fixed lower bound on $k$ to include
$4\max(1,d)$ and $c_d\sum_{j=0}^q(j!)^2$. Since
$4R_{\rm new}\leq k^6$, it gives
\[
 A_d(C_{\rm new},R_{\rm new})\leq k^{6q+7},\qquad
 R_d(R_{\rm new})\leq k^6.
\]
The sum of one and the amplitude and radius bounds in $L^2$ and
supremum norm is at most
$5k^{6q+7}\leq k^{6q+8}\leq k^{10(d+4)}$.
Thus $K_{\rm new}=k^{10(d+4)}$ is a single valid label invariant for
all three updated fields in both norms. Differentiating these
fields gives the frame coefficients
$F=I+DA_{\rm new}$, $F_1=DU_{\rm new}$, and $F_2=DW_{\rm new}$,
with amplitudes $1+K_{\rm new}^2,K_{\rm new}^2,K_{\rm new}^2$ and
radius $4K_{\rm new}$. The determinant-one inverse argument above gives
the corresponding polynomial bound for the inverse frame of the updated
physical particle map.
\section{The active three-plane and ambient amplification}
All cross products in this section are taken in the oriented Euclidean
three-plane $E$ and define an orthonormal basis of $E$. Write
$M=\nabla u(t,0)$ for the central gradient of the preceding solution and
$H=\nabla^2P(t,0)$. Oddness fixes the central
particle at zero. Reflection implies that $M,H$ are block diagonal on $E\oplus E^\perp$.
The active block need not be trace free. Norms of $M,H,B$ in this section are suprema in time at the central
particle. Full-space coefficient bounds enter through the polynomial
particle-label estimates, and the spatial history inverse uses the global
one-sided pressure bound.

The geometric input at a stage is the decomposition
\begin{equation}\label{eq:central-decomposition}
 M=B+hqp^T+\mathcal E,\qquad |\mathcal E|\leq E_*,\qquad h(t_0)=h_*,
\end{equation}
where $p,q$ are the normalized transported normal and primary velocity of
the preceding packet. They lie in $E$ and are orthogonal; put $n=p\times q$.
The shear is the preceding packet amplitude,
$h=c_0|m_{\rm old}||v_{\rm old}|$ with constant $c_0>0$.
Consequently $\dot h/h=-B_{pp}-B_{qq}$, by the two transport equations
for the preceding normal and velocity.
Here $B$ is the gradient of the preceding background solution.
On the stage interval $|B|\leq C G_*$ and $|\dot B|\leq C G_*^2$.
The transported normal and velocity satisfy
\begin{equation}\label{eq:ray-velocity}
 \dot m=-M^Tm,\qquad
 \dot v=-Mv+2m\frac{m\cdot Mv}{|m|^2},\qquad m\cdot v=0.
\end{equation}
The coefficient 2 follows by differentiating the last constraint. These
ambient equations preserve $E$ for central data in $E$.

At activation put
\begin{gather*}
 a=p\cdot Bq\in[1/2,2],\quad \beta=(n\cdot Bq)/a,
 \quad \tfrac12\leq\beta x_{\rm prev}^2\leq2,\\
 \epsilon=\sqrt{a/h_*},\quad
 \tau=a(t-t_0)/\epsilon,\quad x=\sqrt\beta\,\tau.
\end{gather*}
The target is $x=x_{\rm tar}\geq2$. Write
$\tau_{\rm tar}=x_{\rm tar}/\sqrt\beta$ and use the scaled endpoint
$\Theta=\tau_{\rm tar}+\Delta\tau$. Here $0<\sqrt\beta\leq1/4$ and
$0<\Delta\tau\leq\Theta^{-60}$, so in particular
$\Theta\geq\max\{2,\beta^{-1/2},\tau_{\rm tar}\}$.
For the next prescribed physical time width $w_{\rm next}$, set
$\Delta\tau=2\sqrt{a h_*}\,w_{\rm next}$, and
$(\epsilon/a)\Delta\tau=2w_{\rm next}$.
The scalar parameter choice provides an envelope $\widehat\Theta\geq\Theta$
and $\Delta\tau\,\widehat\Theta^{60}\leq1$; monotonicity of reciprocal
powers proves the stated bound for $\Theta$.
The geometric error parameter is
\begin{equation}\label{eq:geometric-error}
 e=C_d\bigl(\epsilon\Theta G_*^2+E_*+
                   \rho K_h^{c_d}/\epsilon\bigr),
 \qquad e\Theta^{40}\leq c_d.
\end{equation}
Here $\rho$ is the physical packet length and $K_h$ bounds the
preceding particle map, its inverse deformation and the required derivatives.
The small constant on the right is fixed before the scales are chosen.

\begin{lemma}[Endpoint selection]\label{lem:actual-endpoint}
On a positive history interval $[0,t_0]$, assume the one-sided
pressure bound $H\leq K_+I$ and $K_+t_0^2/2\leq1/2$, provided by
the mean reserve. Suppose $\|M\|_\infty\leq C_Mh_*$,
$\|H\|_\infty\leq C_Hh_*h_{\rm old}$, $1\leq h_{\rm old}\leq h_*$,
and $t_0^{-1}\leq h_*$. Suppose also
$\overline B=B(t_0)+\mathcal E(t_0)$ has
$\overline B_{pp}<0$ and $\|\overline B\|\leq\zeta h_*$.
For sufficiently small $\zeta$, depending only on $C_M,C_H$, an endpoint
$Y\in E\cap n^\perp$ of size at most $C/h_*$ gives a stationary displacement
whose terminal primary velocity satisfies
\[
 v_q=1,\qquad -C\leq v_p\leq0,\qquad v_{E^\perp}=0.
\]
The displacement is stationary in the full transverse space.
\end{lemma}
\begin{proof}
Choose $m_0=F(t_0,0)^Tn/|F(t_0,0)^Tn|$ and transport it by $F^{-T}$.
Let $\Lambda$ be the endpoint operator of Lemma~\ref{lem:endpoint-symmetry}.
Integration by parts in its stationary equation gives
\[
 Y\cdot\Lambda Z=\int_0^{t_0}
    (\eta_Y'\cdot\eta_Z'-H\eta_Y\cdot\eta_Z)\,dt.
\]
The stationary paths start at zero. The inequality
$\int|\eta|^2\leq(t_0^2/2)\int|\eta'|^2$ therefore implies that their
action is nonnegative even for a nonzero terminal value. Thus $\Lambda$
is symmetric and positive semidefinite. To bound it, put
$L_0=t_0^{-1}+\|M\|_\infty+\|H\|_\infty^{1/2}+1$.
A comparison path vanishes before $t_0-L_0^{-1}$ and then equals
$L_0(t-t_0+L_0^{-1})\operatorname{proj}_{m(t)^\perp}Y$.
Since $\|\partial_t\operatorname{proj}_{m^\perp}\|\leq C\|M\|$,
its action is at most
$C(L_0+\|M\|_\infty^2/L_0+\|H\|_\infty/L_0)|Y|^2
\leq C_0h_*|Y|^2$.
Minimality yields $0\leq\Lambda\leq C_0h_*I$ on the full transverse
space, and reflection reduces its restriction to $E\cap n^\perp$ to a
symmetric two-by-two matrix.

Write $b=\Lambda_{pq}$, $d_0=\Lambda_{qq}$ and
$u_*=\Lambda_{pp}-\overline B_{pp}>0$. The terminal velocity
$v=\eta'-M\eta$ has active coordinates
\[
 \binom{v_p}{v_q}=
 \begin{pmatrix}u_*&b-\overline B_{pq}\\
 b-h_*-\overline B_{qp}&d_0-\overline B_{qq}\end{pmatrix}
 \binom{Y_p}{Y_q}.
\]
If $b\leq h_*/2$, take $Y=-p$. Then $v_p=-u_*\leq0$ and
$v_q=h_*+\overline B_{qp}-b\geq h_*/3$.
If $b>h_*/2$, positivity implies
$\Lambda_{pp}\geq b^2/d_0\geq h_*/(4C_0)$.
Take $Y=q-(b-\overline B_{pq})p/u_*$. Its norm is bounded by a
constant depending on $C_0$, and $v_p=0$. Moreover
\[
 v_q=\left(d_0-\frac{b^2}{u_*}\right)-\overline B_{qq}
 +\frac{bh_*+b(\overline B_{pq}+\overline B_{qp})
       -h_*\overline B_{pq}-\overline B_{qp}\overline B_{pq}}{u_*}.
\]
The first parenthesized term is nonnegative because $u_*\geq\Lambda_{pp}$.
The numerator is at least $h_*^2/2-C(C_0+1)\zeta h_*^2$.
Since $u_*\leq(C_0+\zeta)h_*$, choosing $\zeta$ small gives
$v_q\geq c h_*$. Divide the displacement, endpoint and velocity by
this positive terminal value. Both cases give the asserted normalization
and endpoint bound. Reflection gives vanishing $E^\perp$ components
before and after normalization. At $t_0=0$, prescribe $v=q$ directly.
\end{proof}

\begin{lemma}[The scalar growing solution]\label{lem:scalar-growth}
Let $\beta>0$ be sufficiently small, and set $P_0=\beta\tau^2$,
$Q_0=-2\beta\tau$ and $D=1+P_0^2$. The equation
\[
 (DV')'=2(1-\beta P_0)V,\qquad V_\lambda(0)=1,
 \qquad V_\lambda'(0)=\lambda\geq0
\]
has a positive solution on every target interval under consideration, and
$V_\lambda(x=1)\geq e^{b/\sqrt\beta}$. For $x\geq1$ put
$r_\lambda=-V_\lambda'/V_\lambda$. There is a bounded nonnegative $z$
such that
\[
 r_\lambda=\frac{\sqrt\beta}{x}-\frac z{x^2},\qquad
 z^2=\frac2{1+x^{-4}}+O(\sqrt\beta)\quad(x\geq2).
\]
Writing $g(\tau)=V_0(\tau)$, the two-component propagator for
$(-V',V)$ obeys
\[ \|\Phi_0(t,s)\|\leq C\Theta^8g(t)/g(s). \]
The estimates are dimension independent.
\end{lemma}
\begin{proof}
On $x\leq1$, integrating twice with $D\leq2$ and
$2(1-\beta P_0)\geq1$ gives
$V_\lambda\geq1+\lambda\tau/2+\tfrac12\int_0^\tau(\tau-s)V_\lambda(s)ds$.
Positive iteration bounds it below by $\cosh(\tau/\sqrt2)$.
Its logarithmic derivative $l$ satisfies $0\leq l$ and
$l'\leq2-l^2$, whence $l\leq\sqrt2\coth(\sqrt2\tau)$.
For $x\geq1$, write $\widehat V(x)=V(x/\sqrt\beta)$
and set $s=1-\rho$, $\rho=x^{-1}$ and
$f(\rho)=\widehat V(1/\rho)/\rho$. Direct substitution gives
\[
 ((1+\rho^4)f_\rho)_\rho=(2/\beta-2\rho^2)f.
\]
Starting from $f(1)>0$, $f_\rho(1)<0$, backward integration preserves
both signs. For $z=-\sqrt\beta f_\rho/f$ and
$\mu=\sqrt{(2-2\beta\rho^2)/(1+\rho^4)}$,
\[
 \sqrt\beta z_s=\mu^2-z^2+
               \sqrt\beta\frac{4\rho^3}{1+\rho^4}z.
\]
This gives a uniform upper bound and preserves nonnegativity.
Indeed its initial value at $\rho=1$ is
$z(1)=\sqrt\beta+V_\lambda'(1/\sqrt\beta)/V_\lambda(1/\sqrt\beta)$,
which is bounded independently of $\lambda$ by the preceding
logarithmic-derivative estimate.
The equation for $z-\mu$ has damping coefficient $z+\mu\geq c>0$
and a bounded forcing multiplied by $\sqrt\beta$. Its integrating
factor gives $|z-\mu|\leq Ce^{-c(1-\rho)/\sqrt\beta}+C\sqrt\beta$,
which proves the direction formula for $\rho\leq1/2$.

Reduction of order gives
$V_\lambda=g(1+\lambda I)$, where $I(\tau)=\int_0^\tau(Dg^2)^{-1}$
and $I(1)\geq c>0$. Positivity before $x=1$, and increase of $xg$
afterwards, give $g(s)/g(u)\leq C\Theta$ for $s\leq u$.
For arbitrary scalar data at $s$,
\[
 V(t)=\frac{g(t)}{g(s)}\left[
 V(s)+D(s)\left(V'(s)-\frac{g'(s)}{g(s)}V(s)\right)
 \int_s^t\frac{(g(s)/g(u))^2}{D(u)}\,du\right].
\]
The integral is at most $C\Theta^3$, $D(s)\leq C\Theta^4$, and
$|g'/g|\leq C$. Near $\tau=0$ the latter bound follows from
$g'(0)=0$ and $l'\leq2-l^2$; away from zero it follows from the
displayed bounds for $l$ and $z$.
Differentiating this formula introduces at most one additional
power of $\Theta$, proving the propagator assertion. This scalar
calculation follows~\cite[\S\S4.4--4.5]{OAI}.
\end{proof}

The selected initial conditions at activation are
$m(t_0)=s_0n(t_0)$, $v_q(t_0)=1$, $v_n(t_0)=0$ and
$-C\leq v_p(t_0)\leq0$. At the central label express
$m=s_0(Pp+\epsilon Qq+Nn)$ and
$v=\epsilon Up+Vq+\epsilon Wn$. The moving-frame connection has
$(S_f)_{12}=B_{pq}$, $(S_f)_{13}=B_{pn}$ and
$(S_f)_{32}=-B_{nq}$; skew symmetry determines its other entries.
Equations~\eqref{eq:ray-velocity} then give
\[
 (P,Q,N)=(\beta\tau^2,-2\beta\tau,1)+O(e\Theta^5),
 \qquad W=-(PU+QV)/N,\quad N\geq1/2.
\]
The reference ray system is $P'=-Q$, $Q'=-2\beta N$, $N'=0$.
Its triangular propagator is bounded by $C\Theta^2$. The error integral
is at most $Ce\Theta^5+Ce\Theta^3\sup|\mathrm{error}|$, and the last
term is absorbed by \eqref{eq:geometric-error}. Substitution in the
velocity equation gives a coefficient error at most $Ce\Theta^{13}$
from the two-component ideal system
\[
 U'=-2V+\frac{2P_0((P_0+\beta)V+Q_0U)}{1+P_0^2},\qquad V'=-U.
\]
Indeed the only unsuppressed contributions to
$J=(m\cdot Mv)/(s_0a)$ are $PB_{pq}V/a$, $h\epsilon^2QU/a$
and $NB_{nq}V/a$. The identity $\dot h/h=-B_{pp}-B_{qq}$ gives
$h/h_*=1+O(\epsilon\Theta G_*)$ on this interval, so
$h\epsilon^2/a=1+O(\epsilon\Theta G_*)$. The remaining contributions
contain either $\epsilon G_*$ or a term controlled by
\eqref{eq:geometric-error}. Transversality bounds
$|W|\leq C\Theta^2(|U|+|V|)$; the denominator $|m|^2/s_0^2$
is at least $1/4$. These estimates bound every rational coefficient
independently of the trace of the active block.
The full matrix estimate, including all transverse coordinates, is
derived in Section~\ref{sec:parameter-envelope}.

Apply Duhamel's formula with weight $g(s)/g(t)$.
Lemma~\ref{lem:scalar-growth} bounds its error operator by $Ce\Theta^{22}$. A geometric series therefore
gives the active propagator bound $2C\Theta^8g(t)/g(s)$.
For the selected data, $\lambda=-v_p(t_0)/\epsilon\geq0$,
\[
 |U-U_\lambda|+|V-V_\lambda|
       \leq Ce\Theta^{30}(1+\lambda)g.
\]
For $\tau\geq1$, $V_\lambda\geq c(1+\lambda)g$, so this is a
relative error. In particular
\begin{equation}\label{eq:positive-flux}
 J/V=P_0+\beta+Q_0r_\lambda+O(e\Theta^{34})>0.
\end{equation}
For $x\leq1$ its leading term is at least $\beta$, since
$Q_0,r_\lambda\leq0$. For $x\geq1$ it equals
$x^2-\beta+2\sqrt\beta z/x$. For $x\leq1$, the error is at most
$C c_d\Theta^{-6}\leq C c_d\beta^3$, hence less than $\beta/2$ after
fixing $c_d$ sufficiently small. For $x\geq1$ the leading term is
bounded below by $1-\beta$, and the same error bound suffices.

\begin{lemma}[Central frame renewal]\label{lem:central-renewal}
At the target set $p_2=m/|m|$, $q_2=v/|v|$, $n_2=p_2\times q_2$,
$a_2=p_2\cdot Mq_2$ and $\beta_2=(n_2\cdot Mq_2)/a_2$. Then
\begin{align*}
 a_2/a&=1+O(x_{\rm tar}^{-4}+\beta x_{\rm tar}^{-2}
       +\sqrt\beta x_{\rm tar}^{-3}+e\Theta^{40}),\\
 \beta_2x_{\rm tar}^2&=1+O(\sqrt\beta+x_{\rm tar}^{-4}+e\Theta^{40}),\\
 p_2\cdot Mp_2&\leq-\frac{c\sqrt{h_*}}{x_{\rm prev}x_{\rm tar}}
                                     +C(G_*+E_*).
\end{align*}
Moreover $A_{\rm tar}=|m||v|$ at the target satisfies
$A_{\rm tar}\geq K_h^{-c}\Theta^{-c}e^{b x_{\rm prev}}$.
\end{lemma}
\begin{proof}
Put $r=U/V$, $w=W/V$, $D_a=|m|^2/s_0^2$ and
$E_v=1+\epsilon^2(r^2+w^2)$. Normalize the two vectors directly:
\[
 a_2/a=\frac{J/V}{\sqrt{D_a}\sqrt{E_v}},\quad
 \beta_2=\frac{[(P,\epsilon Q,N)\times(\epsilon r,1,\epsilon w)]
                    \cdot(M(v/V)/a)}{(J/V)\sqrt{E_v}}.
\]
In the second numerator substitute
$P=x^2$, $Q=-2\sqrt\beta x$, $r=\sqrt\beta/x-z/x^2$.
The terms of size $\beta x^2$ and $\sqrt\beta xz$ cancel, leaving
\[
 -1+(1+x^{-4})z^2+\beta x^{-2}-2\sqrt\beta z x^{-3}
       =1+O(\sqrt\beta).
\]
The perturbation contributes at most $Ce\Theta^{36}$.
Also $J/V=x^2-\beta+2\sqrt\beta z/x+O(e\Theta^{34})$,
$D_a=1+x^4+O(e\Theta^7)$ and $E_v=1+O(\epsilon^2\Theta^4)$.
These identities prove the first two estimates. The rank-one part of
\eqref{eq:central-decomposition} gives
$p_2\cdot Mp_2=h\epsilon QP/D_a+O(G_*+E_*)$.
Since $QP\leq-c\sqrt\beta x^3$, $D_a\leq Cx^4$ and
$h\epsilon\asymp\sqrt{h_*}$, this proves compression.
Finally $V_\lambda(x_{\rm tar}/\sqrt\beta)\geq x_{\rm tar}^{-1}
V_\lambda(1/\sqrt\beta)\geq x_{\rm tar}^{-1}e^{b/\sqrt\beta}$ and
$s_0^{-1}\leq K_h^c$; the relative-error estimates give the lower bound
for $A_{\rm tar}$.
\end{proof}

Let $\mathcal P\geq1$ bound the source quantities
$K_h,\Theta,\epsilon^{-1},h_{\rm child},\delta^{-1}$ and the fixed
coefficient bounds and inverse estimates. These quantities are estimated in
Section~\ref{sec:parameter-envelope} before the frequency is selected.
The choice $\alpha=\delta h_{\rm child}/A_{\rm tar}$ therefore satisfies
$\alpha\leq\mathcal P^c e^{-b x_{\rm prev}}$.
At the central target the cutoff equals one, the phase is zero and
$f_\delta'(0)=\delta^{-1}$, so the leading gradient is exactly
$\alpha\delta^{-1}v\otimes m=h_{\rm child}q_2p_2^T$.
Section~\ref{sec:parameter-envelope} proves the corresponding full
transverse propagator and pressure-sign estimates in a neighborhood of
the central label. The following consequence of that comparison gives
a bound uniform over the stages of the iteration.

\begin{lemma}[Uniform target-size comparison]\label{lem:uniform-packet-size}
For the selected primary, on every label in the quantitative packet
neighborhood, the full ambient relative estimates of
Section~\ref{sec:parameter-envelope} give
\[
 \frac{|m(t,y)||v(t,y)|}{A_{\rm tar}}\leq C_d
       \qquad(1\leq\tau\leq\tau_{\rm end}).
\]
On the history interval and on $0\leq\tau<1$ the same ratio is at most
$K_h^{c_d}\Theta^{c_d}\epsilon^{-c_d}e^{-b x_{\rm prev}}$.
For $g(\tau)=V_0(\tau)$ and
$\alpha=\delta h_{\rm child}/A_{\rm tar}$, one also has
$\sup\alpha g\leq\mathcal P^{c_d}$ and
$\alpha\leq\mathcal P^{c_d}e^{-b x_{\rm prev}}$.
\end{lemma}
\begin{proof}
The full relative comparison, including the $E^\perp$ components, and
$\epsilon\Theta^2\ll1$ imply, for $\tau\geq1$,
\[
 |m(t,y)||v(t,y)|\asymp_d
       s_0\sqrt{1+x^4}\,V_\lambda(x/\sqrt\beta).
\]
The small perturbations of the initial label and of $s_0$ are included
in the same constants. At the central target this equivalence bounds
$A_{\rm tar}$ from below by a fixed multiple of the right side.
These constants do not grow with $\Theta$ or the stage.

For $0\leq x\leq1$, $V_\lambda$ is increasing. For $x\geq1$ the
reciprocal-variable proof of Lemma~\ref{lem:scalar-growth} shows that
$xV_\lambda(x/\sqrt\beta)$ is increasing. Hence for
$1\leq x\leq x_{\rm tar}$,
\[
 (1+x^2)V_\lambda(x/\sqrt\beta)
 \leq (x+x^{-1})x_{\rm tar}
             V_\lambda(x_{\rm tar}/\sqrt\beta)
 \leq2x_{\rm tar}^2V_\lambda(x_{\rm tar}/\sqrt\beta).
\]
For $x\leq1$ the same conclusion follows by comparison first with
$V_\lambda(1/\sqrt\beta)$ and then with the target. Since
$\sqrt{1+x^4}\asymp1+x^2$, this proves the constant size comparison
up to the target. Afterwards the logarithmic derivative in $\tau$ of
$\sqrt{1+x^4}V_\lambda$ is
\[
 \frac{2\sqrt\beta x^3}{1+x^4}-\frac{\sqrt\beta}{x}
                                          +\frac z{x^2},
\]
which is bounded for $x\geq x_{\rm tar}\geq2$. The retained scaled
interval has length at most $\Theta^{-60}$, so the additional factor is at
most two. The full relative errors preserve this fixed bound.

The stationary history inverse retains the affine endpoint contribution.
If $C_{\rm hist}$ is its polynomial norm from continuous forcing to
continuous high velocity, the affine-plus-forced construction gives
\[
 \sup_{0\leq t\leq t_0}|b_Y(t,y)|
 \leq \bigl(2t_0^{-1}\|F_t\|_\infty C_{\rm hist}
                 +t_0^{-1}\|F\|_\infty\bigr)|Y|.
\]
Here $Y$ is the fixed transverse coordinate, so
$|Y|\leq\|F(t_0,0)^{-1}\|\,|Y_{\rm phys}|\leq C K_h/h_*$.
The inverse norm is polynomial in the known frame quantities and
$t_0,t_0^{-1}$. Using $t_0\leq1$, $t_0^{-1}\leq h_*$ and
$h_*=a\epsilon^{-2}\leq2\epsilon^{-2}$ bounds this expression by
$K_h^{c_d}\epsilon^{-c_d}$ after enlarging the fixed exponent.
This also accounts for the generally nonzero initial high velocity.
On $\tau\leq1$, the scaled system has bounded
coefficients and initial norm at most $C(1+\lambda)$, with
$\lambda\leq C/\epsilon$; the coordinate transformation is therefore bounded by
the displayed polynomial in the preceding-stage quantities. Divide these estimates by
the exponential target lower bound from
Lemma~\ref{lem:central-renewal} to obtain the early and history ratio.
Finally $V_0\leq V_\lambda$. The preceding monotonicity argument gives
$\sup g\leq C x_{\rm tar}V_\lambda(x_{\rm tar}/\sqrt\beta)$,
whereas
$A_{\rm tar}\geq c_d s_0x_{\rm tar}^2
V_\lambda(x_{\rm tar}/\sqrt\beta)$.
Their growing factors cancel in $\alpha g$. The remaining factors
$\delta,h_{\rm child},s_0^{-1},x_{\rm tar}$ have the fixed polynomial
source bounds. The target exponential lower bound similarly proves
the assertion for $\alpha$ itself.
\end{proof}
\section{The finite packet and its exact Euler correction}\label{sec:exact-packet}

Fix a smooth reference solution, called the parent, on the closed
interval $[0,T]$. Constants may depend on this solution, on $d$,
and on the prescribed seed. The higher-order high-frequency coefficients
are constructed by the variational history and forward-tail method of
\cite[\S3]{OAI}. All source constants are fixed before the packet frequency
is chosen. No uniform bound on these constants over an infinite sequence
of parents is assumed.

Let $U$ be the parent, let $P_U$ be its Newton pressure, and let $\Phi_t$ be
its volume-preserving flow, with $\Phi_0=\mathrm{id}$. In particle labels put
\begin{equation}\label{eq:packet-frame}
 F=D\Phi_t,\qquad M=(DU)\circ\Phi_t=F_tF^{-1},\qquad
 H=(D^2P_U)\circ\Phi_t=-F_{tt}F^{-1}.
\end{equation}
We assume that these flow and coefficient fields have bounded smooth
spatial derivatives and continuous coefficient paths on $[0,T]$.
For the quantitative assertions we assume a Gevrey bound on the known
coefficients $F,F^{-1},F_t,F_{tt}$:
\begin{equation}\label{eq:packet-known-jets}
 \sup_{t,y}\|D_y^j C(t,y)\|\le C_C R_C^j(j!)^2,
 \qquad C\in\{F,F^{-1},F_t,F_{tt}\},\quad j\ge0.
\end{equation}
These conditions concern the parent. The lower singular-value bound of
$F$ follows from its inverse.
All derivative norms in this section are operator norms unless specified
otherwise.

Fix a unit vector $m_0\in\R^d$ and set
\[
 m(t,y)=F(t,y)^{-\top}m_0.
\]
Then $m\ne0$ and $m_t=-M^{\top}m$. The phase variable $\theta$ belongs
to $\mathbb T_P=\R/P\mathbb Z$, with $P=2\pi$ when a specific periodic seed
is used. Every vector field below has $d$ components, and the transverse
space $m^\perp$ has dimension $d-1$.

For physical Sobolev persistence one may use
$s_d=\lfloor d/2\rfloor+3>d/2+1$. Packet
estimates on the cylinder $\R^d\times\mathbb T_P$ use the fixed base order
\begin{equation}\label{eq:packet-base-index}
 b_d=2(d+2)+1=2d+5.
\end{equation}
The cylinder embedding, product, and commutator estimates at this order
follow by Fourier series in $\theta$ and the ordinary Sobolev estimates in
$y$. All their constants depend on $d$ and $P$, not on the number of
additional derivatives.

\subsection{The mean inverse}
Let $\mathcal H=L^2_\sigma(\R^d;\R^d)$ be the closed solenoidal subspace.
Choose a nonnegative even smooth cutoff $\chi$, equal to one near zero and
supported in the unit ball, and put $\chi_\ell(y)=\chi(\ell y)$.
There is a bounded map from the homogeneous first-order tensor space to
$L^2$ defined first on compact smooth tensors by
\[
 (T_{\chi_\ell}Q)_j
   =\sum_{i=1}^d\partial_i\{\chi_\ell(Q_{ij}-Q_{ji})\},
 \qquad A_{\chi_\ell}=T_{\chi_\ell}T_{\chi_\ell}^{*}.
\]
The tensor space here is the completion for the square sum of all
$\partial_kQ_{ij}$, with $1\le i,j,k\le d$. Hardy's inequality for
$d\ge3$ controls the term $(\partial_i\chi_\ell)Q_{ij}$ and extends $T$
continuously. Its bound is uniform for $0<\ell\le1$ after dilation.
Consequently $A_{\chi_\ell}$ is positive and symmetric, its range is
solenoidal, and every output vanishes off $\operatorname{supp}\chi_\ell$.
The construction requires only boundedness and spatial support of the
outputs.

We record the localization inequality used below. There are constants
$C_{1,d},C_{2,d}$ and a fixed $a_d>0$ such that, if
\begin{align*}
 \langle M(0,y)z,z\rangle&\ge-B_e|z|^2
       &&\text{for }|y|\ge a_dr/\ell,\\
 \langle M(0,y)z,z\rangle&\ge-B_c|z|^2
       &&\text{for }|y|<a_dr/\ell,
\end{align*}
where $0\le r\le1/4$, then, for $z\in\mathcal H$ and
$\mathcal L\ge C_{1,d}B_c$,
\begin{equation}\label{eq:packet-boundary-lower}
 \langle M(0)z,z\rangle_{L^2}
 +\mathcal L\langle A_{\chi_\ell}z,z\rangle_{L^2}
 \ge-(B_e+C_{2,d}B_cr^d)\|z\|_2^2.
\end{equation}
Here and below constants absorb the fixed normalization of the
cutoff plateau. To obtain \eqref{eq:packet-boundary-lower}, set
$Q=T_{\chi_\ell}^{*}z$. On the plateau, subtract one half of the uncut
antisymmetric divergence of $Q$ from $z$. The difference is weakly
harmonic, by the defining Riesz identity for $Q$. Its $L^2$ mass on the
concentric ball of relative radius $r$ is bounded by $C_dr^d$ times its
mass on the plateau ball. This follows from the interior harmonic
estimate, obtained by nested cutoff energy estimates and Sobolev
embedding, and then by dilation. The other part is controlled by
$\|Q\|_{\dot H^1}^2=\langle A_{\chi_\ell}z,z\rangle$.
Combining this local bound with the exterior lower bound proves the
inequality in every dimension $d\ge3$.

Assume $H(t,y)\le K I$ in quadratic-form order and
\begin{equation}\label{eq:packet-mean-guard}
 K\frac{T^2}{2}+(B_e+C_{2,d}B_cr^d)T\le\frac12.
\end{equation}
For $v\in L^2(0,T;\mathcal H)$ define
\[
 Iv(t)=-\int_t^T v(s)\,ds,\qquad X_v=F Iv,\qquad
 D_Fv=\partial_t X_v=F_tIv+Fv.
\]
The mean form is
\begin{align}
 \mathfrak a(v,w)
 &=\int_0^T\langle D_Fv,D_Fw\rangle\,dt
   -\int_0^T\langle HX_v,X_w\rangle\,dt\notag\\
 &\quad+\langle (M(0)+\mathcal L A_{\chi_\ell})X_v(0),
                         X_w(0)\rangle.\label{eq:packet-mean-form}
\end{align}
Spatial integrations are included in each inner product. The form need
not be symmetric because $M(0)$ need not be symmetric. It is bounded and
coercive. Indeed $X_v(T)=0$, hence
\[
 \|X_v\|_{L^2_tL^2_y}^2\le\frac{T^2}{2}\|D_Fv\|_{L^2_tL^2_y}^2,
 \qquad \|X_v(0)\|_2^2\le T\|D_Fv\|_{L^2_tL^2_y}^2.
\]
Moreover $X_v(0)=Iv(0)\in\mathcal H$, because $F(0)=I$.
Equations \eqref{eq:packet-boundary-lower} and
\eqref{eq:packet-mean-guard} give
$\mathfrak a(v,v)\ge\frac12\|D_Fv\|_2^2$.
The inverse frame gives a polynomial norm comparison. If
$I=\|F^{-1}\|_\infty$ and $C_1=\|F_t\|_\infty$, then
$v=F^{-1}D_Fv-F^{-1}F_tF^{-1}X_v$, so
$\|v\|_2\le I(1+C_1IT)\|D_Fv\|_2$.
Thus the Hilbert-space coercive inverse constructs the unique $v$ with
\begin{equation}\label{eq:packet-mean-weak}
 \mathfrak a(v,w)=-\int_0^T\langle f,X_w\rangle\,dt
 \quad\text{for every }w\in L^2(0,T;\mathcal H).
\end{equation}
Set $B=Fv$.

\begin{lemma}[Strong mean and its pressure]\label{lem:packet-strong-mean}
For a forcing $f$ with continuous $L^2$ spatial derivatives of every
order, the preceding construction gives continuous paths $B$ and $B_t$
in every spatial Sobolev space and a spatially smooth scalar $q$,
whose spatial derivatives depend continuously on time locally in space,
normalized by $q(t,0)=0$, such that
\begin{equation}\label{eq:packet-mean-pde}
 B_t+MB+F^{-\top}\nabla_yq=f,
 \qquad F^{-1}B\in\mathcal H.
\end{equation}
These identities hold pointwise, with one-sided time derivatives at both
endpoints. In addition, there exists $z_0\in\mathcal H$ such that
\begin{equation}\label{eq:packet-mean-initial}
 B(0)=\mathcal L A_{\chi_\ell}z_0,
 \qquad\operatorname{supp}B(0)\subset\overline B(0,\ell^{-1}).
\end{equation}
\end{lemma}
\begin{proof}
Translate the coefficients and forcing in the spatial variable. The
coercive inverse is differentiable under each translation; differentiating
its inverse identity expresses each new derivative using previously
constructed derivatives and the known coefficient derivatives. This
constructs the spatial derivatives in $L^2$. The same argument applies to
the translated cutoff operator on the fixed homogeneous tensor space.
The argument yields derivatives of every finite order.

Let $Q_t:\mathcal H\to L^2$ be $Q_tz=F(t)z$. Its Gram operator
$Q_t^*Q_t$ has a positive lower bound from $F^{-1}$. The weak equation,
$F_{tt}=-HF$, and integration by parts give
\[
 v_t=(Q_t^*Q_t)^{-1}Q_t^*\{f-2F_tv\}.
\]
First this is an $L^2$ time identity. The $H^1$ trace construction and
continuity of the coefficients then give continuous representatives of
$v$ and $v_t$, with the same spatial derivatives. The residual
$r=f-2F_tv-Fv_t$ satisfies $Q_t^*r=0$; equivalently,
$F^{\top}r$ lies in the closed $L^2$ gradient space. Its smooth
representative is curl-free; hence
\[
 q(t,y)=\int_0^1\langle F(t,sy)^{\top}r(t,sy),y\rangle\,ds
\]
satisfies $q(t,0)=0$ and $\nabla q=F^{\top}r$. This proves
\eqref{eq:packet-mean-pde}. The pointwise equations follow from equality
of continuous representatives of the same $L^2$ field. The Gram
construction and closedness give $v(t)\in\mathcal H$ at every time in
the closed interval. In particular, $B(0)=v(0)\in\mathcal H$. With $z_0=Iv(0)$,
the boundary identity in \eqref{eq:packet-mean-weak} is
\[
 P_\sigma\{-F_t(0)z_0-B(0)+(M(0)+\mathcal L A_{\chi_\ell})z_0\}=0.
\]
Since $F_t(0)=M(0)$ and both $B(0)$ and $A_{\chi_\ell}z_0$ are
solenoidal, this yields \eqref{eq:packet-mean-initial} as an $L^2$
identity. Spatial continuity gives the asserted pointwise support.
\end{proof}

For a transverse
endpoint problem the map $Q_t$ is replaced by $F(t)R_\perp$, where
$R_\perp:\R^{d-1}\to m_0^\perp$ is a linear isometry onto the whole
orthogonal complement. Its Gram lower bound is again supplied by
$F^{-1}$. The energy, trace, and coercive-inverse arguments are unchanged;
there are $d-1$ transverse variables throughout.

\subsection{The joined high inverse and the full tensor corrector}
Fix an activation time $0<t_0<T$. Choose an isometry
$R_\perp:\R^{d-1}\to m_0^\perp$ onto the whole orthogonal complement
and put $Q=FR_\perp$. Its range is $m(t,y)^\perp$, and
$Q_t=MQ$, $Q_{tt}=-HQ$. The fixed coordinate space $\R^{d-1}$ is used at
every label and time.

For an angular-mean-zero forcing
$f\in C([0,T];H^q(\R^d\times\mathbb T_P;\R^d))$ for every finite $q$,
first solve on $[0,t_0]$ the following displacement problem. Find
$z\in H^1_0(0,t_0;L^2(\R^d\times\mathbb T_P;\R^{d-1}))$ such that
\begin{equation}\label{eq:packet-high-history}
 \int_0^{t_0}\!\bigl[\langle (Qz)_t,(Q\zeta)_t\rangle
            -\langle HQz,Q\zeta\rangle\bigr]dt
       =-\int_0^{t_0}\!\langle f,Q\zeta\rangle dt
 \quad(\zeta\in H^1_0).
\end{equation}
Both displacement traces are zero. This condition does not prescribe the
initial high velocity. Set $A=Qz_t$ on the history interval.

\begin{lemma}[The bounded joined high inverse]\label{lem:packet-joined-high}
Under the known frame bounds and $Kt_0^2/2\le1/2$, problem
\eqref{eq:packet-high-history} has a unique solution. Its velocity $A$,
its time derivative, and their spatial derivatives are continuous on
$[0,t_0]$. At $t_0$ its velocity trace can be continued by the
forward high equation. This gives a unique joined high operator on
$[0,T]$ with the specified history boundary conditions. The history
inverse and trace bounds are polynomial in the known coefficient norms,
inverse lower bound, $t_0$ and $t_0^{-1}$. No exponential forward
estimate on $[0,t_0]$ is used.
\end{lemma}
\begin{proof}
Write $X=Qz$. The one-sided Poincar\'e estimate and the curvature upper
bound give
\[
 \int_0^{t_0}(|X_t|^2-\langle HX,X\rangle)dt
       \ge\tfrac12\|X_t\|_{L^2}^2.
\]
If $Q^*Q\ge cI$, then $z=(Q^*Q)^{-1}Q^*X$, and direct differentiation,
using $X_t=Q_tz+Qz_t$, gives the polynomial estimate
\[
 \|z_t\|_2\le
 c^{-1/2}(1+t_0 c^{-1/2}\|Q_t\|_\infty)\|X_t\|_2.
\]
Thus the form is coercive on the fixed Hilbert space of coordinate
velocities $z_t$ with zero time integral. Its bounded coercive inverse
constructs $z$. Integration by parts and $Q_{tt}=-HQ$ give
\[
 (Q^*Q)z_{tt}=Q^*(f-2Q_tz_t).
\]
The positive Gram inverse therefore gives $z_{tt}$ and the
continuous $H^1$ traces. Repeated differentiation under spatial
translation is justified on this same fixed Hilbert space, exactly as
for the mean inverse. It gives all the stated derivatives and preserves
both homogeneous displacement boundary conditions.

At order zero, a physical-space estimate is available directly.
Put $M_* =\sup\|M\|$. Testing the form with $z$
gives $\|X_t\|_2\le2t_0\|f\|_2$. Since $A=X_t-MX$, it follows that
\[
 \|A\|_2\le2t_0(1+M_*t_0)\|f\|_2.
\]
The velocity equation derived below implies
$\|A_t\|_2\le3M_*\|A\|_2+\|f\|_2$. The elementary Hilbert-valued
trace estimate consequently gives
\begin{equation}\label{eq:packet-history-polynomial-trace}
 \|A\|_{C_tL^2}\le
 \bigl[2t_0(1+M_*t_0)(t_0^{-1/2}+3M_*\sqrt{t_0})
                         +\sqrt{t_0}\bigr]\|f\|_{L^2_tL^2}.
\end{equation}
This also bounds the initial datum for the forward segment.
\end{proof}

The Gram equation and tangency show that the same $A$ satisfies
\begin{equation}\label{eq:packet-high-ode}
 A_t=\left(-M+2\frac{m\otimes M^{\top}m}{|m|^2}\right)A
       +\left(I-\frac{m\otimes m}{|m|^2}\right)f.
\end{equation}
Here $(a\otimes b)v=a\langle b,v\rangle$. Indeed the residual
$f-A_t-MA$ is perpendicular to the whole range of $Q$, hence is a
multiple of $m$. Differentiating $m\cdot A=0$ gives its coefficient.
For $t\ge t_0$ use the fundamental evolution of
\eqref{eq:packet-high-ode} and Duhamel's formula with datum $A(t_0)$.
The velocities and their derivatives match at $t_0$ by the common
right-hand side, so the joined field is strongly differentiable there
as well as at the two closed endpoints.

Let $g$ be a positive continuous function with $g(t_0)=1$, and assume
the supported forward bound
\begin{equation}\label{eq:packet-tail-propagator}
 \|\mathcal U(t,s;y)\|\le C_{\rm f}\,g(t)/g(s),
 \qquad t_0\le s\le t\le T,
 \quad y\text{ in the fixed support neighborhood}.
\end{equation}
This hypothesis concerns the known parent evolution on the activation
tail. The polynomial history estimate is supplied separately by
\eqref{eq:packet-history-polynomial-trace}.

Define the normalized periodic primitive $\partial_\theta^{-1}$ and put
\begin{equation}\label{eq:packet-high-pressure}
 \pi=\partial_\theta^{-1}
       \frac{\langle m,f\rangle-2\langle m,MA\rangle}{|m|^2}.
\end{equation}
Then $A_t+MA+m\partial_\theta\pi=f$. The inverse, trace, and forward
maps commute with the angular integral. They therefore preserve zero
angular mean. They also commute with multiplication by a fixed spatial
cutoff, since these linear high equations contain no spatial derivative.
Their uniqueness proves support preservation and covariance. Thus the
joined high field and its pressure are smooth in $(y,\theta)$, periodic,
and supported where the forcing is supported. Their spatial derivatives
are continuous in time; the high velocity also has the continuous time
derivative given by \eqref{eq:packet-high-ode}.

The primary field uses a nonzero terminal displacement. To distinguish its
physical and coordinate forms, let the selected physical endpoint be
$Y_{\rm phys}\in m(t_0,0)^\perp$, and fix
\[
 Y=R_\perp^*F(t_0,0)^{-1}Y_{\rm phys}\in\R^{d-1}.
\]
Then $Q(t_0,0)Y=Y_{\rm phys}$ and
$|Y|\le\|F(t_0,0)^{-1}\|\,|Y_{\rm phys}|$. This same coordinate
vector is used at every label.

The cylinder problem has terminal datum
$\alpha\chi_{\rm in}(y)f_\delta(\theta)Y$, which belongs to every
cylinder Sobolev space, and initial displacement zero. Subtract the
affine lift $\alpha(t/t_0)\chi_{\rm in}f_\delta Y$. Using
$Q_{tt}=-HQ$, its remaining zero-trace problem is exactly
\eqref{eq:packet-high-history} with forcing
$-2\alpha t_0^{-1}\chi_{\rm in}f_\delta Q_tY$. This forcing is
continuous in every cylinder Sobolev space and has the prescribed
compact spatial support.

Equivalently, at each fixed label solve the finite-dimensional homogeneous
displacement problem with endpoints $0,Y$, and write $b_Y(t,y)=Qz_t$
for its velocity, continued homogeneously on the tail. The uniform frame
bounds justify this pointwise solution and its label derivatives. Linearity
and uniqueness identify the cylinder velocity with the compact
periodic primary
\begin{equation}\label{eq:packet-joined-primary}
 A_1(t,y,\theta)=\alpha\chi_{\rm in}(y)f_\delta(\theta)b_Y(t,y).
\end{equation}
Its initial value generally depends on $y$. The first stage,
whose activation time is zero, instead uses a prescribed homogeneous
forward datum. The same recursion applies in both cases. For higher grades
the displacement endpoints in \eqref{eq:packet-high-history} are both
zero, but their initial velocities are retained.

For every such $A$ define all $d^2$ entries of the potential
\begin{equation}\label{eq:packet-tensor-potential}
 Q_{ij}=\partial_\theta^{-1}
       \frac{m_iA_j-m_jA_i}{|m|^2},\qquad Q_{ij}=-Q_{ji},
 \qquad
 C_j=\sum_{i=1}^d D_yQ_{ij}[F^{-1}e_i].
\end{equation}
If $\kappa=k^{-1}$ and $\psi=\Phi_t^{-1}$, the chain rule gives
\begin{equation}\label{eq:packet-pair-curl}
 A_j(\psi(x),km_0\cdot\psi(x))+
 \kappa C_j(\psi(x),km_0\cdot\psi(x))
 =\kappa\sum_{i=1}^d\partial_{x_i}
    Q_{ij}(\psi(x),km_0\cdot\psi(x)).
\end{equation}
Indeed the phase part on the right is
$\sum_i m_i\partial_\theta Q_{ij}=A_j$, by tangency.
Taking divergence and interchanging the two ordinary derivatives cancels
all terms by antisymmetry. This identity holds for every pair $(i,j)$.
Applying the volume-preserving Piola transformation shows that
$F^{-1}(A+\kappa C)$ has zero divergence on every phase-shifted label
graph. Integration over the shift gives its membership in the cylinder's
closed divergence-free $L^2$ space. The potentials are smooth and have
the same compact spatial support as $A$, so the integrations by parts
are legitimate. The time derivative of $C$ is obtained by differentiating
\eqref{eq:packet-tensor-potential}: it includes the derivatives of $F^{-1}$,
$m$, $|m|^{-2}$, the primitive, and $A$. The same formula gives the
one-sided time derivatives at the endpoints.

\subsection{Recursive construction from preceding grades}
For vector fields $a,b$ on the time-label-angle cylinder write
\[
 S(a,b)=D_yb[F^{-1}a],\qquad
 T(a,b)=\langle m,a\rangle\partial_\theta b,\qquad
 L a=\partial_ta+Ma.
\]
Let $A_1$ be the homogeneous joined primary
\eqref{eq:packet-joined-primary} (or the first-stage forward primary).
Put $B_1=q_1=0$, form $C_1$ by
\eqref{eq:packet-tensor-potential}, and set every grade-zero field to zero.
Suppose grades below $p\ge2$ have been constructed. Set
\[
 V_i=A_i+B_i+C_{i-1}\quad(1\le i<p),\qquad
 \widehat V_p=C_{p-1},
\]
and extend this known sequence by zero above $p$. The forcing determined
by the preceding grades is
\begin{equation}\label{eq:packet-known-force}
 K_p=-LC_{p-1}-F^{-\top}\nabla_y\pi_{p-1}
       -\sum_{i+j=p}S(\widehat V_i,\widehat V_j)
       -\sum_{i+j=p+1}T(\widehat V_i,\widehat V_j).
\end{equation}
Every term uses only the strict prefix. Let an overline denote normalized
angular averaging. First solve the mean problem with forcing
$\overline K_p$, obtaining $B_p,q_p$. Then apply the joined high
inverse of Lemma~\ref{lem:packet-joined-high} with forcing
\begin{equation}\label{eq:packet-new-high-force}
 f_p=K_p-\overline K_p-\langle m,B_p\rangle\partial_\theta A_1,
\end{equation}
obtaining $A_p,\pi_p$, and form $C_p$ from $A_p$.
The final term of \eqref{eq:packet-new-high-force} has zero angular mean.
Terms in $K_p$ containing only means are angularly constant and disappear
from $K_p-\overline K_p$. Every remaining term contains a high field or
corrector, so it has the same compact spatial support as the seed. The
mean forcing has all continuous $L^2$ spatial derivatives by the
Sobolev product estimates. The high forcing is smooth, periodic, mean
zero, and spatially supported. These facts inductively verify the domains
of the two inverse operators.

In the quadratic
coefficient, replacing $\widehat V_p$ by
$\widehat V_p+A_p+B_p$ adds only
$\langle m,B_p\rangle\partial_\theta A_1$.
The other potentially new fast terms vanish because $A_p,A_1$ are
tangent and $B_p$ is angularly constant. This proves both strict-prefix
coherence and the grade-$p$ momentum equation. In particular, families
constructed for different finite cutoffs are restrictions of one family.

\subsection{Finite truncation, divergence, and the residual identity}
For $N\ge1$ define
\begin{align}
 V^{[N]}&=\sum_{i=1}^N\kappa^i(A_i+B_i+\kappa C_i),\notag\\
 p^{[N]}&=\sum_{i=1}^N\kappa^i(q_i+\kappa\pi_i).
                                                        \label{eq:packet-finite-pairs}
\end{align}
Equivalently the velocity grades are $V_i=A_i+B_i+C_{i-1}$ for
$1\le i\le N$, followed by the terminal grade $V_{N+1}=C_N$.
The pressure grades are $p_i=q_i+\pi_{i-1}$ for $1\le i\le N$,
followed by $p_{N+1}=\pi_N$. Grades outside this range are zero.
In particular the terminal grade does not include unsolved
$A_{N+1}$ or $B_{N+1}$.

Each high-corrector pair in \eqref{eq:packet-finite-pairs} is exactly
solenoidal after graph evaluation by \eqref{eq:packet-pair-curl}.
Each mean has $F^{-1}B_i\in\mathcal H$ by
Lemma~\ref{lem:packet-strong-mean}; its constant-angle embedding therefore
belongs to the same cylinder solenoidal space. Thus the complete finite
sum satisfies the closed divergence constraint. The constant-angle
embedding of an ordinary $L^2$ gradient lies in the cylinder's closed
gradient space: approximate it by gradients of compactly supported
smooth scalars and use $\nabla_\kappa\varphi=\kappa\nabla_y\varphi$
for an angle-independent scalar. This applies to the mean pressure
gradients. The compact periodic high scalar pressures themselves give
gradient generators. Their finite sum therefore belongs to the required
closed gradient space.

For any $V,p$ define the label momentum residual by
\begin{equation}\label{eq:packet-raw-residual}
 \mathcal R_\kappa(V,p)=LV+F^{-\top}\nabla_yp+
 \kappa^{-1}m\partial_\theta p+S(V,V)+\kappa^{-1}T(V,V).
\end{equation}
For the finite grade sequences above, let
\[
 R_n=LV_n+F^{-\top}\nabla_yp_n+m\partial_\theta p_{n+1}
       +\sum_{i+j=n}S(V_i,V_j)+\sum_{i+j=n+1}T(V_i,V_j).
\]
The mean and high equations give $R_n=0$ for $0\le n\le N$.
Expanding the finite sums with their derivatives, including the
one-sided time derivatives, consequently gives the exact identity
\begin{equation}\label{eq:packet-tail-identity}
 \mathcal R_\kappa(V^{[N]},p^{[N]})
   =\sum_{n=N+1}^{2N+2}\kappa^nR_n.
\end{equation}
This finite identity involves no time derivative of $p$.
In particular, $R_{N+1}$ retains $LC_N$ and
$F^{-\top}\nabla_y\pi_N$.

\subsection{Quantitative bounds at one known parent}
We give the estimates in a form that makes the derivative counting
explicit. For an integer $a\ge0$ set
$\mathcal G_a(j)=R^{j+a}((j+a)!)^2$.
Use ordered cylinder derivatives with alphabet $\{1,\ldots,d+1\}$.
More precisely, the block of order $j$ is
\[
 \operatorname{Block}_j(u)=
 \sum_{|w|=j}\sum_{l=0}^{b_d}\sum_{|v|=l}
     \|\partial^v\partial^w u\|_{C([0,T];L^2)}.
\]
At fixed time the corresponding word Sobolev norm is equivalent to the
usual $H^{b_d}$ norm, with fixed constants depending only on $d,b_d$.
A positive time weight is inserted by dividing the field by that
weight before taking spatial or angular derivatives. For one sufficiently
large $R$, depending only on the known coefficients, seed, and inverse
bounds, the preceding recursion gives
\begin{align}
 \operatorname{Block}_j(B_p,B_{p,t},\nabla_yq_p)
   &\le H_0^{2p-2}\mathcal G_{100p-140}(j),&&p\ge2,\notag\\
 \operatorname{Block}_j(A_p,A_{p,t},C_p,C_{p,t},\nabla_y\pi_p)
   &\le \gamma(t)H_0^{2p-2}\mathcal G_{100p-80}(j),&&p\ge1.
                                                        \label{eq:packet-grade-bounds}
\end{align}
Here each entry has the stated bound separately. We use
$\gamma(t)=\alpha$ for $0\le t\le t_0$ and
$\gamma(t)=\alpha g(t)$ on the tail, with
$H_0\ge\max(1,\sup\gamma)$. The notation with $\gamma(t)$ means the
uniform-time block of the field divided by $\gamma$. The supported
forward constant in \eqref{eq:packet-tail-propagator}, the history trace
bound, and the terminal datum are included among the known source
constants controlling $R$. In particular the higher initial high fields
are bounded with the factor $\alpha$, not set to zero. This is the
family used below and in Section~\ref{sec:initial-summability}.

For completeness, the derivative estimates underlying
\eqref{eq:packet-grade-bounds} are as follows. The word Leibniz rule is
bounded by a binomial convolution. Its Gevrey convolution closes because
\[
 \binom jl(l!)^2((j-l)!)^2
       =\frac{(j!)^2}{\binom jl},
 \qquad \sum_{l=0}^j {\binom{j}{l}}^{-1}\le3.
\]
Fixed-order Sobolev products add a constant depending on $b_d,d$.
On the history interval, differentiate the fixed-space coercive
inverse and then use the Gram equation and its trace. On the tail,
differentiate the Duhamel equation and use
\eqref{eq:packet-tail-propagator}. The history terminal trace is the
datum in this Volterra recurrence. The joined high velocity and its time
derivative require shifts of at most three grades. Differentiating the
coercive inverse, followed by its Gram reconstruction and the $H^1$ time
trace, requires shifts of at most three grades for the mean, its time
derivative, and its pressure gradient. These are estimates on word blocks with the full
alphabet. The angular mean is a bounded $L^2$ map with norm $P^{-1/2}$,
and the constant-angle embedding has norm $P^{1/2}$; the factors cancel.
Thus converting mean input to mean output does not multiply the unknown
word radius at each grade. Only the known coefficient radius is enlarged
once. Also the mean time weight is the constant $H_0^{2p-2}$, so no
commutation of a time-dependent weight with the mean inverse is claimed.
The history estimate is applied to $f/\alpha$ and the tail estimate
to $f/(\alpha g)$; linearity preserves the common factor $\alpha$.
The tail estimate uses the weighted Duhamel operator without differentiating
$g$. The tensor potential and its time derivative require one spatial
derivative and a fixed $d$-dependent multiplier bound; the pressure
source \eqref{eq:packet-high-pressure} and one slow derivative require
shifts of at most two additional grades. Enlarging $R$ once absorbs these
fixed constants.
The grade counting can be made explicit. The mean force has shift at
most $100p-150$, and the high force has shift at most $100p-90$.
For example a slow high--high term with $i+j=p$ has shift
$100p-159$. In a fast term with $i+j=p+1$, the first high factor is
annihilated by $m\cdot A_i=0$. A first corrector factor has shift
$100i-180$, so the resulting high term has shift at most $100p-159$;
a first mean factor gives shift $100p-119$, and has zero angular average.
The new term $(m\cdot B_p)\partial_\theta A_1$ has the latter shift as
well. These shifts satisfy the stated forcing bounds. The profile powers close
because $\gamma\le H_0$; in particular a mean--high fast product has
the required factor $\gamma H_0^{2p-2}$.
The at most quadratic number of summands is absorbed by the additional
factorial shifts: for $a\ge0$,
$\mathcal G_{a+2}(j)/\mathcal G_a(j)
 =R^2(j+a+1)^2(j+a+2)^2$.
The low grades satisfy the bounds after one enlargement of the known radius.
The mean inverse and the joined high inverse, followed by the corrector
and pressure estimates, require fewer than ten additional shifts.
Strong induction proves \eqref{eq:packet-grade-bounds} for this
joined family using only previously constructed fields.

For the correction use the truncated weighted norm
\[
 \|u\|_{r,b_d;J}
 =\sum_{j=0}^J\frac{r^j}{(j!)^2}
     \sum_{|w|=j}\sum_{l=0}^{b_d}\sum_{|v|=l}\|\partial^v\partial^wu\|_2.
\]
A block bound
$C\mathcal G_a(j)$ gives a weighted norm bound of at most
$2C(4R)^a(a!)^2$ when $4Rr\le1/2$.
This conversion is performed after the word recursion, not at every
recursive step. In particular, at every derivative cutoff,
\[
 \|V_i\|_{r,b_d;J}\le
  6H_0^{2i}(4R)^{100i}((100i)!)^2\quad(i\ge1),
\]
with analogous bounds for the time derivatives and pressure
forces at a smaller fixed radius. To quantify the dependence on the
expansion order, choose a known-source constant
$C_0\ge1$ and put $S_N=C_0(N+1)^{200}$. The inequality
$(100i)!\le(100(N+1))^{100i}$ for $i\le N+1$ bounds each displayed
grade by $S_N^i$ after enlarging $C_0$. The time derivatives
and pressure terms obey the same bound; one slow or angular derivative
is controlled by one fixed radius reduction, whose reciprocal factor is
included in $C_0$. A slow product with $i+j=n$ is bounded by
$C S_N^n$, and a fast product with $i+j=n+1$ by $C S_N^{n+1}$.
There are at most
$(N+2)^2$ terms. The terminal linear terms have the same bound, including
the $C_N$ time derivative and the slow derivative of $\pi_N$.
Thus $C_1(N+2)^2S_N^{n+1}$ bounds the entire residual grade. Taking
$B_N=C_2(N+1)^{202}$ with a sufficiently large fixed $C_2$ absorbs this
factor for every $n\ge N+1$. Thus the fixed tail degree is $202$, and
\begin{equation}\label{eq:packet-tail-base}
 B_N\le C_*(N+1)^{202},\qquad
 \|R_n\|_{r,b_d;J}\le B_N^{n+1}\quad(N+1\le n\le2N+2).
\end{equation}
The constant $C_*$ is fixed by the current parent and seed, and is
independent of $N,k,J$. Spatial derivatives of the high pressure, the
time derivative of $C_N$, and the terminal linear terms are all included
in this estimate.

Put
\begin{equation}\label{eq:packet-coordinate-change}
 Z=kF^{-1}V^{[N]},\quad Q^a=k^2p^{[N]},\quad
 \nabla_\kappa=\kappa\nabla_y+m_0\partial_\theta,
 \quad G=F^{-1}F^{-\top}.
\end{equation}
Then the coordinate residual is $\mathcal E=kF^{-1}\mathcal
R_\kappa(V^{[N]},p^{[N]})$. If $B_N\le k^{1/100}$ and the known inverse
multiplier bound is at most $k^{1/100}$, the geometric sum in
\eqref{eq:packet-tail-identity} yields, for $N\ge X-1$, $X\ge6$,
\begin{equation}\label{eq:packet-small-residual}
 \|\mathcal E\|_{r,b_d;J}
       \le \exp(-\tfrac7{10}X\log k).
\end{equation}
The field $Z$ and its first cylinder derivatives have fixed weighted
bounds $B_0,B_1$. More sharply, its cylinder transport vector is
\[
 \mathcal V_\kappa Z=(\kappa Z,m_0\cdot Z),
 \qquad \|\mathcal V_\kappa Z\|_{r,b_d;J}\le D_*/k.
\]
The order-one high field is tangent, so its fast transport component
vanishes exactly. All other contributions carry at least one inverse
frequency. The radius estimate uses this small transport norm rather than the
full $O(1)$ norm of $Z$.

\subsection{An exact correction on the cylinder}
The coordinate equation is
\begin{equation}\label{eq:packet-coordinate-equation}
 W_t+(W\cdot\nabla_\kappa)W+2F^{-1}F_tW
    +\kappa\sum_{i=1}^d W_iF^{-1}(\partial_iF)W
    +G\mathbf p=0,
 \qquad \nabla_\kappa\cdot W=0.
\end{equation}
Here $\mathbf p$ is a vector pressure variable in
$\mathscr G_\kappa=\overline{\{\nabla_\kappa\varphi:
\varphi\text{ smooth, periodic, compact in }y\}}^{L^2}$.
The construction requires only membership in this closed space, without
a global periodic scalar potential or a closed-range property for the
degenerate $d$-component operator on the $(d+1)$-dimensional cylinder.
The equation contains all $d$ quadratic coefficient matrices.
Substituting $Z$ and the approximate pressure vector
$\mathbf p^a=\nabla_\kappa Q^a$ gives exactly the residual $\mathcal E$
above. The correction pressure is solved in this closed subspace of
vector-valued cylinder $L^2$.
The positive matrix $G$ gives a coercive variational pressure inverse on
that subspace. Its derivatives are controlled by the coefficient
jets and commutators.

Seek $W=Z+E$ with $E(0)=0$. Use the inverse metric
$G^{-1}=F^{\top}F$ in the finite weighted Sobolev energy. The top
pressure term cancels since
$\langle G^{-1}E,G\mathbf p^e\rangle
 =\langle E,\mathbf p^e\rangle=0$ for $\mathbf p^e\in\mathscr G_\kappa$.
The top transport term is integrated by parts using the closed
solenoidal constraint. Differentiated metric terms and lower products
have a known coefficient bound. To specify the energy convention, put
$\|u\|_{F^{\top}F}=(\int\langle F^{\top}Fu,u\rangle)^{1/2}$ and set
\[
 \mathcal A_J=\sum_{j=0}^J\frac{r^j}{(j!)^2}
   \sum_{|w|=j}\sum_{l=0}^{b_d}\sum_{|v|=l}
       \|\partial^v\partial^w E\|_{F^{\top}F}.
\]
The quantity $\mathcal A_J$ is a sum of metric norms and is
uniformly equivalent to the preceding weighted norm by the known frame
bounds. Let $\mathcal L_J$ be the same sum with the additional factor
$j$ in every summand. Derivatives of norms at zero can be justified by
first replacing each norm by its quadratic regularization and then
letting the regularization parameter decrease to zero.

The orthogonal gradient projection commutes with every cylinder
derivative. In the differentiated pressure equation, each nontrivial
commutator with $G$ lowers the derivative order on the pressure by at
least one. This compensates for the one derivative in the transport
source. After the fixed-base Sobolev pressure estimate, its remaining
coefficient-word convolutions are absorbed by choosing the known radius
small enough. Thus no unknown derivative of order $J+b_d+1$ is left in
the energy bound. The word product convolution and the metric transport
cancellation consequently give a single radius-loss factor. In fact, in
a transport commutator with $1\le j\le n$, put $l=n-j+1$. The ratio of
Gevrey weights, including its Leibniz coefficient, is
\[
 \frac{r^n}{(n!)^2}\binom nj
   \frac{(j!)^2(l!)^2}{r^j r^l}
   =\frac{l^2}{r\binom nj}\le\frac lr,
\]
since $\binom nj\ge l$. The excluded term $j=0$ is precisely the
cancelled principal transport term. Thus the loss is the factor $l$
in $\mathcal L_J$, with no additional derivative beyond the cutoff. We obtain
\begin{equation}\label{eq:packet-correction-energy}
 \mathcal A_J'
 \le C\{\mathcal A_J+\mathcal A_J^2+\varepsilon\}
  +\left\{\frac{r'}{r}
       +C(r^{-1}+R_c)(D_*/k+\mathcal A_J)\right\}\mathcal L_J.
\end{equation}
The same $C,R_c$ work for every external derivative cutoff $J$.
Here $\varepsilon=2\exp(-7X\log k/10)$ includes the metric conversion
of the residual after a fixed enlargement of constants.

Choose a positive $r_0$ from the known coefficient and packet radii so
that all pressure absorptions hold, in particular $r_0R_c\le1$, and set
\[
 \eta=\exp(-\sqrt X),\qquad
 r(t)=r_0-2C(D_*/k+\eta)t.
\]
If
\begin{equation}\label{eq:packet-correction-guards}
 2\varepsilon e^{3CT}\le\eta/2,
 \qquad 2C(D_*/k+\eta)T\le r_0/2,
\end{equation}
then $r_0/2\le r(t)\le r_0$. In the bootstrap range
$\mathcal A_J\le\eta\le1$ the loss coefficient in
\eqref{eq:packet-correction-energy} is nonpositive. The scalar integral
inequality gives the strict bound $\mathcal A_J\le\eta/2$ and closes
the bootstrap.

For existence, we follow the pressure projection
and vanishing-viscosity argument of \cite[\S3.6]{OAI}. Let
$\mathscr G_\kappa$ be the closed gradient space above and let
$P_{\mathscr G}$ be its orthogonal projection in cylinder $L^2$.
On $\mathscr G_\kappa$ the map
$g\mapsto P_{\mathscr G}(G(t)g)$ is coercive, since
$\langle g,G(t)g\rangle\ge c\|g\|_2^2$ with $c>0$ supplied by $F$.
Set
\[
 \mathcal T_t=(P_{\mathscr G}G(t)|_{\mathscr G_\kappa})^{-1}
                         P_{\mathscr G},\qquad
 \mathcal P_t=I-G(t)\mathcal T_t.
\]
Then $P_{\mathscr G}\mathcal P_t=0$,
$\mathcal P_tu=u$ for $u\in\mathscr G_\kappa^\perp$, and
$\mathcal P_tG(t)g=0$ for $g\in\mathscr G_\kappa$.
Differentiating the coercive pressure equation gives the Sobolev bounds
for $\mathcal T_t$ and $\mathcal P_t$ at every fixed order. Their constants
may depend on that fixed order; the previously derived weighted energy
constants do not depend on the external derivative cutoff.

Write the unprojected correction source as
\[
 \mathcal N(E)=b_E\cdot\nabla_{y,\theta}E+\mathcal Z(E)+\mathcal E,
 \qquad b_E=\mathcal V_\kappa(Z+E),
\]
where $\mathcal Z(E)$ contains the terms $E\cdot\nabla_\kappa Z$,
$2F^{-1}F_tE$ and the differences of the two quadratic coefficient
terms. It has no derivative of $E$. For $0<\nu\le1$, solve
\begin{equation}\label{eq:packet-viscous-correction}
 E_t^\nu-\nu\Delta_{y,\theta}E^\nu
       =-\mathcal P_t\mathcal N(E^\nu),\qquad E^\nu(0)=0.
\end{equation}
The right side is locally Lipschitz $H^q\to H^{q-1}$ for
$q\ge b_d$, and the Fourier heat estimate
\[
 \|e^{\nu t\Delta}f\|_{H^q}
       \le C_q(1+(\nu t)^{-1/2})\|f\|_{H^{q-1}}
\]
gives a local mild solution by contraction. Every component of
$\nabla_\kappa$ has constant coefficients and commutes with the full
cylinder Laplacian. The heat semigroup therefore preserves the fixed
space $\mathscr G_\kappa^\perp$. Since the non-heat term in
\eqref{eq:packet-viscous-correction} takes values in this space, the mild
solution is exactly divergence free. No time derivative of
$\mathcal P_t$ and no commutation of $\mathcal P_t$ with the heat
semigroup is required.

The pressure gradient is
$-\mathcal T_t\mathcal N(E^\nu)$. The viscosity contributes no pressure,
because $\Delta E^\nu$ is still divergence free. Heat maximal regularity
justifies testing differentiated equations by
$F^{\top}F\partial^wE^\nu$: on each compact existence interval,
$E^\nu\in L^2_tH^{q+1}$ and $E_t^\nu\in L^2_tH^{q-1}$.
The differentiated metric-viscosity term satisfies
\[
 \nu\int\langle F^{\top}F u,\Delta u\rangle
       \le-\tfrac{c\nu}{2}\|\nabla u\|_2^2+C\nu\|u\|_2^2.
\]
This follows by integration by parts and Young's inequality, using the
known first derivatives of $F^{\top}F$. Thus it adds only a bounded
multiple of the lower energy to \eqref{eq:packet-correction-energy},
uniformly for $\nu\le1$. In the choice of $C$ above we include this
known metric-viscosity constant before selecting the radius or frequency;
it depends only on the fixed source and not on $J$ or $\nu$.
The same norm regularization justifies summing this estimate for
$\mathcal A_J$. The bootstrap extends every viscous solution
to $T$. Uniqueness in the low Sobolev class identifies the solutions
constructed at different higher orders. In particular they have uniform
bounds in $C_tH^q$ for each finite $q$.

The viscosity limit follows from a Cauchy estimate in the global
Sobolev spaces. For two viscosities, subtract their equations and use the
same metric $L^2$ difference energy. The transport and pressure leading
terms cancel as before. The remaining viscous defect is
$(\nu-\mu)\Delta E^\mu$, uniformly bounded in $L^2$ by a constant
of the fixed source times $|\nu-\mu|$. With zero initial difference,
Gronwall gives
\[
 \|E^\nu-E^\mu\|_{C_tL^2}\le C_*|\nu-\mu|.
\]
The dissipative term $\nu\Delta(E^\nu-E^\mu)$ is handled by the same
integration by parts. Interpolation with the uniform higher-order
bounds now gives a strong limit in $C_tH^q$ for every finite $q$.
The nonlinear source and the pressure inverse pass to this limit
one order lower; $\nu\Delta E^\nu\to0$ there. The gradient and
solenoidal constraints are closed, the finite weighted energies pass to
the limit, and the initial trace remains zero. The limiting equation
supplies the continuous strong time derivative at every Sobolev order,
including the within-interval endpoint derivatives.

These steps use the heat Fourier multiplier on
$\R^d\times\mathbb T_P$, a closed Hilbert gradient space, and the
Sobolev product threshold \eqref{eq:packet-base-index}. In the
three-dimensional proof the corresponding cylinder has dimension four;
here the only changes are the $d+1$ Fourier variables, all $d$ quadratic
coefficient matrices, and the dimension-dependent Sobolev constants.
The argument uses the full transverse space and tensor coefficients.

There is an explicit finite frequency choice. With
\[
 X=k^{1/203000},\qquad N=\lfloor X\rfloor,
\]
\eqref{eq:packet-tail-base} meets $B_N\le k^{1/100}$ at sufficiently
large $k$. Denoting by $I_*$ the inverse-multiplier
bound, the packet and correction smallness conditions hold if
\begin{equation}\label{eq:packet-explicit-threshold}
 k\ge\left(64+e+2^{202}C_*+I_*+12CT
       +\frac{8CTD_*}{r_0}
       +\left(\frac{8CT}{r_0}\right)^2\right)^{203000}.
\end{equation}
Indeed the $203000$th root controls $X$, the polynomial tail, $\log k$,
and the two terms in the radius loss; the negative exponential in
\eqref{eq:packet-small-residual} dominates $e^{3CT}$ and
$e^{-\sqrt X}$. The constants in \eqref{eq:packet-explicit-threshold}
depend only on the fixed source; no bound in terms of a subsequent
parent's parameters is assumed.

\subsection{The physical child and the signed pressure term}
Define the perturbed solution, called the child, by
\begin{equation}\label{eq:packet-exact-child}
 U^+(t,x)=U(t,x)+\kappa F(t,y)
          W(t,y,km_0\cdot y),\qquad y=\Phi_t^{-1}(x).
\end{equation}
The chain rule, $\Phi_t'=U\circ\Phi_t$, and
\eqref{eq:packet-frame} transform
\eqref{eq:packet-coordinate-equation} into the physical momentum equation
with pressure force
\[
 \nabla P_U(t,x)+\kappa F(t,y)^{-\top}
             \mathbf p(t,y,km_0\cdot y).
\]
This force is a gradient on physical space. Indeed, closed-gradient
membership implies
\[
 D_{\kappa,i}\mathbf p_j-D_{\kappa,j}\mathbf p_i=0
\]
in distributions: it holds for each compact smooth gradient generator
and is preserved by the $L^2$ limit. All spatial derivatives of the
constructed pressure vector are continuous, so the same identity holds
pointwise. The pullback chain rule for a one-form gives
\[
 \partial_{x_i}(\kappa F^{-\top}\mathbf p)_j
 -\partial_{x_j}(\kappa F^{-\top}\mathbf p)_i
 =\sum_{a,b}(F^{-1})_{bi}(F^{-1})_{aj}
       (D_{\kappa,b}\mathbf p_a-D_{\kappa,a}\mathbf p_b)=0
\]
after graph evaluation. The terms differentiating $F^{-1}=D\Phi_t^{-1}$
cancel by the symmetry of its ordinary second derivatives.

Graph evaluation uses the smooth periodic representatives and one extra
angular Sobolev trace estimate; the parent change of variables is volume
preserving. It gives all physical $L^2$ spatial derivatives and their
continuous time paths, both for the velocity and for this pressure force.
A curl-free $L^2$ vector field on $\R^d$ lies in the ordinary closed
gradient space, by its Fourier characterization. The Piola
identity also transforms the closed cylinder solenoidal constraint into
$\operatorname{div}U^+=0$. Taking the ordinary divergence of the physical
momentum equation and applying the Helmholtz decomposition identifies
its pressure force with the gradient of the Newton pressure in
Lemma~\ref{lem:newton}. We choose that normalized scalar pressure.
The smooth Sobolev local theory then identifies the constructed field
with the Euler evolution on the whole interval. The mean and finite
angular scalar pressures provide normalized
representatives, while the correction is constructed directly in the
closed gradient space.

The same field has the decomposition
\[
 U^+=U+V^{[N]}(t,y,km_0\cdot y)
        +\kappa F(t,y)E(t,y,km_0\cdot y).
\]
Thus its initial increment is exactly the finite packet, and the
correction initially vanishes. The correction estimates at all derivative orders give
physical velocity-gradient and pressure-Hessian errors bounded by fixed
current-source constants times $\eta$. One fast derivative of the velocity
correction is paired with its prefactor $\kappa$. The physical pressure
force of the correction is
$\kappa F^{-\top}\mathbf p^e$ evaluated on the graph. Its one
physical derivative, which is the pressure Hessian, has the same
cancellation of $\kappa$ against one fast derivative. For the finite
scalar primary pressure the corresponding factor is its prefactor
$\kappa^2$ against two fast derivatives. All remaining powers of
$\kappa$ are nonnegative. A fixed number of the uniform weighted
estimates controls the required coefficient and angular derivatives;
the stated constants are independent of $k$.

For both the joined and initial forward packets, set $P=2\pi$ and use
the explicit profile
\[
 f_\delta(\theta)=\arctan\frac{\sin\theta}{1+\delta-\cos\theta},
 \qquad 0<\delta\le1.
\]
Its first denominator is at least $\delta$, so the profile is smooth,
odd and $2\pi$-periodic. Oddness gives zero mean over $[-\pi,\pi]$,
and periodicity gives zero mean over any full period. Also
$|f_\delta|\le\pi/2$. Direct differentiation gives
\[
 f_\delta'(\theta)=\frac{N_\delta(\theta)}{D_\delta(\theta)},\qquad
 N_\delta=(1+\delta)\cos\theta-1,\qquad
 D_\delta=(1+\delta)^2-2(1+\delta)\cos\theta+1\ge\delta^2.
\]
At zero the quotient equals $\delta^{-1}$. The identities
\[
 D_\delta+N_\delta=(1+\delta)(1+\delta-\cos\theta)>0,
 \qquad
 D_\delta-\delta N_\delta
   =(1+\delta)(2+\delta)(1-\cos\theta)\ge0
\]
give the bounds used in the pressure estimate:
\[
 f_\delta'(0)=\delta^{-1},\qquad -1\le f_\delta'\le\delta^{-1}.
\]
The dependence of the Gevrey bounds on the width is explicit.
All derivatives of
$N_\delta$ have magnitude at most $3$, and the positive-order
derivatives of $D_\delta$ have magnitude at most $4$.
Differentiate $D_\delta D_\delta^{-1}=1$ and use the Gevrey binomial
convolution to obtain, by induction,
\[
 \|(D_\delta^{-1})^{(n)}\|_\infty
 \le 10\delta^{-2}(40\delta^{-2})^n(n!)^2.
\]
Multiplying by $N_\delta$ in the formula for $f_\delta'$ uses the same
radius; the convolution factor is at most $3$. Since $40\delta^{-2}\ge1$,
the resulting bound at order $n+1$, together with the bound
$|f_\delta|\le\pi/2$ at order zero, yields
\[
 \|f_\delta^{(n)}\|_\infty
 \le100\delta^{-2}(40\delta^{-2})^n(n!)^2\quad(n\ge0).
\]
Both the amplitude constant and radius are polynomial in $\delta^{-1}$.

Use this profile for the joined primary \eqref{eq:packet-joined-primary}.
The uncut physical high velocity is $b_Y$, and the amplitude is $\alpha$.
At
$y=\Phi_t^{-1}(x)$ and $\vartheta=km_0\cdot y$, the leading velocity
and pressure Hessian tensors are respectively
\begin{align}
 \mathcal S&=\alpha\chi_{\rm in}(y)f_\delta'(\vartheta)\,
                   b_Y\otimes m,\notag\\
\mathcal Q&=-2\alpha\chi_{\rm in}(y)
       \frac{\langle m,Mb_Y\rangle}{|m|^2}
       f_\delta'(\vartheta)\,m\otimes m.
                                                        \label{eq:packet-leading-tensors}
\end{align}
Direct differentiation of the primary velocity gives
\[
 D_x\{k^{-1}\alpha\chi_{\rm in}(y)f_\delta(\vartheta)b_Y(t,y)\}
 =\mathcal S+k^{-1}\alpha f_\delta(\vartheta)
                    D_y(\chi_{\rm in}b_Y)(t,y)F^{-1}(t,y).
\]
The slow coefficient therefore contributes a constant times $k^{-1}$.
The corresponding pressure-force coefficient is
$-2\alpha\chi_{\rm in}\langle m,Mb_Y\rangle m/|m|^2$;
its derivative is the other slow primary remainder. To bound these
derivatives, it suffices to bound the carrier and its first
label derivative on the ball of radius $1/2$. Both coefficients
vanish off the ball of radius $1/4$, so their derivatives vanish
locally outside that support. The history inverse and the supported
tail estimates above supply these carrier bounds. This gives uniform
spatial remainder constants without any sharp shear estimate away
from the central particle.
The finite higher grades, derivatives of slow coefficients, and
exact correction give
\begin{equation}\label{eq:packet-low-errors}
 \|DU^+-DU-\mathcal S\|\le C_v/k+C'_v\eta,
 \quad
 \|D^2P_{U^+}-D^2P_U-\mathcal Q\|\le C_p/k+C'_p\eta,
\end{equation}
uniformly on the closed time interval and physical space. The pressure
identity follows first for the finite scalar pressure and then for
the correction pressure; equality of their gradient with the
canonical Euler pressure gradient transfers the identity to
$P_{U^+}$.

If $\langle m,Mb_Y\rangle\ge0$, the lower slope bound gives the one-sided
estimate
\begin{equation}\label{eq:packet-signed-curvature}
 \langle\mathcal Qz,z\rangle
 \le2\alpha\chi_{\rm in}\langle m,Mb_Y\rangle |z|^2.
\end{equation}
Without this sign condition, the absolute estimate is
\[
 \|\mathcal Q\|\le
   2\alpha\|M\|\chi_{\rm in}|m|\,|b_Y|/\delta.
\]
The one-sided estimate retains the gain from the narrow profile in the
mean curvature bound at the next stage.
All tensors in \eqref{eq:packet-leading-tensors} act on the entire
$\R^d$, and the norm identities used here hold in every dimension.

For spatial localization the preceding construction is applied to the
rescaled parent
\[
 U^\rho(t,x)=\rho^{-1}U(t,\rho x),\qquad
 P_{U^\rho}(t,x)=\rho^{-2}P_U(t,\rho x),\qquad 0<\rho\le1.
\]
Direct substitution proves the Euler equations on the same time interval.
Its flow is $\rho^{-1}\Phi_t(\rho y)$, so its frame, strain and
pressure Hessian in labels are exactly $F(t,\rho y)$,
$M(t,\rho y)$ and $H(t,\rho y)$. Undoing the scaling after constructing
the child gives the physical field
$\rho(U^\rho)^+(t,x/\rho)$. Velocity gradients and pressure Hessians
are unchanged in size by these two inverse scalings, and the support
radius of a compact oscillatory coefficient is multiplied by $\rho$.
Thus the low error bounds and leading tensors above transfer to the
original parent with no factor $\rho^{-1}$ in those two derivative
orders. Higher-order initial derivatives retain their scaling
factors as recorded in Section~\ref{sec:initial-summability}.

\subsection{Physical localization of the mean boundary condition}
The oscillatory localization scale and the mean cutoff scale are chosen
separately. With a fixed physical cutoff parameter $\ell_*$ and physical
unscaling $x=\rho y$, use the normalized cutoff parameter
$\ell=\rho\ell_*$. Then
$\chi(\ell y)=\chi(\ell_*x)$, and
\eqref{eq:packet-mean-initial} places each physical mean initial
increment in $\overline B(0,\ell_*^{-1})$. If
$\operatorname{supp}\chi_{\rm in}\subset\overline B(0,1/4)$, the high
and corrector initial increments are supported in
$\overline B(0,\rho/4)$. The exact correction has initial value zero.
Lemma~\ref{lem:fixed-cutoff} proves propagation of the mean reserve
with this fixed physical cutoff, and Lemma~\ref{lem:common-support}
gives the resulting common compact support through the iteration.


\subsection{The packet construction for one parent}
The following lemma summarizes the construction for use in the iteration.
Its assumptions concern only the fixed source.

\begin{lemma}[Exact joined packet child]\label{lem:exact-packet-child}
Fix $d\ge3$ and a smooth Sobolev parent on $[0,T]$ with flow
\eqref{eq:packet-frame}, the known coefficient bounds
\eqref{eq:packet-known-jets}, and the mean reserve
\eqref{eq:packet-mean-guard}. Fix its transverse terminal datum, compact
Gevrey cutoffs, and the positive forward profile with
\eqref{eq:packet-tail-propagator}. There is a finite constant determined
by these data such that every sufficiently large $k$ gives the child
\eqref{eq:packet-exact-child} from the joined strict-prefix family.
It solves the full $d$-dimensional Euler equation with its Newton pressure
on the whole closed interval and belongs to
$C_tH^m\cap C_t^1H^{m-1}$ for every integer $m\ge1$.

Its leading velocity-gradient and pressure-Hessian increments are
exactly the two tensors in \eqref{eq:packet-leading-tensors}. After an
explicit enlargement of the known-source frequency threshold, each
error in \eqref{eq:packet-low-errors} is at most $k^{-1/4}$, uniformly in
time and space. Its initial increment is the entire finite sum
\eqref{eq:packet-finite-pairs} at time zero, including every higher-order
high-frequency initial velocity and the terminal corrector $C_N$; the exact correction
has initial value zero.

If the parent is odd and the cutoffs are even, then the child is odd.
If in addition a fixed orthogonal coordinate reflection $\mathscr R$
obeys
\[
 U(t,\mathscr Rx)=\mathscr R U(t,x),\quad
 \mathscr R m_0=m_0,\quad \mathscr R R_\perp Y=R_\perp Y,
\]
and preserves the cutoffs, then
\[
 U^+(t,\mathscr Rx)=\mathscr R U^+(t,x),\qquad
 DU^+(t,\mathscr Rx)=\mathscr R DU^+(t,x)\mathscr R^{-1}.
\]
For the fixed physical mean cutoff described above, every
initial increment is supported in
$\overline B(0,\max(1/4,\ell_*^{-1}))$. Thus the child's initial support
is contained in the union of this ball and the parent's initial support.
\end{lemma}
\begin{proof}
The preceding constructions give existence, the Euler equation,
the normalized pressure, and all Sobolev orders. The residual identity
includes the terminal corrector and pressure.

To obtain the stated error bound, let
$\beta=203000$ and let $C_{\rm err}\ge1$ dominate the four constants in
\eqref{eq:packet-low-errors}. Increase the lower bound for $X=k^{1/\beta}$
to include $(2\beta)^4$ and $(4C_{\rm err})^2$, as well as the quantity
inside parentheses in \eqref{eq:packet-explicit-threshold}.
For $X\ge1$, $\log X\le4X^{1/4}$, and hence
\[
 \tfrac14\log k\le\tfrac12\sqrt X,
 \qquad \log(4C_{\rm err})\le\tfrac12\sqrt X.
\]
It follows that
$C_{\rm err}e^{-\sqrt X}\le k^{-1/4}/4$ and
$C_{\rm err}/k\le k^{-1/4}/4$. The sum satisfies the asserted error
bound. The enlarged threshold remains a fixed power of a sum of known
source constants.

For symmetry, uniqueness of the parent's flow first gives
$\Phi_t(-y)=-\Phi_t(y)$ and the corresponding reflection covariance.
The mean form commutes with the induced orthogonal action, since its
cutoff tensor operator is constructed using all tensor indices and the
cutoff is invariant. Uniqueness of its coercive inverse gives the mean
parities. For the joined high operator the action preserves the fixed
zero-trace displacement space; for the primary it also preserves the
terminal datum. Uniqueness of the history inverse and then of the
forward initial-value equation proves the high parities. The normalized
angular primitive, the full tensor corrector, and each product
in the strict-prefix force commute with the same actions. Induction
therefore gives odd vector fields and even scalar pressures for the
whole finite family, together with the stated reflection covariance.

The projected coordinate equation and its closed solenoidal and gradient
spaces are also invariant under these actions. Its unique correction
with zero initial value has the same parities. Transforming back gives
the child symmetry; differentiating this identity gives the
claimed full derivative covariance. Finally, the initial support
follows from \eqref{eq:packet-mean-initial}, support preservation of the
joined high inverse and the tensor potential, and $E(0)=0$.
\end{proof}
\section{A fixed compact support and the full mean reserve}
The initial mean boundary cutoff and the oscillatory packet have
different spatial scales. In
physical coordinates fix two parameters $0<\ell_*\leq1$ and
$0\leq r_*\leq1/4$. Write $c_d=C_{2,d}$ for the nonnegative
volume-error constant in the localized mean boundary inequality, and
$C_{1,d}$ for its boundary-operator constant. These dimensional constants
are denoted separately from the stage-dependent bound on the second
frame coefficient.
Let $s_{\rm cut,d}>0$ be the fixed dimensional cutoff scale and put
$a_d=(3s_{\rm cut,d})^{-1}$. For the current
solution horizon $L>0$, define the mean reserve by
\begin{equation}\label{eq:full-mean-reserve}
 \mathfrak M(L)=\tfrac12 K L^2+(B_e+c_d B_c r_*^d)L.
\end{equation}
The parameters $K,B_c,B_e$ are nonnegative. Here $K$ bounds the pressure
Hessian above, while $B_c$ and $B_e$ bound
the negative part of the initial symmetric gradient inside and outside
the fixed ball of radius $a_d r_*/\ell_*$, respectively. The boundary
penalty $\mathcal L$ is chosen to be at least $C_{1,d}B_c$.

\begin{lemma}[Fixed-cutoff propagation]\label{lem:fixed-cutoff}
For a constructed child let $e_p\geq0$ be its upper pressure-Hessian
increment and $e_i\geq0$ its uniform initial-gradient increment. Retain
the same physical $\ell_*,r_*$, replace
\[
 K\mapsto K+e_p,\quad B_e\mapsto B_e+e_i,\quad
 B_c\mapsto B_c+e_i,\quad \mathcal L\mapsto C_{1,d}(B_c+e_i).
\]
All corresponding mean inequalities hold whenever the resulting
$\mathfrak M(L)$ is less than $1/2$, and
\begin{equation}\label{eq:mean-reserve-increment}
 \mathfrak M_{\rm child}(L)-\mathfrak M_{\rm parent}(L)
   =\tfrac12 e_p L^2+(1+c_d r_*^d)e_iL.
\end{equation}
\end{lemma}
\begin{proof}
The pressure upper bound adds $e_p$. If the difference of the initial
gradients has operator norm at most $e_i$, its quadratic form is at
least $-e_i|z|^2$. The old interior and exterior inequalities therefore
persist on exactly the same regions with the indicated new constants.
The new boundary penalty meets the same coercivity requirement.
Expanding \eqref{eq:full-mean-reserve} proves
\eqref{eq:mean-reserve-increment}. On any shorter horizon the reserve
only decreases. The updated parameters are used for the subsequent mean
problem; the child's flow, inverse, time derivatives, and pressure remain
those obtained from the packet construction.
\end{proof}

The comparisons in this lemma are made at the parent's horizon $L$.
Passing to a shorter horizon decreases the reserve. Each nonnegative
increment may be estimated using a common upper horizon, while the
preceding reserve is evaluated on the current solution's horizon.

The spatial normalization preserves this full functional exactly.
At physical localization length $\rho$, the normalized mean scale is
$\rho\ell_*$ and its splitting parameter is still $r_*$. Unscaling
changes the scale back to $\ell_*$ and leaves $K,B_c,B_e,r_*$ unchanged.
Since $\ell_*\leq1$, the final cutoff normalization makes no further
change. Thus the physical interior radius remains $a_dr_*/\ell_*$,
and both the normalized and physical mean reserves equal
\eqref{eq:full-mean-reserve}. Changing auxiliary source bounds does
not alter these mean parameters.

The two initialization packets permit distinct initial increments $e_c,e_e$
on the interior and exterior regions. The same proof then gives
\[
 \mathfrak M_{\rm child}(L)-\mathfrak M_{\rm parent}(L)
      =\tfrac12 e_pL^2+(e_e+c_dr_*^de_c)L.
\]
For the seed the fixed physical choice is
$\ell_*=1$, $r_*=X^{-1000}$. Its source and frequency bounds
make the full child reserve at most $1/8$ by increasing $X$ after
$J$ is fixed. The first normal forward packet contributes at most $1/16$
to the pressure-time term and $1/16$ to the initial-time term.
Consequently the first positive-history state has mean reserve at most
$1/4$. These estimates include both the interior contribution and
the nonzero exterior mean fields with the same fixed cutoff.

For subsequent stages, the one-sided pressure estimate is as follows.
Let $w$ be the transported primary carrier, including its selected
amplitude $\alpha$, let $m$ be the transported normal, and let
$M$ be the parent gradient in particle labels. On the packet support,
$\|M\|\leq M_0$. Write $b=\chi|m||w|$, with the nonnegative
spatial cutoff $\chi$. The principal pressure Hessian is
\[
 Q=-2\chi\frac{\langle m,Mw\rangle}{|m|^2}
                   f'_\delta\,m\otimes m.
\]
For the chosen periodic profile $f'_\delta\geq-1$ and
$|f'_\delta|\leq\delta^{-1}$ when $0<\delta\leq1$.
On the good-time set the flux is nonnegative; hence
$Q\leq2\chi\langle m,Mw\rangle I\leq2M_0bI$ as quadratic forms.
The supported localization estimates give $b\leq\delta G$ there.
On the complementary times they give $b\leq\delta E$, and the
absolute estimate yields $\|Q\|\leq2M_0E$. Outside the cutoff $Q=0$.
The child's canonical pressure Hessian therefore has the global upper
increment
\[
 e_p\leq2M_0(\delta G+E)
       +\frac{C_p(\mathcal W)}k
       +C_{\rm corr}(\mathcal W)\delta_{\rm corr}(X_{\rm exp}).
\]
The remainder polynomials follow from the carrier and
coefficient derivative estimates, including the correction that eliminates
the residual. Every carrier norm is controlled by these estimates.
The supported bounds on the selected profile give
$G\leq C_d h_{\rm child}$ and
$E\leq C_d(\mathcal P)h_{\rm child}e^{-x_n/8}$, where
$\mathcal P$ is the source scale determined by the parent. The
coefficient of $\delta h_{\rm child}$ is controlled by this scale,
which excludes the angular factor $\delta^{-1}$.

The initial-gradient increment satisfies
$e_i\leq C_i(\mathcal W)(\alpha+k^{-2})$.
Retain the factor $1+c_dr_*^d$ in
\eqref{eq:mean-reserve-increment}, and bound it above in the new
error terms by the fixed constant $B_{\rm weight}=1+C_{2,d}$.
There are five error contributions: the principal pressure
terms on the good-time set and its complement, the initial term
proportional to $\alpha$, the pressure remainder, and the mean initial
term proportional to $k^{-2}$. The prescribed
expansion satisfies $\delta_{\rm corr}(X_{\rm exp})\leq k^{-1}$, so
both pressure remainders are included in the same frequency-dependent
contribution. The two initial contributions include $B_{\rm weight}$.
The preceding reserve retains the factor $r_*^d$.

The coefficient and damping estimates in
Section~\ref{sec:parameter-envelope} and Section~\ref{sec:iteration}
control these five contributions at the prescribed frequencies. In particular,
$\delta h_{\rm child}\leq e^{-x_n/(2(J+n)^3)}$ and
$h_{\rm child}e^{-x_n/8}\leq e^{-x_n/16}$, while the selected
$\alpha$ and the frequency remainders have their own packet-source
bounds. The finite parameter selection bounds the sum of each of these
five contributions by $1/64$.
At every finite positive-history stage the mean reserve is bounded
by the partial sum
\[
 \frac14+\sum_{i<n}\sum_{a=1}^5\operatorname{cost}_{a,i}
 \leq\frac14+\frac5{64}=\frac{21}{64}<\frac12.
\]
Including the nonnegative index-zero contribution in this partial sum
preserves the upper bound, since the two initialization reserves have
already been estimated separately. The dimensional multiplier
and the seven additional initialization contributions belong to the finite
list fixed before choosing $J$. The remaining lower bounds are met
by enlarging $X$ before constructing any packet. This changes the
frequency values through their fixed formulas, while all construction
exponents, error contributions, and the chosen $J$ remain fixed.

\begin{lemma}[One compact set for all initial increments]\label{lem:common-support}
With this fixed-cutoff construction every initial increment is supported
in $\overline B(0,\max(1/4,\ell_*^{-1}))$, irrespective of its
frequency, localization length, or finite truncation order.
\end{lemma}
\begin{proof}
At a stage with physical localization length $0<\rho\leq1$, normalized
labels are $y=x/\rho$, and the mean cutoff is
$\chi((\rho\ell_*)y)$. The mean boundary condition expresses
its initial velocity as the compactly supported solenoidal boundary
operator with this cutoff, multiplied by its scalar penalty.
Hence the normalized mean initial field vanishes for
$|y|>(\rho\ell_*)^{-1}$. Physical unscaling replaces $y$ by $x/\rho$,
so it vanishes for $|x|>\ell_*^{-1}$. This holds for every
mean grade. The oscillatory profiles and their full tensor correctors
are supported in $|y|\leq1/4$, hence after unscaling in
$|x|\leq\rho/4\leq1/4$. The exact correction has zero initial value.
Summing the finite family, including its
terminal corrector, proves the assertion. Adding the fixed compact
support of the base datum gives one compact set for all finite-stage
initial velocities and their pointwise limit.
\end{proof}
\section{Joined source estimates and a uniform parameter envelope}
\label{sec:parameter-envelope}

We fix the dimension throughout this section. The Sobolev base order
$q=2d+5$ is fixed before any stage or truncation is chosen. Increasing this
order changes constants and polynomial degrees only.

\subsection{The variational history inverse}

For every forced high grade, the inverse consists of a Dirichlet problem
before activation and a forward evolution after activation. The Dirichlet
coercivity estimate gives polynomial dependence on the parent coefficients
on the history interval. On the forward interval this dependence follows
from the relative propagator estimate. Compact-interval existence of a
forward fundamental matrix alone does not provide either quantitative bound.

The quantitative estimate is established in the transverse space
$U=m_0^\perp$, of dimension $d-1$. Let
\[
 Q(t,x)=F(t,x)|_U,\qquad Q_t=Q_1,\qquad
 Q_{tt}=-\mathcal H Q,\qquad
 \|Qv\|^2\geq c\|v\|^2,\quad c>0.
\]
The symmetric matrix $\mathcal H$ is the pressure Hessian along
the preceding particle map. For a positive activation time $0<t_0\leq T$, suppose on $[0,t_0]$
\[
 \langle\mathcal H v,v\rangle\leq K\|v\|^2,\qquad
 0\leq K,\qquad Kt_0^2/2\leq1/2.
\]
These are the inherited global one-sided pressure bounds. The absolute
coefficient norms below are obtained from the particle-label fields:
$F_1=D_x(\partial_t\Phi)$, $F_2=D_x(\partial_t^2\Phi)$,
$\mathcal H=-F_2F^{-1}$, and $\det F=1$ expresses $F^{-1}$ by cofactors.
The fixed Sobolev shift and factorial product bounds therefore make all
these norms polynomial in the label amplitude and radius. The sharper
central bounds of the preceding section are used for central
endpoint selection and renewal. All norms below can also be taken in
spatial-angular $L^2$; the multiplication coefficients have uniform
pointwise bounds.

For a forcing $f$, solve for $z\in H^1_0((0,t_0);U)$:
\begin{equation}\label{eq:joined-history-form}
 \int_0^{t_0}\!
 \big\langle(Qz)',(Q\zeta)'\big\rangle
 -\big\langle\mathcal H Qz,Q\zeta\big\rangle\,dt
 =-\int_0^{t_0}\langle f,Q\zeta\rangle\,dt
 \quad(\zeta\in H^1_0((0,t_0);U)).
\end{equation}
Its physical Jacobi displacement is $y=Qz$; its physical high velocity is
\begin{equation}\label{eq:joined-history-high}
 b=Qz',
\end{equation}
which is generally different from $y'$. In particular, $z(0)=z(t_0)=0$
does not assert that $b(0)$ or $b(t_0)$ vanishes.

\begin{lemma}[Polynomial history inverse]\label{lem:joined-history-polynomial}
The solution of \eqref{eq:joined-history-form}, its high velocity,
its time derivative and its terminal high velocity have bounds
polynomial in
\[
 1+t_0+t_0^{-1}+c^{-1}+
 \|Q\|_\infty+\|Q_1\|_\infty+\|\mathcal H\|_\infty .
\]
Their fixed-order spatial-angular Sobolev bounds are also polynomial in
the corresponding fixed-order coefficient bounds. The assertion holds
in every finite transverse dimension.
\end{lemma}
\begin{proof}
Put $C_0=\|Q\|_\infty$, $C_1=\|Q_1\|_\infty$ and
$L=(Q^*Q)^{-1}Q^*$. Then $LQ=I$ and
\[
 \|L\|\leq c^{-1}C_0,\qquad
 \|L'\|\leq c^{-1}C_1+2c^{-2}C_0^2C_1.
\]
Consequently, with
\[
 A=1+c^{-1}C_0+
       t_0(c^{-1}C_1+2c^{-2}C_0^2C_1),
\]
the identity $z=Ly$ and the elementary inequality
$\|y\|_{L^2}\leq t_0\|y'\|_{L^2}$ give
$\|z'\|_{L^2}\leq A\|y'\|_{L^2}$.
The sharper primitive bound
$\|y\|_{L^2}^2\leq(t_0^2/2)\|y'\|_{L^2}^2$ gives
\[
 \int_0^{t_0}(|y'|^2-\langle\mathcal Hy,y\rangle)
 \geq\tfrac12\|y'\|_{L^2}^2
 \geq(2A^2)^{-1}\|z'\|_{L^2}^2.
\]
Thus the coercive variational theorem on the fixed Hilbert space $H^1_0$
constructs the inverse. Testing with $z$ also gives
\[
 \|y'\|_{L^2}\leq2t_0\|f\|_{L^2},\qquad
 \|z'\|_{L^2}\leq2At_0\|f\|_{L^2}.
\]
Both bounds are polynomial in $C_0$, $C_1$ and $\|\mathcal H\|$.

Integration by parts and $Q''=-\mathcal HQ$ give
\[
 Q^*(Qz''+2Q_1z')=Q^*f.
\]
Therefore the acceleration is
\begin{equation}\label{eq:joined-gram-acceleration}
 z''=(Q^*Q)^{-1}\{Q^*f-2Q^*Q_1z'\},
\end{equation}
and
\[
 \|z''\|_{L^2}\leq
 c^{-1}C_0(\|f\|_{L^2}+2C_1\|z'\|_{L^2}).
\]
The time trace inequality
\[
 \|z'\|_{C_t}\leq
 t_0^{-1/2}\|z'\|_{L^2}+t_0^{1/2}\|z''\|_{L^2}
\]
bounds both endpoints polynomially. For continuous $f$,
\eqref{eq:joined-gram-acceleration} gives continuous $z''$, and
$b'=Q_1z'+Qz''$ is the derivative on the closed interval.
In particular the forward initial datum $b(t_0)$ is a bounded trace
of the solution of \eqref{eq:joined-history-form}.

Spatial derivatives are obtained by differentiating the operator on
the fixed space $H^1_0((0,t_0);U)$. Every commutator contains only
coefficient derivatives and a lower derivative of the same solution.
Applying the same coercive inverse proves the fixed-order estimates
by finite induction. For example, a sufficient bound for the inverse is
\[
 I_0=I,\qquad I_{r+1}=I+I_r+2^rB I_r^2,
\]
where $I$ bounds the base inverse and $B$ the coefficient block.
This is a polynomial after the fixed number $q$ of steps.
The Gram inverse in \eqref{eq:joined-gram-acceleration} has the same
positive-order commutator proof. These arguments use adjoints, norms
and Hilbert-space coercivity in arbitrary transverse dimension.
\end{proof}

On $[t_0,T]$ solve the forward transverse equation with initial
datum $b(t_0)$ obtained above. Its value and derivative match those of
the history, by the same equation and forcing at the junction.
The algebraic normal pressure is reconstructed from this joined field.
The profile is $\gamma(t)=\alpha$ on the history interval and
$\gamma(t)=\alpha g(t)$ afterwards, where $g(t_0)=1$ and $\alpha>0$.
Normalization by $\gamma$ in spatial word estimates never differentiates
$\gamma$ in time.

\subsection{The full transverse relative estimate}
\label{subsec:full-transverse-relative}

We derive the estimate for the full transverse space. At activation
let $p,q$ be the preceding unit normal and unit velocity, and put
\[
 a=\langle Bq,p\rangle>0,\qquad
 \beta=\frac{\|\Pi_{\{p,q\}^{\perp}}Bq\|}{a},\qquad
 \sigma=\sqrt\beta,\qquad \varepsilon=\sqrt{a/h}.
\]
Here $h$ is the preceding shear height. Assume $0<\sigma\leq1/4$,
$1/2\leq a\leq2$ and $0<\varepsilon\leq1$.
Choose the third unit vector $n$ in the direction of
$\Pi_{\{p,q\}^{\perp}}Bq$; the positive tilt makes it nonzero.
Write $E=\operatorname{span}\{p,q,n\}$ and
$\mathcal E=E^\perp$. The center reflection fixes $E$ and reverses
$\mathcal E$. The actual center coefficient therefore preserves this
splitting. Neighboring coefficients need not preserve it.

Use the inherited bounds $G\geq1$, $\|B\|\leq G$ and
$\|B_t\|\leq G^2$. Transport the entire orthonormal frame by its skew
connection $S$. In this frame, for the preceding background matrix $B$,
\[
 S=p\otimes B^*p-B^*p\otimes p+
       q\otimes\Pi_{\{p,q\}^{\perp}}Bq
       -\Pi_{\{p,q\}^{\perp}}Bq\otimes q .
\]
Consequently $\|S\|\leq4G$ if $\|B\|\leq G$, and the exact column identity is
\begin{equation}\label{eq:full-connection-column}
 (B+S)q=2\langle Bq,p\rangle p+\langle Bq,q\rangle q.
\end{equation}
The normal-normal connection is zero. Differentiating the transported
background gives norm at most $G^2+2(4G)G=9G^2$.
At activation its column outside the $q$ diagonal is
$a(p+\beta n)$. At later physical time $t$ its defect is therefore at
most $9G^2|t-t_0|$. These statements concern the complete vector column,
including all additional components.

Use fast time $\tau=a(t-t_0)/\varepsilon$. The estimate is restricted to
$0\leq\tau\leq\tau_{\rm end}$; assume
$\tau_{\rm end}\leq\Theta$, where $\Theta\geq1$ is the numerical horizon.
The estimate does not require extending the coefficients.
At a neighboring label the actual gradient in this same frame is
\[
 M=B+h(t)\,q\otimes p+\mathcal R,\qquad
 \|\mathcal R\|\leq e_{\rm old}+\rho L_M .
\]
The rank-one term uses the center pair; the neighboring discrepancy is
included in $\mathcal R$.

Define the diagonal maps
\[
 D_r=\operatorname{diag}(1,\varepsilon,1,I_{\mathcal E}),\qquad
 D_v=\operatorname{diag}(\varepsilon,1,\varepsilon,
                                      \varepsilon I_{\mathcal E}).
\]
They satisfy $D_rD_v=\varepsilon I$ and both have norm at most one.
The actual ray and high velocity, including every extra coordinate, are
\begin{equation}\label{eq:ambient-scaled-states}
 m=s_0 O(t)D_r y,\quad y=(P,Q,N,r),\qquad
 v=O(t)D_v z,\quad z=(U,V,W,Z),\quad r,Z\in\mathcal E.
\end{equation}
The primary solution uses one fixed coordinate endpoint at every label.
More precisely, if the endpoint lemma selects the physical vector
$Y_{\rm phys}\in m(t_0,0)^\perp$, set
\[
 Y_{\rm coord}=R_\perp^*F(t_0,0)^{-1}Y_{\rm phys}.
\]
At the center, $F(t_0,0)R_\perp Y_{\rm coord}=Y_{\rm phys}$.
At every other label the history has the same coordinate endpoint
$Y_{\rm coord}$. Its norm is at most
$\|F(t_0,0)^{-1}\|\|Y_{\rm phys}\|$. Thus $Y_*$ in the history
Lipschitz estimates below includes this inverse-frame factor.

For $t_0>0$, the unlocalized history is solved pointwise in the finite-dimensional
transverse space and has uniform spatial derivative bounds. Its
cylinder $L^2$ realization is obtained after multiplication by the
compact seed. To verify this directly, put
$Z(y,\theta)=\chi_{\rm in}(y)f_\delta(\theta)Y_{\rm coord}$
and $z_{\rm lin}(t,y,\theta)=tZ(y,\theta)/t_0$.
The zero-trace remainder $w=z-z_{\rm lin}$ solves
\eqref{eq:joined-history-form} with forcing
\[
 f_{\rm lin}=-(2/t_0)Q_1Z.
\]
This identity follows from $Q_{tt}=-\mathcal HQ$ and
$z_{\rm lin,tt}=0$. The forcing is a continuous path in every
cylinder Sobolev space, is angular-mean zero, and has compact spatial
support. Thus the Hilbert-space inverse applies. Its uniqueness
identifies the result with the localized pointwise history. All
cylinder estimates concern this localized field. For $t_0=0$ the
prescribed forward primary is used instead of the history construction.

Here $s_0>0$ and $O(t)$ is the transported orthogonal frame.
The velocity $v$ is the unamplified selected high solution. The packet
uses $\alpha v$, with $\alpha>0$ determined by its target size below.
Set
\[
 A=a^{-1}D_rMD_v,\qquad C=a^{-1}D_r(M+S)D_v .
\]
Direct substitution into the ray and high equations gives
\begin{equation}\label{eq:ambient-exact-generator}
 y'=(C^*-2A^*)y,\qquad
 z'=\mathscr B_\varepsilon(y)z,\qquad
 \mathscr B_\varepsilon(y)
  =-C+2\frac{D_r^2y}{|D_ry|^2}\otimes A^*y .
\end{equation}
For example, the pressure numerator transforms by
$\langle D_ry,MD_vz\rangle=a\langle y,Az\rangle$.
Thus \eqref{eq:ambient-exact-generator} is the rescaled equation
for all transverse components.

Let
\[
 A_0=(p+\beta n)\otimes q+q\otimes p,\qquad
 C_0=2p\otimes q+q\otimes p .
\]
The large shear term scales to
$(\varepsilon^2h(t)/a)q\otimes p$.
For any matrix $L$, separating its $q$ input column gives
\[
 a^{-1}D_rLD_v
 =(\varepsilon/a)D_rL+
          ((1-\varepsilon)/a)(D_rLq)\otimes q.
\]
Together with \eqref{eq:full-connection-column}, this proves
\begin{equation}\label{eq:ambient-matrix-gap}
 \max(\|A-A_0\|,\|C-C_0\|)
 \leq \left|\frac{\varepsilon^2h(t)}a-1\right|
 +\frac{11\varepsilon G+18G^2|t-t_0|+
                    e_{\rm old}+\rho L_M}{a}.
\end{equation}
One obtains the slightly better coefficient 3 in place of 11 for $A$;
the displayed common bound suffices. The norm of an additional input
column has the small factor $\varepsilon$: the shear is zero on
$\{p,q\}^\perp$. Both diagonal maps have norm at most one, so matrix
scaling introduces no factor $\varepsilon^{-1}$.

The logarithmic derivative of $h(t)$ is bounded by $2G$.
For $G\varepsilon\Theta/a$ small, its contribution to
\eqref{eq:ambient-matrix-gap} is $O(G\varepsilon\Theta/a)$.
Choose a single error parameter $e\geq\varepsilon$ large enough that
the right side of \eqref{eq:ambient-matrix-gap} is at most $2e$,
and that the initial ray error is at most $e$.
It is bounded by a fixed constant times
\begin{equation}\label{eq:ambient-error-input}
 \varepsilon\Theta G^2+
 e_{\rm old}+
 \rho\left(L_M+\frac{3L_m}{s_0\varepsilon}
          +\frac{2L_{\rm hist}Y_*}{\varepsilon}\right).
\end{equation}
The lower bound on $a$ was used here. The constant $L_m$ bounds the
spatial Lipschitz norm of the normal at activation. The constant
$L_{\rm hist}$ bounds, in operator norm, the spatial derivative of the
linear map $Y_{\rm coord}\mapsto b_{Y_{\rm coord}}(t_0,y)$.
Finally $Y_*$ bounds the norm of the selected coordinate endpoint.
Consequently the spatial variation of its high velocity is at most
$\rho L_{\rm hist}Y_*$.

The reference ray is
\[
 y_0(\tau)=(\beta\tau^2,-2\beta\tau,1,0).
\]
Its full generator $C_0^*-2A_0^*=-p\otimes q-2\beta q\otimes n$
is nilpotent of order three. Its evolution is exactly
\[
 I-(\tau-s)p\otimes q-2\beta(\tau-s)q\otimes n+
                  \beta(\tau-s)^2p\otimes n
\]
and has norm at most $5\Theta^2$.
Duhamel's formula and absorption in the supremum norm give, if
$60e\Theta^3\leq1$,
\begin{equation}\label{eq:ambient-ray-error}
 \|y-y_0\|\leq310e\Theta^5.
\end{equation}
Indeed the difference generator has norm at most $6e$; the resulting
bound is
$10\Theta^2(e+30e\Theta^3)\leq310e\Theta^5$.
Under the stronger smallness condition used below, this is at most
$1/2$. Thus
\[
 N\geq\tfrac12,\qquad |D_ry|^2\geq\tfrac14,\qquad
 \|y\|\leq6\Theta^2.
\]
In particular, \eqref{eq:ambient-ray-error} controls the extra ray components $r$.

Tangency is $y\cdot z=0$, so the eliminated coordinate is exactly
\begin{equation}\label{eq:extra-tangent-elimination}
 W=-\frac{PU+QV+r\cdot Z}{N}.
\end{equation}
Let $\zeta=(U,V,Z)$, let $\Pi$ omit the $n$ coordinate, and let
$E_y\zeta$ insert the value \eqref{eq:extra-tangent-elimination}.
Then $\|E_y\|\leq1+2\|y\|\leq13\Theta^2$, and the actual reduced
equation on the full $d-1$ dimensional transverse space is
$\zeta'=\Pi\mathscr B_\varepsilon(y)E_y\zeta$.
It retains all couplings to $Z$.

The pressure denominator admits the following operator-norm estimate on
the full ambient space. If
$r_*=\|y-y_0\|\leq1/2$, $R_*\geq\max(\|y\|,\|y_0\|)$ and
$\eta=\max(\|A-A_0\|,\|C-C_0\|)\leq1$, then
\begin{align}
 \|\mathscr B_\varepsilon(y)-\mathscr B_0(y_0)\|
 &\leq \eta+
 8R_*(4+8R_*^2)(r_*+\varepsilon R_*)\nonumber\\
 &\hspace{18mm}+2R_*(\eta R_*+3r_*).              \label{eq:ambient-generator-gap}
\end{align}
To prove it, put $u_\varepsilon=D_r^2y/|D_ry|^2$.
The preceding lower denominator bounds and
$\|D_ry-D_0y_0\|\leq r_*+\varepsilon R_*$
give
$\|u_\varepsilon-u_0\|\leq
(4+8R_*^2)(r_*+\varepsilon R_*)$.
Also
$\|A^*y-A_0^*y_0\|\leq\eta R_*+3r_*$,
$\|A^*y\|\leq4R_*$ and $\|u_0\|\leq R_*$.
Insert these estimates in the difference of the two rank-one
operators in \eqref{eq:ambient-exact-generator}.
All estimates are Euclidean operator estimates on the complete space.

Using \eqref{eq:ambient-ray-error}, $\eta\leq2e$,
$R_*\leq6\Theta^2$, and $\varepsilon\leq e$, the right side of
\eqref{eq:ambient-generator-gap} is at most
$5\cdot10^6e\Theta^{11}$.
The ideal generator has zero $n$ input column.
Therefore
$\Pi\mathscr B_0(y_0)(E_y-E_{y_0})=0$, and the actual reduced generator
differs from the ideal reduced generator by at most
\begin{equation}\label{eq:ambient-reduced-gap}
 65\cdot10^6e\Theta^{13}.
\end{equation}
The ideal reduced generator is explicitly
\[
 \begin{pmatrix}
 \dfrac{2P_0Q_0}{1+P_0^2}&
       -2+\dfrac{2P_0(P_0+\beta)}{1+P_0^2}\\[2mm]
 -1&0
 \end{pmatrix}
 \oplus0_{\mathcal E},
 \qquad P_0=\beta\tau^2,\quad Q_0=-2\beta\tau .
\]
Hence its active equation is
\[
 ((1+\beta^2\tau^4)V')'=2(1-\beta^2\tau^2)V,\qquad U=-V',
\]
and every ideal extra coordinate is constant.

Let $g=V_0$ be the positive scalar solution with $g(0)=1$, $g'(0)=0$.
The scalar relative propagator and the ratio estimate
$g(s)/g(t)\leq C\Theta$ bound the complete ideal reduced evolution by
\[
 K_d\frac{g(t)}{g(s)},\qquad K_d=C_d\Theta^8 .
\]
The extra identity block is included using this ratio estimate;
$g$ need not be nondecreasing. Whenever the reduced
coefficient gap is $\Delta$, the weighted Duhamel estimate gives
\[
 \frac{g(s)}{g(t)}\|U(t,s)\|
 \leq K_d+K_d\Delta
       \int_s^t\frac{g(s)}{g(r)}\|U(r,s)\|\,dr .
\]
Thus $K_d\Theta\Delta\leq1/2$ implies
$\|U(t,s)\|\leq2K_dg(t)/g(s)$.
This estimate covers arbitrary transverse initial data and forcing,
including additional modes.

It also bounds the joined high evolution in the fixed transverse
coordinates used in the history problem. At each endpoint the actual
physical field and reduced coordinates are related by
\[
 b=O D_vE_y\zeta,\qquad
 \zeta=\Pi D_v^{-1}O^*b.
\]
Their operator norms are at most $13\Theta^2$ and
$\varepsilon^{-1}$, respectively. If $R_\perp:U\to m_0^\perp$
is an isometric identification and $b=F R_\perp a_{\rm hist}$, then
\[
 a_{\rm hist}=R_\perp^*F^{-1}b,\qquad
 b=F R_\perp a_{\rm hist}.
\]
These two endpoint transformations have norms at most $\|F^{-1}\|$ and $\|F\|$.
Thus the evolution in the actual $a_{\rm hist}$ coordinates is bounded
by $26K_d\Theta^2\varepsilon^{-1}
\sup\|F^{-1}\|\sup\|F\|\,g(t)/g(s)$.
Every factor has a known polynomial bound in the source parameters.
Applying the same endpoint maps to the variation-of-constants formula gives the forced
tail bound with these factors. This is the full $d-1$ dimensional
forward operator joined to the history trace, including its action
on all additional transverse components.

For the selected primary let $\lambda\geq0$ be its actual normalized
initial slope and let $V_\lambda(0)=1$, $V_\lambda'(0)=\lambda$.
The ideal reduced initial vector is $(-\lambda,1,0)$.
At a neighboring label the discrepancy of $(U,V,Z)$ is at most
$\rho L_{\rm hist}Y_*/\varepsilon$, hence at most $e$ after the choice
\eqref{eq:ambient-error-input}. In particular the extra primary
initial component is bounded by the same history Lipschitz estimate. The remaining initial component is
determined by \eqref{eq:extra-tangent-elimination}.

The weighted Duhamel estimate applied to the difference from
$\zeta_\lambda=(-V_\lambda',V_\lambda,0)$ yields
\begin{equation}\label{eq:ambient-primary-relative}
 \|\zeta-\zeta_\lambda\|
 \leq C_de\Theta^{30}(1+\lambda)g(\tau).
\end{equation}
Indeed the initial term is $K_de\,g(\tau)$ and the integral term is
at most $2K_d^2\Theta\Delta(1+\lambda+e)g(\tau)$.
The exponent is $2\cdot8+1+13=30$.
Restoring the eliminated coordinate introduces at most two more powers
of $\Theta$, by \eqref{eq:extra-tangent-elimination} and
\eqref{eq:ambient-ray-error}.

The late-interval comparison is uniform in $\lambda$.
The second fundamental solution satisfies
\[
 \frac{V_1(\tau)-g(\tau)}{g(\tau)}
   =\int_0^\tau\frac{ds}{(1+\beta^2s^4)g(s)^2}.
\]
On $[0,1]$ the scalar equation has uniformly bounded coefficients,
so this integral at one has a fixed positive lower bound.
Consequently, for $\tau\geq1$,
$V_\lambda(\tau)\geq c(1+\lambda)g(\tau)$.
The actual normalized flux is $J=\langle y,Az\rangle$; the physical
flux for the unamplified field equals $s_0aJ$; the amplified field has
flux $\alpha s_0aJ$.
The ideal flux is
$J_0=Q_0(-V_\lambda')+(P_0+\beta)V_\lambda$.
Equations \eqref{eq:ambient-ray-error}--\eqref{eq:ambient-primary-relative}
give
\begin{equation}\label{eq:ambient-flux-relative}
 |J-J_0|/V_\lambda\leq C_de\Theta^{34}
 \quad(1\leq\tau\leq\tau_{\rm end}).
\end{equation}
For example, expand its difference into
$\langle y,(A-A_0)z\rangle+
 \langle y-y_0,A_0z\rangle+
 \langle y_0,A_0(z-z_\lambda)\rangle$.
The final term introduces two powers of $\Theta$ beyond the restored-state
bound; the other two terms are smaller.

The scalar formula gives $J_0/V_\lambda\geq\beta=\sigma^2$.
Below the threshold $x=\sigma\tau=1$, both the
$Q_0(-V_\lambda'/V_\lambda)$ contribution and $P_0$ are nonnegative.
Above the threshold $x=1$, its equivalent expression is
$x^2-\beta+2\sqrt\beta\,z/x\geq1-\beta\geq3/4$, with the
nonnegative scalar Riccati variable $z$.
Let $x_{\rm tar}\geq2$ be the target coordinate and assume
$\tau_{\rm tar}=x_{\rm tar}/\sigma\leq\Theta$.
Then $\beta\geq4/\Theta^2$.
The geometric smallness condition is
\begin{equation}\label{eq:ambient-forty-guard}
 C_de\Theta^{40}\leq1.
\end{equation}
Indeed the flux error is then at most a fixed small multiple of
$\Theta^{-6}$, which is at most $\beta/2$ after fixing the dimensional
constant. This gives a strictly positive lower bound for the late flux.
The ray, denominator, reduced-propagator and relative-state estimates above
involve smaller powers and follow from the same smallness condition.
Enlarge the fixed $C_d$ to cover their finitely many numerical constants, including
the full transverse dimensional perturbation bound. The power $60$
in the separate short-extension estimate concerns the numerical time
envelope; it is not an additional geometric-error hypothesis.

The same condition transfers the ideal uniform-size estimate to the actual
neighbor. Indeed the unamplified size is
$s_0\|D_ry\|\|D_vz\|$; the packet size is $\alpha$ times this quantity.
Subtract the corresponding central reference expression for the packet
size and apply the two state bounds. At the target the actual $q$ velocity component is at least
$V_\lambda(\tau_{\rm tar})/2$ and the normal component $N$ is at least
$1/2$. Hence the unamplified target size is at least
$s_0V_\lambda(\tau_{\rm tar})/4$.
The actual normalization $\alpha\,\mathrm{size}_{\rm target}
=\delta h_{\rm child}$ and the scalar ratio
$g(\tau)/g(\tau_{\rm tar})\leq C\Theta$ make this size error at most
$C_de\Theta^{40}\delta h_{\rm child}$, hence at most
$\delta h_{\rm child}$.
On $0\leq\tau\leq1$, the reference states are at most $C(1+\lambda)$
and the perturbation bound adds at most the same size.
The scalar target growth gives
$\alpha s_0(1+\lambda)\leq
 C\Theta\delta h_{\rm child}e^{-b/\sigma}$.
Here $b>0$ is a fixed constant from the scalar growth lemma, decreased
if needed. Thus the early size retains its exponentially small factor and a
fixed polynomial in $\Theta$ (the larger $\Theta^5$ bound suffices).
All these statements concern the actual short interval ending at
$\tau_{\rm end}$; the ideal scalar equation alone is defined globally.

\paragraph{The bounded profile height.}
The preceding normalization also establishes the polynomial bound for the
profile constant $H_0$ through cancellation of the scalar growth
against the terminal normalization.
The forward scalar profile is $g=V_0$. In physical time this means
$\gamma(t)=\alpha$ for $t\leq t_0$ and
$\gamma(t)=\alpha V_0(a(t-t_0)/\varepsilon)$ for $t\geq t_0$.
Since $V_0\leq V_\lambda$,
the scalar ratio estimate on the retained interval through the target
and its stipulated short extension gives
$g(\tau)/V_\lambda(\tau_{\rm tar})\leq C\Theta$.
Together with the actual target-size lower bound this yields
\begin{equation}\label{eq:actual-profile-height}
 H_0=\max(1,\sup\gamma)
 \leq1+C\Theta\,\delta h_{\rm child}s_0^{-1}.
\end{equation}
The constant $s_0$ is that of the actual unit initial phase:
$s_0^{-1}=\|F(t_0)^*n_{\rm new}\|\leq\|F(t_0)\|$.
It therefore has a bound in terms of the preceding particle-label parameters.
Likewise the scalar target lower bound gives
\[
 \alpha\leq
 C\Theta\,\delta h_{\rm child}s_0^{-1}
              (1+\lambda)^{-1}e^{-b/\sigma}.
\]
Thus both the polynomial height and the small history amplitude come
from normalization at the same target; the resulting source majorant
has polynomial dependence on the parent parameters.

Finally, \eqref{eq:ambient-error-input} explains how to choose a
quantitative spatial neighborhood. Only preceding unnormalized coefficient
and history Lipschitz bounds occur. Their positive-word inverse and
endpoint estimates are polynomial in the preceding particle-label data.
These bounds depend only on the parent geometry, its selected history
and $\varepsilon^{-1}$. In particular, they are independent of the new
angular width $\delta_n$. The prescribed $\rho_n$ below satisfies
the resulting neighborhood conditions.

\subsection{A common radius and polynomial source bounds}

For an ordered-word Gevrey block put
$\mathfrak m_{R,b}(n)=R^{n+b}((n+b)!)^2$. The fixed-base coefficient
constant is bounded by the explicit polynomial
\[
 S_{d,q}(R_c,C)=2^q C\sum_{r=0}^q(4(d+1)R_c)^r(r!)^2.
\]
Here $R_c$ is the known coefficient radius. The full word alphabet is
$\{1,\ldots,d+1\}$.

Indeed, at external order $n$ and base order $0\leq r\leq q$,
\[
 (n+r)!^2\leq4^{n+r}(n!)^2(r!)^2.
\]
The $(d+1)^{n+r}$ ordered coordinate choices therefore give a
coefficient-block bound
\[
 S_{d,q}(R_c,C)\,[4(d+1)R_c]^n(n!)^2.
\]
Thus the known coefficient radius is replaced once by
$4(d+1)R_c$. In the commutator argument below $R_c$ denotes this
enlarged known radius. This conversion is applied only to coefficient
tensor bounds; the forcing is already measured by its full ordered-word
block.

For a positive continuous profile $\gamma$ and a fixed base order $q$,
the normalized block used here is
\[
 \mathcal B_n^\gamma(f)=
 \sum_{I\in\{1,\ldots,d+1\}^{n}}
 \bigl\|\gamma^{-1}\partial_I f\bigr\|_{C([0,T];H^q)}.
\]
The finite sum is outside the time norm. At each fixed time the block
dominates the corresponding sum of Sobolev norms and hence the
weighted norms used in the correction estimates. Equivalent choices
of the base $H^q$ norm change only constants depending on $d,q$.

\begin{lemma}[Same-radius joined bounds]\label{lem:joined-same-radius}
There is a threshold $R_*\geq1$, polynomial in the known source coefficients,
their reciprocal coercivity bounds, inverse time lengths and the relative
propagator constant of Section~\ref{subsec:full-transverse-relative}, such that
for every $R\geq R_*$ the joined high inverse maps a forcing satisfying
$\mathcal B_n^\gamma(f)\leq A\mathfrak m_{R,b}(n)$ to velocity and
time derivative bounds
\[
 \mathcal B_n^\gamma(h_f)\leq C_{\rm v}A\mathfrak m_{R,b+3}(n),\qquad
 \mathcal B_n^\gamma(\partial_t h_f)\leq C_{\rm t}A\mathfrak m_{R,b+3}(n).
\]
Here $h_f$ denotes the joined high field. The constants and radius
threshold are independent of $A$, $b$ and $R\geq R_*$. The
mean inverse, pressure inverse and tensor corrector have analogous
fixed-shift bounds at this same radius.
\end{lemma}
\begin{proof}
The history base estimate is Lemma~\ref{lem:joined-history-polynomial}.
Its terminal trace supplies the actual forward datum. The forward
base estimate is the relative estimate in Section~\ref{subsec:full-transverse-relative}, on the
fixed packet support. Products with the known frame turn coordinate
velocities into physical velocities, and the Gram equation controls
the true time derivative.

We prove the all-order estimate by induction at the same radius. Differentiating
the actual equation produces positive-order coefficient commutators.
The elementary inequality
\[
 \binom nj(j!)^2((n-j+b)!)^2\leq((n+b)!)^2
\]
holds for every $b\geq0$. If $M$ is the fixed-base inverse or forward
constant and $B R_c^j(j!)^2$ is the known coefficient block, the normalized
positive-order sum is bounded by
\[
 MB\sum_{j\geq1}(R_c/R)^j\leq\tfrac12
 \quad\hbox{when}\quad R\geq R_c(1+2MB).
\]
Thus a factor two closes the induction at all orders, without changing
the unknown forcing radius. The subsets of an ordered word give exactly
the binomial convolution above. A dimensional factor is permitted when
converting the known coefficient tensor bounds into word bounds;
it is absorbed once into $R_c$.

The fixed-base inverse is the finite polynomial recursion in
Lemma~\ref{lem:joined-history-polynomial}. The mean form and pressure
equation have the same commutator identity and use their actual coercive
inverses. The corrector uses known coefficient products, one angular
primitive, and one spatial derivative. The primitive has operator norm bounded by
the fixed period and the coordinate sums are finite ($\leq d^3$
terms). Taking the sum of these finitely many required radii gives
a polynomial common radius. Linearity, applied after dividing the
forcing and datum by $A$, proves that the radius is independent of
$A$; uniqueness gives the zero-amplitude case. The time profile itself
is never differentiated.
\end{proof}

These estimates define an explicit polynomial source majorant. Replace
known primitive bounds by $W\geq1$, coefficient blocks by $S_{d,q}$,
fixed inverses by their $q$-step polynomial recursion, and maxima by
sums. Reciprocal coercivity, ray and Gram constants are first replaced
by the bounds proved above. Every remaining operation in the
source construction is addition, multiplication, a fixed
power or one of these fixed recursions. The resulting positive
polynomial $\Phi_d(W)$ gives the bound
\[
 \Phi_d(W)\leq A_dW^{e_d},\qquad
 e_d=\deg\Phi_d,\qquad A_d=1+\sum_j|[W^j]\Phi_d|.
\]
The exponent is therefore determined by a fixed finite recursion.
The mean and high grade recursion uses this common radius at every grade.

\subsection{The explicit nonlinear threshold}

Let $p=202$ be the tail-polynomial degree proved in
\eqref{eq:packet-tail-base}, and put
$\vartheta=1/(1000(p+1))=1/203000$. Write
$Z=k^\vartheta$, $N=\lfloor Z\rfloor$. Let $C,I,G,D,T,r_{\mathrm a}^{-1}$
be the tail, inverse, correction-growth, drift, time and reciprocal
radius constants of that packet, with $C,I,G,D,T\geq0$.
Here $r_{\mathrm a}>0$ is its analytic
radius; it is distinct from the physical localization length $\rho$ and
from the fixed mean cutoff parameter $\ell_*$. Define
\begin{equation}\label{eq:source-cap}
 H=64+\exp(1)+2^pC+I+12GT+\frac{8GTD}{r_{\mathrm a}}
                   +\left(\frac{8GT}{r_{\mathrm a}}\right)^2,\qquad
 k_*=H^{1000(p+1)}.
\end{equation}
For $k\geq k_*$ one has $H\leq Z$, $Z\geq64$ and $\log k\geq1$.
In particular
\[
 C(N+1)^p\leq C(2Z)^p\leq Z^{p+1}=k^{1/1000}.
\]
The residual and correction tolerance obey
\[
 2[2e^{-(7/10)Z\log k}]e^{3GT}\leq\tfrac12e^{-\sqrt Z},
 \qquad
 2G(D/k+e^{-\sqrt Z})T\leq r_{\mathrm a}/2.
\]
Indeed $3GT\leq Z/4$, $\sqrt Z\leq Z/8$, $\log8\leq Z/8$ prove
the first inequality. For the second, use
$8GTD\leq r_{\mathrm a}k$ and
$8GT/r_{\mathrm a}\leq\sqrt Z\leq e^{\sqrt Z}$.
This gives a polynomial frequency threshold for the exponential
correction-growth condition.

If these constants are bounded by $Q\geq1$, the explicit bound is
\begin{equation}\label{eq:polynomial-frequency-threshold}
 k_*\leq(2^p+200)^{1000(p+1)}Q^{6000(p+1)}.
\end{equation}
For $p=202$ the power of $Q$ is $1\,218\,000$. This exponent is
independent of the stage and the truncation.

\subsection{The next particle-label bound}

Use the lifted and physical fields constructed in
Lemmas~\ref{lem:graph-invariant}--\ref{lem:physical-labels} and
\eqref{eq:graph-l2-jet}. If $u$ is the corrected coordinate packet,
its lifted velocity is
\[
 \mathcal V=(k^{-1}u,\langle m_0,u\rangle),\qquad \|m_0\|=1,
\]
and its material acceleration is the field
$(\mathcal V_t+D\mathcal V\,\mathcal V)\circ\Psi$.
The cylinder time-derivative estimate and the finite product and
composition bounds define the source polynomial $C_{\rm rel}(W)$
of the lifted-flow construction. At the prescribed scales its
smallness conditions are
\[
 16C_{\rm rel}(W)\leq k^{1/2},\qquad
 \rho^{-1}\leq k^{1/2},\qquad
 \delta_{\rm corr}(k^\vartheta)\leq k^{-1}.
\]
Here $\delta_{\rm corr}$ is the correction tolerance, distinct from
the profile width. The drift cancellation gives
$BRT\leq C_{\rm rel}(W)(k^{-1}+\delta_{\rm corr})\leq1/8$.
The graph trace and physical scaling therefore give common
relative-field amplitude $M\leq k$ and radius $S\leq k^2$.
These bounds hold for displacement, material velocity and material
acceleration in $L^2$, and for displacement and material velocity
also in supremum norm. All fields come from the same lifted flow. The
preceding conditions require only the prescribed source smallness exponents.

Write $X_{\rm parent}=\mathrm{id}+D_0$, with material velocity
$V_0$ and acceleration $W_0$. In original physical labels the relative
map is $Y=\mathrm{id}+a=\rho\psi^{\rm norm}(\cdot/\rho)$, with
time derivatives $b,c$. The full particle map is $X_{\rm parent}\circ Y$,
so its three fields are
\[
 \begin{split}
 D_0^+&=D_0\circ Y+a,\\
 V_0^+&=V_0\circ Y+b+(DD_0\circ Y)b,\\
 W_0^+&=W_0\circ Y+2(DV_0\circ Y)b+
        (D^2D_0\circ Y)[b,b]+c+(DD_0\circ Y)c.
 \end{split}
\]
The parent fields are estimated in the original physical labels, with
common $L^2$ and supremum amplitude and radius $K\leq k$. Volume
preservation of $Y$ retains the parent $L^2$ factors in each composition term. The product and composition
estimates proved in the lifted-flow section give common child
$L^2$ amplitude and radius
\[
 C^+=K+M+9K^2M+36K^3M^2,\qquad
 R^+=(1+S)((1+M)16K+2)+S.
\]
For the fixed lower frequency bound $k\geq69$, these satisfy
$C^+\leq47k^5\leq k^6$ and $R^+\leq69k^4\leq k^5$.
The fixed Sobolev evaluation shift in
\eqref{eq:fixed-sobolev-jet-shift} has order $q_{\rm ev}=d+2$,
separate from the packet base order $q=2d+5$. After absorbing its fixed
dimensional constants at the prescribed frequency, it yields
supremum amplitude at most $k^{6q_{\rm ev}+7}$ and radius at most
$k^6$. Thus the single bound
\[
 K^+=k^{10(d+4)}
\]
controls all three child fields in both norms, since the sum of one,
the two $L^2$ bounds and the two supremum bounds is at most
$5k^{6q_{\rm ev}+7}\leq k^{6q_{\rm ev}+8}\leq k^{10(d+4)}$.
In particular, the acceleration has the supremum jet bounds
needed for the next $F_2^+$ source. These are the fixed
source-polynomial and frequency conditions included in the single
scale choice.

The full inverse deformation is polynomial as well. Since
$\det F^+=1$, its inverse is the adjugate. If
$\|D^nF^+\|\leq A R^n(n!)^2$, expansion of each cofactor into
$(d-1)!$ products and the Gevrey product constant 3 give
\[
 \|D^n(F^{+,-1})\|
 \leq d(d-1)!\,3^{d-2}A^{d-1}R^n(n!)^2.
\]
The factor $d$ bounds operator norm by the largest entry.
The time derivative satisfies
$(F^{+,-1})_t=-F^{+,-1}F_1^+F^{+,-1}$.
Also $F_1^+=D_xV_0^+$ and $F_2^+=D_xW_0^+$.
Thus the source coefficients have fixed polynomial bounds in the
new particle-label parameter. Fix, for instance,
\[
 D_d=10(d+4)
\]
for the particle-label invariant $K^+\leq k^{D_d}$.
The inverse and coefficient bounds are fixed polynomials in $K^+$,
possibly with exponents different from $D_d$.

\subsection{Uniform scales and a single frequency choice}

Set $s_n=J+n$ and
\[
 x_0=X,\quad x_{n+1}=s_n^2x_n,\quad
 k_n=e^{x_n/s_n^2},\quad h_n=e^{x_n/s_n^5},\quad
 \delta_n=e^{-x_n/s_n^3},\quad\rho_n=e^{-x_n/s_n^{7/2}}.
\]
Here $n=j-J$ and $x_n=x_{j-1}$ in the notation of
Section~\ref{sec:iteration}; $\rho_n$ is the physical localization,
not the fixed mean cutoff $\ell_*$. For $n\geq1$ the identities
\[
 \log k_{n-1}=\frac{x_n}{(s_n-1)^4},\quad
 \log h_{n-1}=\frac{x_n}{(s_n-1)^7},\quad
 \log\rho_{n-1}^{-1}=\frac{x_n}{(s_n-1)^{11/2}}
\]
are exact. If $K_n\leq k_{n-1}^{D_d}$, every preceding geometric constant in
\eqref{eq:ambient-error-input} is bounded by
$Cs_n^a x_n^b e^{c x_n/(s_n-1)^4}$.
For sufficiently large fixed $J$ and then $X$,
\[
 \rho_n Cs_n^a x_n^b e^{c x_n/(s_n-1)^4}
 \leq e^{-x_n/(2s_n^{7/2})}.
\]
The additional factors in the ambient smallness condition have the same
inherited scale: $\Theta_n\leq C(1+s_n^2x_n^2)$,
$s_0^{-1}\leq\|F(t_0,0)\|$, and
$\varepsilon^{-1}\leq C\sqrt{h_{n-1}}$.
They therefore contribute only polynomial factors in $s_n,x_n$ or an
exponential with denominator $(s_n-1)^7$, which is already absorbed
by the displayed denominator $(s_n-1)^4$. The terms
$\varepsilon\Theta G^2$ and $e_{\rm old}$ in
\eqref{eq:ambient-error-input} are controlled by the separate shear
separation and preceding frequency-error comparisons in
Section~\ref{sec:iteration}; they are not made small by localization.
At $n=0$ the activation time is zero, and the forward high construction
is used. This case contains no reciprocal history-length factor.
The seed predecessor has polynomial bounds in $X$. The normal-stage
factor $e^{-X/J^{7/2}}$ dominates those bounds for the same final large choice of $X$.

Consequently all neighborhood conditions hold at the prescribed
localization length $\rho_n$.

The full source list additionally includes $\delta_n^{-1}$ and the
new scalar height. Its positive-polynomial bounds therefore have the
form
\begin{equation}\label{eq:uniform-parameter-envelope}
 \mathcal P_n=Cs_n^a x_n^b
              \exp\left(\frac{c x_n}{(s_n-1)^3}\right).
\end{equation}
In particular the radius, $H_0$, reciprocal analytic radii, history
and mean inverse constants, and the constants in \eqref{eq:source-cap}
are bounded by fixed powers of \eqref{eq:uniform-parameter-envelope}.

\begin{lemma}[Uniform domination]\label{lem:uniform-source-domination}
For a finite list of fixed $A,C,c,a,b,N>0$ and $\gamma>0$, one common
choice of $J$ and then $X$ ensures
\[
 A\mathcal P_n^N\leq k_n^\gamma\qquad(n\geq0)
\]
for every entry of the list.
\end{lemma}
\begin{proof}
Choose $J$ so that
$Nc/(s-1)^3\leq\gamma/(2s^2)$ for $s\geq J$;
$J\geq16Nc/\gamma+2$ suffices. Increase $X$ until
\[
 \log(AC^N)+Na\log s_n+Nb\log x_n
 \leq\frac{\gamma x_n}{2s_n^2}
\]
uniformly in $n$. To see the uniformity, use
$x_n=X\prod_{r<n}(J+r)^2$.
The sequences
$s_n^2\log s_n/\prod_{r<n}(J+r)^2$ and
$s_n^2\log(\prod_{r<n}(J+r)^2)/\prod_{r<n}(J+r)^2$
are bounded, since the denominator eventually dominates every
polynomial in $n$ while its logarithm grows at most
$2n\log(J+n)$. The resulting bound is
$O_J((1+\log X)/X)$ and tends to zero.
Adding the two inequalities proves the assertion. Take the maximum
of the finitely many choices of $J,X$.
\end{proof}

Apply this lemma simultaneously to the source constants, the neighborhood
conditions, \eqref{eq:polynomial-frequency-threshold}, the lifted-flow smallness
conditions and their numerical constants. It gives the actual nonlinear
correction at the preselected $k_n$ and then
$K_{n+1}\leq k_n^{D_d}$. The exponent $D_d$ is fixed first;
the source polynomials and their degrees are then fixed; $J,X$ are
chosen once for the resulting finite list. The same exponents and starting
index apply to every subsequent child.

The zero-history initialization adds seven fixed constants to this same
finite list: the actual forward family, weighted carrier localization,
support geometry, gradient and pressure errors, relative particle
fields, and the initial amplitude ratio. Their exponents and constants
are fixed before the single choice of $J$. The finitely many remaining
seed and first forward-stage thresholds depend on $d,J$ and
are met by increasing $X$ alone. The seed cutoff normalization, the two
initial mean reserves, and the preserved core difference are specified
in Section~\ref{sec:iteration}. Thus the first positive-history state
and every successor use the same $J,X$ and the same prescribed
frequencies. The induction then applies to the joined variational-history
family and the quantitative neighborhood constructed in this section.
\section{Initial increments at every Sobolev order}
\label{sec:initial-summability}

We use the joined displacement-history and forward inverse of
Lemma~\ref{lem:packet-joined-high}. When the activation time is positive,
the forced high coefficients generally have nonzero initial velocities;
all of them enter the estimates below.

Fix the dimension $d$ and an admissible integer base Sobolev order $s$ for the
packet estimates. Constants below may depend on $d$ and on a fixed
derivative order, but the frequency condition will not depend on that
derivative order. Write $\mathcal P\ge1$ for a bound on the source primitives,
including the coefficient and inverse constants, the terminal datum, the
reciprocal phase scale, the fixed mean cutoff constant, and the common Gevrey
radius. Section~\ref{sec:parameter-envelope} bounds these primitives for the
constructed sources, uniformly along the prescribed scale sequence.

\subsection{The initial coefficient sums}
Put $\kappa=k^{-1}$ and let $N$ be the chosen truncation. At time zero,
the parent flow is the identity and the nonlinear correction vanishes.
The initial increment is
\begin{align}
 u_{\mathrm{osc},0}(a)
 &=\rho\sum_{p=1}^{N}
    \bigl[\kappa^p A_p+\kappa^{p+1}C_p\bigr]
       (0,a/\rho,km_0\cdot a/\rho),\label{eq:initial-osc-sum}\\
 u_{\mathrm{mean},0}(a)
 &=\rho\sum_{p=2}^{N}\kappa^p B_p(0,a/\rho),
 \qquad
 \Delta u_0=u_{\mathrm{osc},0}+u_{\mathrm{mean},0}.
 \label{eq:initial-mean-sum}
\end{align}
Here $C_p$ is the full ambient antisymmetric-potential corrector and $B_p$
is independent of the angle. The estimates include every high coefficient, its tensor corrector,
and every mean coefficient.
Formula~\eqref{eq:initial-osc-sum} extends
\cite[\S3.9, (3.60)]{OAI} to the full ambient tensor construction.

Let $\alpha>0$ be the selected primary amplitude. The joined profile is
\[
 \gamma(t)=\alpha\quad(0\le t\le t_0),\qquad
 \gamma(t)=\alpha g(t)\quad(t_0\le t\le T),\qquad g(t_0)=1,
\]
and put $H_0=\max(1,\alpha,\sup\alpha g)$. The coefficient induction
assigns the oscillatory profiles $\gamma H_0^{2p-2}$ and the mean profiles
$H_0^{2p-2}$.

The factor $\alpha$ on the history interval follows from linearity.
The primary terminal displacement is $\alpha$ times the normalized compact
datum, and the displacement inverse is linear. In every higher zero-mean
force, terms containing only means cancel under angle-mean subtraction.
Every remaining term contains a previous oscillatory factor or a derivative
of one. The term involving the newly solved mean is
$-(m\cdot B_p)\partial_\theta A_1$, which also contains this factor.
Extra oscillatory factors are bounded by the indicated powers of $H_0$.
The history inverse and the angular tensor primitive therefore
retain $\alpha$ at time zero. This factor is preserved at each step of the graded construction
by linearity of its high solver.

Suppose first that $t_0>0$. In fixed transverse coordinates, every
forced history solves
$\xi(0)=\xi(t_0)=0$ and has velocity $A=FR_\perp\xi_t$.
The primary instead has $\xi(0)=0$ and
$\xi(t_0)=\alpha\chi_{\rm in}(y)f_\delta(\theta)\xi_T$, where $\xi_T$ is the
selected terminal coordinate. It is obtained from the physical
central endpoint by the inverse frame; its norm is included in $\mathcal P$.
Subtracting the affine extension
$(t/t_0)\alpha\chi_{\rm in}f_\delta\xi_T$ reduces this primary problem to the
same zero-endpoint inverse. All spatial and phase words of the forced
problem still have zero endpoint values. For the primary, differentiating
this explicit affine extension supplies exactly the differentiated
boundary datum; the differentiated boundary traces are therefore explicit.
The forward solution starts with the trace $FR_\perp\xi_t(t_0)$.
It agrees there with the history solution, and the common velocity equation
then gives equality of the time derivatives at the junction. These are
the boundary conditions of Lemma~\ref{lem:packet-joined-high}.
When $t_0=0$, the primary instead has its prescribed forward datum
and the forced high fields have zero initial velocity. The forward
estimates give the same bounds below, without a history inverse or
reciprocal factor $t_0^{-1}$. This applies both to the exceptional seed
and to the first normal stage $j=J$; positive history begins at $j=J+1$.

The all-word estimate uses the same inverse. After
translation in the label and angle, write the equation as
$L(a)u(a)=f(a)$. At a fixed admissible base Sobolev order let its inverse constant
be $M$, and let the known coefficient derivative bounds be
$B R_c^j(j!)^2$. Differentiating a word of length $n$ gives the exact
Leibniz identity
\[
 L\partial^\omega u=\partial^\omega f-
 \sum_{\varnothing\ne S\subset\{1,\ldots,n\}}
 (\partial^{\omega_S}L)\partial^{\omega_{S^c}}u.
\]
The homogeneous boundary space is unchanged, so the same inverse applies
to each term. In the block norm this yields a binomial convolution.
For $b\ge0$,
\[
 \binom nj(j!)^2((n-j+b)!)^2\le((n+b)!)^2.
\]
Consequently a forcing bound $F R^{n+b}((n+b)!)^2$ gives
$2MF R^{n+b}((n+b)!)^2$ for the solution as soon as
$R\ge R_c(1+2MB)$: the positive-order convolution contributes at most
$MB\sum_{j\ge1}(R_c/R)^j\le1/2$.
The initial finite base inverse is obtained by the finite Sobolev
commutator recursion. Its order is fixed by $d$, not by the packet grade.
Time traces follow from the $H^1$ equation and have a polynomial bound
in the reciprocal history length and known frame bounds. The same estimate
on the forward interval uses its relative propagator and the
history trace. This proves preservation of the unknown forcing radius;
known coefficient conversion enlarges $R_c$ only once.

The displacement argument uses the full transverse space of dimension
$d-1$. If $R_\perp$ has orthonormal columns, $\eta=FR_\perp\xi$, and
$\eta(t_0)=0$, then
\[
 \int_0^{t_0}\bigl(|\eta_t|^2-\langle H\eta,\eta\rangle\bigr)\,dt
 \ge \left(1-\frac{K_+t_0^2}{2}\right)
                       \int_0^{t_0}|\eta_t|^2\,dt
 \ge\frac12\int_0^{t_0}|\eta_t|^2\,dt.
\]
The one-sided Poincar\'e inequality here follows by integrating
$\eta(t)=-\int_t^{t_0}\eta_t(r)\,dr$. Conversion to $\xi$ uses only the
known bounds for $F^{-1}$ and its derivative. Thus the inverse constant is
polynomial in those bounds in every transverse dimension.

\subsection{Separating derivative order from expansion order}
For a fixed integer $m\ge0$, the coefficient estimates need derivatives
only through $m+s+1$. Their shift is at most $100p$; replacing this fixed
coefficient by a larger dimensional constant makes no change to the
argument. For $b=m+s+1$,
\[
 (b+100p)!\le2^{b+100p}b!(100p)!,\qquad
 ((b+100p)!)^2\le C_b(Cp)^{200p}.
\]
Also $R^{b+100p}=R^b(R^{100})^p$. Consequently all dependence on $m$
can be put in a prefactor $C_{d,m}\mathcal P^{c_{d,m}}$, while the
expansion-order bound has the form
\begin{equation}
 B_N^{p+1},\qquad B_N=\mathcal P^{c_d}(C_dN)^{L_d}.
 \label{eq:initial-grade-base}
\end{equation}
with one finite exponent $L_d$ independent of $m$.
For the displayed shifts one may take $L_d=300$ after enlarging constants.
Choose $N=\lfloor k^\vartheta\rfloor$ with
$L_d\vartheta<1/100$. A single fixed source-to-frequency condition then
ensures $B_N\le k^{1/100}$.

At time zero the cylinder Sobolev norms of $A_p$ and $C_p$ are bounded by
\[ \alpha C_{d,m}\mathcal P^{c_{d,m}}B_N^{p+1}. \]
The mean has the same bound without $\alpha$. For $k\ge4$,
\[
 \sum_{p=1}^{N}k^{-p}B_N^{p+1}\le C,
 \quad
 \sum_{p=1}^{N}k^{-p-1}B_N^{p+1}\le C,
 \quad
 \sum_{p=3}^{N}k^{-p}B_N^{p+1}\le Ck^{-2}.
\]
These follow by summing the geometric series with ratio $k^{-99/100}$.
In the last sum its first power is $k^{-3+4/100}\le k^{-2}$.
The grade-two mean is estimated separately at its fixed order, giving
$C_{d,m}\mathcal P^{c_{d,m}}k^{-2}$ as well. This fixed-order estimate preserves
the power $k^{-2}$ in the mean bound.

For $T_kf(y)=f(y,km_0\cdot y)$ and $|m_0|=1$, the chain rule and
$H^1(\mathbb T)\hookrightarrow L^\infty(\mathbb T)$ give
\[
 \|T_kf\|_{H^m(\R^d)}
       \le C_{d,m}(1+k)^m
                    \|f\|_{H^{m+1}(\R^d\times\mathbb T)}.
\]
Indeed each spatial derivative is a product of the commuting operators
$\partial_{y_i}+km_{0,i}\partial_\theta$. Apply the periodic one-dimensional
estimate to each resulting derivative at fixed $y$, then integrate in $y$.
Exactly one additional angle derivative is needed. For the mean, angle
independence eliminates the additional graph derivative.
Finally, physical scaling multiplies the $L^2$ norm of an order-$r$
derivative by $\rho^{1+d/2-r}\le\rho^{-m}$ for $r\le m$ and
$0<\rho\le1$. We obtain the estimates
\begin{align}
 \|u_{\mathrm{osc},0}\|_{H^m}
 &\le C_{d,m}\rho^{-m}k^m\alpha \mathcal P^{c_{d,m}},\label{eq:joined-initial-high}\\
 \|u_{\mathrm{mean},0}\|_{H^m}
 &\le C_{d,m}\rho^{-m}k^{-2}\mathcal P^{c_{d,m}}.
 \label{eq:joined-initial-mean}
\end{align}
The selected central amplitude satisfies
\[
 \alpha\le \mathcal P^{c_d}e^{-b x}
\]
for a fixed $b>0$, by the central target growth, with
the normalized terminal and reciprocal phase constants included in $\mathcal P$.
In the scalar indexing below, $n=j-J$ corresponds to normal stage $j$
of Section~\ref{sec:iteration}: $a_n=j$, $x_n=x_{j-1}$,
$k_n=k_j$, and $\rho_n=\rho_j$. The inherited geometric parameter
is therefore $x_n$, while the new target is $x_{n+1}$. The amplitude
factor is $e^{-b x_n}$. This exponential factor applies to
all initial high coefficients, including those produced by positive
history intervals.

\subsection{One scale sequence gives all orders of summability}
\begin{lemma}[Scalar summability]\label{lem:initial-scalar-sum}
Fix $J\ge2$, $X\ge1$, and define
\[
 a_n=J+n,\quad x_0=X,\quad x_{n+1}=a_n^2x_n,
 \quad k_n=e^{x_n/a_n^2},\quad \rho_n=e^{-x_n/a_n^{7/2}}.
\]
Here $\rho_n$ is the physical packet length; it is distinct from the
fixed physical mean-cutoff parameter $\ell_*$ and from the Gevrey
weight radius used in the cylinder estimates. Let
\[
 \mathcal P_n=C a_n^p x_n^q
           \exp\left(\frac{c x_n}{(a_n-1)^3}\right)
\]
with fixed $C\ge1$, $c\ge0$, $p,q\in\N$ and $b>0$. For every fixed
$m,\nu\in\N$, both series
\[
 \sum_n\rho_n^{-m}\mathcal P_n^\nu k_n^m e^{-b x_n},
 \qquad
 \sum_n\rho_n^{-m}\mathcal P_n^\nu k_n^{-2}
\]
converge.
\end{lemma}
\begin{proof}
We have $x_n\ge X4^n$ and
$\log x_n\le\log X+2n\log(J+n)$. Therefore, for every fixed $r\ge0$,
\[
 \frac{a_n^r(1+\log a_n+\log x_n)}{x_n}\longrightarrow0.
\]
The logarithm of the first summand, divided by $x_n$, tends to $-b$:
its other exponential terms are $m/a_n^2$, $m/a_n^{7/2}$ and
$c\nu/(a_n-1)^3$, and its logarithmic terms vanish. Eventually the summand
is at most $e^{-b x_n/2}$, a summable sequence.

For the second summand, divide its logarithm by $x_n/a_n^2$. The result
tends to $-2$, since
\[
 ma_n^{-3/2}\to0,\qquad
 c\nu\frac{a_n^2}{(a_n-1)^3}\to0,
\]
and the logarithmic terms again vanish. Eventually it is at most
$e^{-x_n/a_n^2}\le e^{-n}$, proving convergence.
\end{proof}

The choices of $J,X$ in this lemma are independent of $m$. Only the
starting index of the convergence comparison depends on $m$; the
construction itself uses the same scales for every Sobolev order.
The same proof permits a localization scale whose reciprocal is bounded
by a polynomial times $\exp(c_0x_n/a_n^\beta)$ for any fixed $\beta>2$.
The quantitative localization estimate is therefore part of the
source construction.

To verify that the prescribed frequency satisfies all construction
conditions, consider any finite list $A\mathcal P_n^\nu\le k_n^\theta$ with
$\theta>0$. Choose $J$ so that
$c\nu a^2/(a-1)^3\le\theta/2$ for all $a\ge J$ and every condition.
Writing $x_n=Xb_n$ with $b_n\ge J^{2n}$, the sequences
$a_n^2/b_n$, $a_n^2\log a_n/b_n$ and $a_n^2\log b_n/b_n$ are bounded.
Hence a single increase of $X$ absorbs all the remaining logarithmic terms,
uniformly in $n$. This proves the conditions at the prescribed $k_n$.
Only finitely many construction conditions enter at fixed $d$; the fixed-$m$
prefactors in \eqref{eq:joined-initial-high}--\eqref{eq:joined-initial-mean}
do not introduce further construction conditions.

\subsection{The smooth compact limit}
Section~\ref{sec:parameter-envelope} bounds the stage-$n$ source
primitives by $\mathcal P_n$ and verifies the prescribed localization
scale for the joined construction. The preceding estimates consequently give
\[
 \sum_n\|\Delta u_n(0)\|_{H^m}<\infty
 \qquad\text{for every integer }m\ge0.
\]
For each spatial order $r$, take $m>r+d/2$. The Sobolev
embedding in dimension $d$ gives uniform convergence of the derivative
series through order $r$. Include the base solution, the seed increment and
the first normal zero-history forward increment in $u_{\mathrm{seed},0}$.
The latter two are fixed finite packet increments, including their
mean fields and terminal tensor correctors. The remaining sum starts
with the first positive-history joined increment; this finite index
shift changes neither the convergence proof nor the already fixed
$J,X$. Consequently the single sum
$u_0=u_{\mathrm{seed},0}+\sum_{n\geq1}\Delta u_n(0)$ is smooth, has the termwise
derivatives, and is divergence free. The fixed physical mean-cutoff
construction puts the supports of all these initial increments
in one compact ball. Enlarging it to contain the base support proves
$u_0\in C^\infty_{c,\sigma}(\R^d)$.
\section{The scale choice and the finite-stage induction}
\label{sec:iteration}
The construction uses one prescribed scale sequence. The quantitative source estimate supplies a fixed exponent $D_d$
for the particle-map invariant and a finite list of polynomial construction
conditions. All choices in this section are made after those exponents have
been fixed.

\subsection{The base flow and the order of choices}
Choose an oriented orthonormal basis $(p,q,n)$ of $E$ and let the full
matrix $L_B$ vanish on $E^\perp$ and satisfy
\[
 L_Bq=p+X^{-2}n,\qquad L_Bp=L_Bn=0.
\]
It is trace free. Lemma~\ref{lem:potential}, with a fixed radial cutoff,
gives a compact, odd, reflection-equivariant solenoidal datum equal to
$L_Bx$ near zero. Its Sobolev and Gevrey bounds are uniform for $X\geq1$.
Lemma~\ref{lem:local} gives its actual Euler solution $u_B$ on a uniform
short interval. Its Gevrey bounds also persist on a uniform shorter
interval: with fixed base order $q_d> (d+1)/2+2$ and radius $r(t)$,
set
\[
 \mathcal X_r=\sum_{n\geq0}\frac{r^n}{(n!)^2}
                  \sum_{|I|=n}\|\partial_Iu_B\|_{H^{q_d}},
 \qquad
 \mathcal Y_r=\sum_{n\geq0}\frac{nr^n}{(n!)^2}
                  \sum_{|I|=n}\|\partial_Iu_B\|_{H^{q_d}}.
\]
At each finite cutoff the differentiated transport cancellation and
factorial product convolution give
$\mathcal X_r'\leq C_d\mathcal X_r^2+
(r'/r+C_d\mathcal X_r/r)\mathcal Y_r$.
Choose $r'$ a sufficiently negative fixed constant. A short-time
bootstrap bounds $\mathcal X_r$ and keeps $r>0$ uniformly in the
cutoff. Passing to the increasing sums proves the claim. The projected
Euler equation controls the time derivative and material acceleration
at a smaller radius. The integral flow equation then bounds the actual
particle map and its first two time derivatives in the same manner.
This proof uses the Sobolev embedding at an order appropriate to the fixed dimension.
Write $K_B$ for a uniform upper bound on its pressure Hessian.

First fix $d$, the packet base order and all finite inverse, coefficient,
flow and source exponents, including $D_d$. Fix $C_M,C_H$ large enough
for the central gradient and Hessian induction below. The global
coefficient bounds are separately supplied by the actual particle-label
invariant; the global pressure bound is one-sided. These fix the
endpoint constant and then its small threshold $\zeta$.
Choose $D_0$ sufficiently large for the first geometric comparison;
$D_0=1000$ suffices for the displayed exponent $60$.
All seed mean-cutoff and boundary bounds are fixed powers of $X$.
Choose $D_k\geq1000$ larger than the resulting seed polynomial exponent
divided by each of the finitely many frequency margins, including the
seed gradient, pressure and particle-label estimates. Fix also
$C_{\rm init}=2C_{\rm lit}$, where
$C_{\rm lit}=1+\pi C_{\chi,d}$ and $C_{\chi,d}$ is the fixed seed cutoff
derivative bound. This bounds the combined contribution of the two
initializations to $B_c-B_e$.

The same finite scale choice includes the twelve later-stage estimates and
seven additional estimates for the first normal forward stage: its finite
family and correction, weighted carrier localization, support geometry,
gradient error, pressure error, relative particle fields, and initial
amplitude ratio. All seven are fixed source polynomials multiplied by
the prescribed frequency or exponential damping. In particular, the last
condition is
\[
 [1+c_d(1+\pi C_{\chi,d})]\,\alpha_J/\delta_J\leq1/64.
\]
Choose the starting index $J$ once for this combined list, and then
increase $X$ to meet the seed, short-time, support and reserve
thresholds. These thresholds depend only on $d,J$ and constants already
fixed before $J$. No second choice of $J$, frequency or exponent is made
in either initialization or the later induction.

Put $h_{J-2}=1$, $x_{J-1}=X$, and $x_j=j^2x_{j-1}$ for $j\geq J$.
The auxiliary value $h_{J-2}$ records the uniform scale of the base
solution. Prescribe
\begin{align}
 h_{J-1}=X^{D_0},\quad k_{J-1}=X^{D_k},\quad
 \rho_{J-1}=s_{\rm cut,d}^{-1}X^{-D_0},\quad \delta_{J-1}=X^{-(D_0+10)},
 \label{eq:seed-scales}\\
 h_j=e^{x_{j-1}/j^5},\quad k_j=e^{x_{j-1}/j^2},\quad
 \delta_j=e^{-x_{j-1}/j^3},\quad
 \rho_j=e^{-x_{j-1}/j^{7/2}}\quad(j\geq J).
 \label{eq:normal-scales}
\end{align}
Here $s_{\rm cut,d}>0$ is the fixed dimensional cutoff scale, normalized so that
$a_d=1/(3s_{\rm cut,d})$. The seed frequency is fixed by its polynomial source
bound in $X$; the source-envelope lemma establishes the required estimates at the prescribed normal
frequencies. Use the explicit physical mean-cutoff parameters
\[
 \ell_*=1,\qquad r_*=X^{-D_0}.
\]
For large $X$, $r_*\leq1/4$. The splitting ball of radius
$a_dr_*/\ell_* = \rho_{J-1}/3$ contains the oscillatory seed support
of radius $\rho_{J-1}/4$. The normalized seed mean has parameters
$\ell=\rho_{J-1}$ and $r=r_*$, hence splitting radius $1/3$.
These parameters are fixed functions of the seed value before the final
choice of $X$; their physical values are retained throughout the
iteration and are not replaced by the shrinking $\rho_j$.
Their reciprocals and all subsequent boundary penalties are bounded by
fixed powers of $X$. The multiplier in the full mean reserve is at most
$1+c_d$, a dimensional constant fixed before the scale choice.
The first normal support condition follows from
$\rho_Js_{\rm cut,d}X^{D_0}\leq1$: for fixed $d,J$, the left side tends to zero
as $X\to\infty$. Thus this condition only enlarges the final lower bound on the seed parameter.

\subsection{Time and geometric comparisons}
Let
\[
 W_j=\frac{3x_jx_{j-1}}{\sqrt{h_{j-1}}},\qquad
 S_{J-1}:=S_{\rm base}=2W_J=6J^2X^{2-D_0/2},\qquad t_{J-1}=0.
\]
At stage $j$, the preceding solution determines $a_j,\beta_j$.
Define
\begin{equation}\label{eq:actual-times}
 t_j=t_{j-1}+\frac{x_j}{\sqrt{\beta_ja_jh_{j-1}}},
 \qquad S_j=t_j+2W_{j+1}.
\end{equation}
When $a_j,\beta_jx_{j-1}^2\in[1/2,2]$,
$W_j/6\leq t_j-t_{j-1}\leq2W_j/3$.
For $j>J$ the exact identities
\[
 \log k_{j-1}=\frac{x_{j-1}}{(j-1)^4},\qquad
 \log h_{j-1}=\frac{x_{j-1}}{(j-1)^7}
\]
show that $W_{j+1}/W_j$ is arbitrarily small uniformly after $J$:
its logarithm is
\[
 2\log j+2\log(j+1)-\frac{x_{j-1}}{2j^5}
                    +\frac{x_{j-1}}{2(j-1)^7}.
\]
At $j=J$ the last term is instead $(D_0/2)\log X$.
The negative term dominates in both cases. Hence $S_j\leq S_{j-1}$,
and with the numerical envelope
$\widehat\Theta_j=C(1+x_jx_{j-1})$ the extra scaled interval obeys
$\Delta\tau_j=2\sqrt{a_jh_{j-1}}W_{j+1}\leq\widehat\Theta_j^{-60}$.
The actual central parameter is
$\Theta_j=H_j^{\rm sc}=x_j/\sqrt{\beta_j}+\Delta\tau_j$
and satisfies $\Theta_j\leq\widehat\Theta_j$. In particular
$\Delta\tau_j\leq\Theta_j^{-60}$. The superscript distinguishes
this scaled horizon from a pressure Hessian.
Also $t_J\geq S_{\rm base}/12$, so at positive-history stages
$t_{j-1}^{-1}\leq12S_{\rm base}^{-1}\leq h_{j-1}$.

Use $E_*=k_{j-1}^{-1/4}$, $G_*=1+h_{j-2}$ and
$K_h=k_{j-1}^{D_d}$. Each part of \eqref{eq:geometric-error} satisfies
its smallness condition at the same scale choice. For the spatial term the
logarithm is at most
\[
 -\frac{x_{j-1}}{j^{7/2}}+
       \frac{c_dD_dx_{j-1}}{(j-1)^4}
       +\frac{Cx_{j-1}}{(j-1)^7}+C\log x_{j-1}
 \leq-\frac{x_{j-1}}{2j^{7/2}}.
\]
This estimate uses the inherited geometric bound, which excludes the new
$\delta_j^{-1}$; using the complete source list here would not imply the required smallness.
The old frequency error has logarithm
$-x_{j-1}/(4(j-1)^4)$. For $j\geq J+2$,
\[
 \log(\epsilon G_*^2)\leq
 -\frac{x_{j-1}}{2(j-1)^7}+
 \frac{2x_{j-1}}{(j-1)^2(j-2)^7}+C
 \leq-\frac{2x_{j-1}}{5(j-1)^7}+C.
\]
For $j=J$, $\epsilon\widehat\Theta_J^{61}\leq C(J)X^{122-D_0/2}$;
for $j=J+1$ the new exponential dominates the preceding polynomial shear.
These cover the two exceptional stages. Multiplying by any of the
fixed powers of $x_{j-1},\widehat\Theta_j$ still gives the required
smallness, including the actual $e\Theta_j^{40}$ condition.
The same comparisons give
\begin{equation}\label{eq:shear-separation}
 h_j/h_{j-1}^2\gg1,\qquad
 \frac{\sqrt{h_{j-1}}}{x_{j-1}x_j}\gg G_*+1,\qquad
 \rho_j<r_{\rm core}.
\end{equation}
Here $r_{\rm core}>0$ is the fixed physical seed core radius.
Each use of $\gg$ in this proof denotes a fixed finite inequality whose
constant was chosen before $J,X$.

\subsection{Summable error bounds}
For positive-history stages let $\mathcal P_j^{\rm par}\geq1$
bound the parent coefficient and geometric bounds, including $C_1,C_I$,
but excluding the new reciprocal phase scale. The particle-label
invariant and determinant-one inverse formula give
\[
 \mathcal P_j^{\rm par}
 \leq Cj^a x_{j-1}^b\exp\!\left(\frac{c x_{j-1}}{(j-1)^4}\right).
\]
At the first normal stage its predecessor bounds are fixed polynomials
in $X$. The full packet envelope $\mathcal P_j$ also includes
$\delta_j^{-1}$ and has denominator $(j-1)^3$ as in
\eqref{eq:uniform-parameter-envelope}.

The relative size estimate for the selected packet bounds the history,
early-time and initial primary errors by a fixed polynomial times
$h_j e^{-b x_{j-1}}$. After enlarging that polynomial to cover all the
actual initial fields, a common bound is
\[
 b_j=C_d\mathcal P_j^{c_d}h_j e^{-b x_{j-1}}
                         \leq e^{-b x_{j-1}/2}.
\]
The logarithm of the prefactor divided by $x_{j-1}$ is at most
$C_d/(j-1)^3+o(1)$; the seed predecessor contributes only a polynomial
in $X$. This comparison is ensured by the same finite choice of $J,X$.

The late upper-pressure estimate uses the actual global parent strain
bound $M_{\rm all}\leq3C_1C_I\leq3(\mathcal P_j^{\rm par})^2$.
The pressure sign \eqref{eq:positive-flux} on the packet support and
$f_\delta'\geq-1$ therefore bound its late contribution by
$C_d(\mathcal P_j^{\rm par})^2\delta_jh_j$.
The large positive part of $f_\delta'$ gives a nonpositive leading
Hessian contribution; only its bounded negative part enters this upper
bound. Outside the support the leading tensor is zero. In particular,
no global sharp bound of size $h_{j-1}$ on the parent gradient is needed.
Since
\[
 \delta_jh_j\leq e^{-x_{j-1}/(2j^3)},
\]
the parent-only polynomial factor is absorbed by the stronger damping
with denominator $j^3$. For large fixed $J$ and then $X$, the resulting
contribution is at most $e^{-x_{j-1}/(4j^3)}$. This absorption would not be
valid with an envelope containing the new $\delta_j^{-1}$.

The exact packet gradient and Hessian remainders are bounded by
$k_j^{-1/4}$. Thus nonnegative bounds for the increments are
\begin{equation}\label{eq:actual-increment-budgets}
 e_{p,j}=C_d\bigl(b_j+(\mathcal P_j^{\rm par})^2\delta_jh_j
                         +k_j^{-1/4}\bigr),\qquad
 e_{i,j}=C_d(b_j+k_j^{-1/4}).
\end{equation}
The latter applies to the complete initial finite family, including all
means and terminal correctors. For each fixed $A>0$,
$x_{j-1}/j^A$ increases at least geometrically for large $J$.
The damped error terms and fixed negative powers of $x_{j-1}$ consequently
have arbitrarily small sums after increasing $X$. The five terms in
the full mean-reserve increment are bounded by $1/64$ each;
their total is at most $5/64$.

The initial core term in that reserve is retained. With the fixed
cutoff above and $B_c=O_d(X^{D_0})$, it satisfies
\[
 c_dB_cr_*^dS_{\rm base}
     =O_d\!\left(J^2X^{2+(1/2-d)D_0}\right)\longrightarrow0.
\]
The remaining seed reserve terms also tend to zero. The
seed solution is chosen with full reserve at most $1/8$, and the first
normal forward step contributes at most another $1/8$. Later finite prefixes
are therefore at most $1/4+5/64=21/64<1/2$, on their actual horizons.

Boundary penalties remain polynomial by an exact difference invariant,
not by a global sharp initial-gradient assertion. After the two
initializations,
\[
 B_c-B_e\leq C_{\rm init}X^{D_0}.
\]
Every subsequent joined update adds the same $e_i$ to $B_c,B_e$, so
this inequality persists. On a positive-history horizon $T$ the full
reserve implies $B_eT\leq1/2$, while $T\geq S_{\rm base}/12$.
For the already chosen large $X$ these give
\[
 B_e\leq2X^{D_0},\quad
 B_c\leq(C_{\rm init}+2)X^{D_0},\quad
 |\mathcal L|\leq C_{1,d}(C_{\rm init}+2)X^{D_0}.
\]
These are the boundary bounds used in the next source envelope.

\subsection{Seed and induction step}
Set $u_{J-2}=u_B$. Apply Lemma~\ref{lem:exact-packet-child} on
$[0,S_{\rm base}]$ with the seed scales, $m_0=p$, $v(0)=q$ and
$\alpha=\delta_{J-1}h_{J-1}=X^{-10}$. The base mean parameters have
$B_c=B_e=C_{B,d}$, $K=K_{B,d}$ and
$\mathcal L=C_{1,d}C_{B,d}$, with fixed nonnegative dimensional
constants. Its positive core radius and cutoff are exactly the
normalized parameters specified above. Their full mean reserve is small on
the short base interval; no boundary term is set to zero.

At time zero $m\cdot M_Bv=1$ throughout the oscillatory support.
The uniformly bounded base equations preserve
$m\cdot M_Bv\geq1/2$ on the chosen short interval. The seed
upper-pressure increment is therefore at most
$C_dX^{-10}+k_{J-1}^{-1/4}$. The forced high initial data vanish in
this zero-time forward construction, but the actual mean initial
fields need not vanish. Retain every finite mean and full tensor
corrector. The correction removing the residual has zero initial value,
so the initial increment is precisely this complete finite sum.

The estimates for that sum give interior and exterior gradient bounds
\[
 e_i=\frac{\alpha}{\delta_{J-1}}
          +\frac{\pi C_{\chi,d}\alpha+C_\nabla}{k_{J-1}},
 \qquad e_o=\frac{C_\nabla}{k_{J-1}}.
\]
Here $C_\nabla$ is the higher-grade source bound, controlled at the
prescribed seed frequency. Update $B_c$ by $e_i$ and $B_e$ by $e_o$;
retain the physical cutoff and set the new penalty to $C_{1,d}(B_c+e_i)$.
The difference obeys $B_c^+-B_e^+\leq C_{\rm lit}X^{D_0}$.
The nonoscillatory initial terms outside the small core are bounded by
$e_o$ and remain in the common fixed compact support. The full reserve
of this seed child is at most $1/8$ by the explicit short-time
and frequency estimates above.

The seed expansion is, at the fixed origin,
\[
 M_{J-1}=M_B+h(t)q(t)p(t)^T+\mathcal E(t),\qquad
 |\mathcal E|\leq k_{J-1}^{-1/4},\qquad h(0)=h_{J-1}.
\]
The normalized frame is transported by the actual base gradient.
Its activation parameters are exactly $a_J=1$ and $\beta_J=X^{-2}$.
The actual lifted flow estimate gives its polynomial particle-map
bound at $k_{J-1}$. Thus the complete input state for stage $J$ consists
of two exact solutions, the actual frame and shear decomposition,
the central gradient and Hessian bounds, the global one-sided pressure bound,
the strict fixed-cutoff reserve and the polynomial label bounds.
The first activation time is zero, so stage $J$ uses the forward primary
with its actual weighted growth and initial amplitude. Its seven estimates
were included in the sole scale choice. Its initial interior and exterior
updates can differ; their difference adds at most
$C_{\rm lit}X^{D_0}$. Thus $C_{\rm init}=2C_{\rm lit}$ covers both
initializations. Its pressure-time and initial-time reserve increments
are each at most $1/16$, leaving the first positive-history state with
full mean reserve at most $1/4$. Its actual activation time is $t_J>0$;
all later steps use the joined history construction.

Assume that state at stage $j$. The actual parent is available on
$[0,S_{j-1}]$, and $S_{j-1}-t_{j-1}=2W_j$ contains the target and
short retained extension. Restrict it to $[0,S_j]$ when forming the
actual central source and packet. The central horizon is precisely
\[
 H_j^{\rm sc}=a_j(S_j-t_{j-1})/\epsilon_j=\Theta_j;
\]
it is bounded by the numerical source envelope $\widehat\Theta_j$.
Coefficient bounds inherited from the longer parent remain valid under
restriction. The source envelope bounds the error contributions; it does
not extend any field past its actual time domain. Duhamel and energy
estimates and the relative target-size comparison are used only on the
actual retained interval, whose scaled extension past the target is
$\Delta\tau_j$. Along the older central solution,
$\dot B=-B^2-H_{j-2}$, which gives the central bounds on $B,\dot B$.
Lemma~\ref{lem:actual-endpoint} supplies the actual stationary history
for $j>J$, and the zero-time prescription supplies it for $j=J$.
The ambient comparison and source-envelope argument verify the
quantitative neighborhood, full transverse propagator and every
coefficient bound at the already prescribed $\rho_j,k_j$.
Take
\[
 \alpha_j=\frac{\delta_jh_j}{|m(t_j,0)||v(t_j,0)|}.
\]
The joined construction in Lemma~\ref{lem:exact-packet-child} then gives
an exact smooth odd reflection-equivariant solution $u_j$ on $[0,S_j]$.
Its full ambient tensor correctors, all means and the residual correction
are part of this same solution. The graph-flow identities give the
next actual particle-map invariant $K_{h,\rm next}\leq k_j^{D_d}$.

At the fixed central particle the expansion gives
\[
 \sup_{t\leq S_j}|M_j(t,0)-M_{j-1}(t,0)|
       \leq C_dh_j+k_j^{-1/4},\qquad
 \sup_{t\leq S_j}|H_j(t,0)-H_{j-1}(t,0)|
       \leq C_dh_jh_{j-1}+k_j^{-1/4}.
\]
Here the leading Hessian uses the sharp central parent strain bound.
Together with \eqref{eq:shear-separation} these preserve
\[
 \sup_{t\leq S_j}|M_j(t,0)|\leq C_Mh_j,\qquad
 \sup_{t\leq S_j}|H_j(t,0)|\leq C_Hh_jh_{j-1}.
\]
These are the central gradient and Hessian bounds needed by the central endpoint and
renewal equations. They are not asserted as full-space sharp norms.
Full-space coefficient and spatial derivative bounds instead follow
from the constructed particle-label fields and the determinant-one
polynomial inverse. The one-sided bounds
\eqref{eq:actual-increment-budgets} preserve the global pressure upper
bound, the initial interior and exterior quadratic-form bounds, and
the full fixed-cutoff reserve. The pressure is not assumed local:
its signed leading expansion and its residual error estimate hold
globally, with the leading tensor vanishing off the packet support.

For $t\in[t_j,S_j]$, the new normalized pair $p_2,q_2$ satisfies
\[
 M_j=M_{j-1}+h(t)q_2p_2^T+\mathcal E_j,
 \qquad |\mathcal E_j|\leq k_j^{-1/4},\qquad h(t_j)=h_j.
\]
Its frame equations have older field $M_{j-1}$, by the exact ray and
velocity equations, and $h'/h=-(M_{j-1})_{p_2p_2}-(M_{j-1})_{q_2q_2}$.
Lemma~\ref{lem:central-renewal} bounds each relative change in $a_j$
by $C_dx_{j-1}^{-4}$, whose sum is arbitrarily small. Since $a_J=1$,
the product of the successive relative factors remains in $[1/2,2]$. The same lemma gives directly
$\beta_{j+1}x_j^2\in[1/2,2]$; that estimate is not multiplied across
stages. Its compression bound and \eqref{eq:shear-separation} give
$p_2\cdot M_{j-1}p_2<-1$, while $|\mathcal E_j|<1$.
Also $(C_Mh_{j-1}+1)/h_j<\zeta$. These prove the two actual activation
conditions for the next endpoint problem. The central gradient and Hessian bounds on the entire
history and $t_j^{-1}\leq h_j$ prove the next endpoint-selection
hypotheses. The global one-sided mean bound and polynomial label
coefficients prove the full spatial history inverse hypotheses.
These estimates establish every component of the induction hypotheses
for the subsequent solution.

\section{The limiting datum and endpoint}
\label{sec:endpoint}
The initial estimates of Section~\ref{sec:initial-summability}, applied
to the same joined finite families and the actual source envelope, give
\[
 \sum_{j\geq J}\|u_j(0)-u_{j-1}(0)\|_{H^m}<\infty
                         \qquad\text{for every fixed }m.
\]
The seed and the first normal forward increment are two fixed
smooth increments; including the latter as the first term of this
series or in the initial finite datum gives the same limit.
Lemma~\ref{lem:common-support}
places all these initial velocities in one compact set. Sobolev embedding
at arbitrarily high fixed orders therefore proves that their single limit
$u_0$ is smooth, compactly supported and divergence free.

The target times increase to a finite positive $T_\infty\leq S_{\rm base}$.
At the central target the exact expansion gives
\[
 |\nabla u_j(t_j,0)|_{\rm F}
       \geq h_j-C_Mh_{j-1}-k_j^{-1/4}\longrightarrow\infty.
\]
Let $u$ be the unique maximal smooth Sobolev solution from $u_0$.
If its lifetime $T^*$ exceeded $T_\infty$, it would be smooth on a
closed interval containing all $t_j$, with bounded $H^{s+1}$ norm.
Apply Lemma~\ref{lem:stability} on the actual varying intervals
$[0,t_j]$ to $u_j$ and $u$. Their initial $H^s$ difference tends to
zero, so their gradients converge uniformly on those intervals.
The displayed divergence of the central gradients would contradict
boundedness of $\nabla u$ there. Thus $0<T^*\leq T_\infty<\infty$.
This argument does not assert equality of $T^*$ and $T_\infty$.
Lemma~\ref{lem:continuation} now supplies the exact Frobenius gradient
limsup and the once-counted antisymmetric-vorticity integral in $H_d$.
Lemma~\ref{lem:newton} fixes the Newton pressure normalization, and
Lemma~\ref{lem:local} supplies every required Sobolev order on each
shorter closed time interval.

The quantitative joined-source estimate and the prescribed scale
sequence verify the hypotheses at every finite stage. The initial limit,
finite lifetime and continuation conclusions therefore prove
Theorem~\ref{target:Hd}, with no additional assumption on the dimension
or on a prospective solution.

\section*{Use of artificial intelligence}
All mathematical derivations in this manuscript were produced using
artificial intelligence.

\begin{thebibliography}{9}
\bibitem{OAI}
OpenAI, \emph{Finite time blowup for the Euler equation}, manuscript,
\url{https://cdn.openai.com/pdf/315b36cd-ec98-4023-8342-93345194ece1/euler.pdf}.
\end{thebibliography}
\end{document}
